%%%%%%%%%%%%%%%%%%%%%%%%%%%%%%%%%%%%%%%%%%%%%%%%%%%%%%%%%%%%%%%%%%%%%%%%%%%%%%%%
%% APPENDIX F — Dataset Documentation, EDA, and Feature Engineering Pipeline
%% Comprehensive data appendix: 7 datasets, sample records, EDA, feature
%% engineering, frequency bands, statistical analysis, and cross-dataset comparison
%%%%%%%%%%%%%%%%%%%%%%%%%%%%%%%%%%%%%%%%%%%%%%%%%%%%%%%%%%%%%%%%%%%%%%%%%%%%%%%%
\normalsize\normalfont\color{black}

\section{Dataset Inventory and Provenance}
\label{app:dataset_inventory}

This appendix documents the seven datasets used in this dissertation, providing sample records, exploratory data analysis (EDA), the complete feature engineering pipeline, statistical analysis of data properties, and cross-dataset comparison. All datasets are publicly available, de-identified secondary data sources; no primary data collection involving human subjects was undertaken.

% MOVED TO CHAPTER 3 MAIN TEXT (Mar 2026)
\iffalse
\needspace{3\baselineskip}
\begingroup\singlespacing\tablefontsize
\noindent Table~\ref{tab:app_dataset_inventory_full} summarises dataset inventory and provenance for appendix review.

\begin{xltabular}
{\textwidth}{@{} p{2cm} p{2.5cm} p{1.2cm} c c c c p{1.2cm} L @{}}
\caption[Dataset inventory and provenance]{Dataset Inventory and Provenance: Seven Datasets Used in This Dissertation}
\label{tab:app_dataset_inventory_full}\\
\toprule
\thead{Domain} & \thead{Source} & \thead{Format} & \thead{N} & \thead{Ch.} & \thead{Hz} & \thead{Seg.} & \thead{Balance} & \thead{Licence} \\
\midrule
\endfirsthead
\endhead
\bottomrule
\endlastfoot
Schizophrenia & RepOD Repository & .eea & 84 & 1 & 128 & 8,400 & 0.87:1 & Institutional DUA \\
\addlinespace
Epilepsy & CHB-MIT (PhysioNet) & .edf & 102 & 23 & 256 & 5,100 & 1:1 & PhysioNet Credentialed \\
\addlinespace
Stress & DEAP + SAM40 & .bdf/.csv & 120 & 32 & 512/128 & 6,000 & 1:1 & EULA / Institutional DUA \\
\addlinespace
Autism (ASD) & KAU / Kaggle & .csv & 300 & 19 & 256 & 3,000 & 1:1 & Institutional DUA \\
\addlinespace
Parkinson  & PhysioNet & .csv & 50 & 32 & 256 & 500 & 1:1 & PhysioNet Credentialed \\
\addlinespace
Depression & OpenNeuro ds003478 & .set & 112 & Multi & 256 & 1,792 & 1.95:1 & CC-BY-4.0 \\
\addlinespace
Cancer & TCIA / Kaggle & .png & 20,000 & --- & --- & 20,000 & Varies & TCIA Data Usage Policy \\
\end{xltabular}
\endgroup

\vspace{0.5em}\noindent\textit{Note.} N = number of subjects or images; Ch.\ = number of channels; Hz = sampling rate; Seg.\ = approximate total segments after windowing; DUA = data use agreement; EULA = end-user licence agreement; CC-BY = Creative Commons Attribution. Cancer dataset uses peripheral blood smear (PBS) images rather than EEG signals.
\fi

%% ═══════════════════════════════════════════════════════════════════════
%% F.2 — SCHIZOPHRENIA DATASET
%% ═══════════════════════════════════════════════════════════════════════
% [SPACE-FIX] clearpage removed to eliminate empty page
\section{Dataset 1: Schizophrenia (RepOD)}
\label{app:data_schizophrenia}

\subsection{Dataset Description}

The schizophrenia dataset was obtained from the RepOD repository, containing single-channel EEG recordings from 84 subjects. Table~\ref{tab:app_f_schiz_params} summarises the key acquisition and cohort parameters.
\needspace{8\baselineskip}
{\tablestripes\tablefontsize
\begin{xltabular}{\textwidth}{@{} L R @{}}
\caption[Schizophrenia dataset parameters]{Schizophrenia Dataset: Key Acquisition and Cohort Parameters}\label{tab:app_f_schiz_params} \\
\toprule
\thead{Parameter} & \thead{Value} \\
\midrule
\endfirsthead
\multicolumn{2}{@{}l}{\textit{(\,Tab.~\ref{tab:app_f_schiz_params} continued from previous page\,)}} \\
\toprule
\thead{Parameter} & \thead{Value} \\
\midrule
\endhead
\bottomrule
\endlastfoot
Source repository & RepOD \\
Total subjects & 84 \\
Healthy controls & 39 \\
Schizophrenia patients & 45 \\
EEG channels & 1 (single-channel) \\
Sampling rate & 128~Hz \\
Segments per subject & $\approx$100 (2,000 samples each, $\approx$15.6~s) \\
Total segments & $\approx$8,400 \\
Class ratio (healthy:schizophrenia) & 0.87:1 \\
\end{xltabular}}
\vspace{0.5em}\noindent\textit{Note.} Recordings were acquired using standard clinical EEG equipment. The near-balanced class distribution (46.4\% healthy, 53.6\% schizophrenia) requires no oversampling.

Schizophrenia is characterised by three altered neural oscillatory patterns:

\begin{itemize}[leftmargin=*, itemsep=3pt]
\item \textbf{Reduced gamma-band synchronisation:} Diminished gamma-band ($>$30~Hz) coherence reflecting impaired cortical binding and sensory integration, a hallmark of disrupted NMDA receptor function.
\item \textbf{Elevated theta power:} Increased theta-band (4--8~Hz) power indicating frontal cortical slowing associated with working memory deficits and negative symptoms.
\item \textbf{Disrupted alpha asymmetry:} Abnormal hemispheric asymmetry in the alpha band (8--13~Hz), reflecting dysregulated thalamocortical oscillations and attentional processing.
\end{itemize}

\noindent These EEG biomarkers provide the discriminative features that the classification pipeline exploits. The dataset was selected because it offers clean binary labelling (healthy vs.\ schizophrenia) with a near-balanced class distribution (46.4\% healthy, 53.6\% schizophrenia).

\subsection{Sample Records (10 Rows)}

Table~\ref{tab:sample_schizophrenia} presents a representative 10-record sample from the schizophrenia dataset, illustrating the EEG signal values and class labels for both healthy controls and schizophrenia patients.
\needspace{3\baselineskip}
{\tablestripes\tablefontsize
\begin{xltabular}{\textwidth}{@{} C C R R R R R C @{}}
\caption[Schizophrenia dataset: 10-record sample]{Schizophrenia Dataset: Sample of 10 Records Showing EEG Signal Values and Labels}\label{tab:sample_schizophrenia} \\
\toprule
\thead{ID} & \thead{Seg.} & \thead{Mean ($\mu$V)} & \thead{Var ($\mu$V$^2$)} & \thead{ZCR} & \thead{$\alpha$ Power} & \thead{$\theta$ Power} & \thead{Label} \\
\midrule
\endfirsthead
\multicolumn{8}{@{}l}{\textit{(\,Tab.~\ref{tab:sample_schizophrenia} continued from previous page\,)}} \\
\toprule
\thead{ID} & \thead{Seg.} & \thead{Mean ($\mu$V)} & \thead{Var ($\mu$V$^2$)} & \thead{ZCR} & \thead{$\alpha$ Power} & \thead{$\theta$ Power} & \thead{Label} \\
\midrule
\endhead
\bottomrule
\endlastfoot
H01 & 1 & $-$0.42 & 18.73 & 0.312 & 4.21 & 2.87 & 0 (Healthy) \\
H01 & 2 & $-$0.38 & 17.95 & 0.308 & 4.35 & 2.91 & 0 (Healthy) \\
H15 & 1 & 0.11 & 20.14 & 0.325 & 3.98 & 3.12 & 0 (Healthy) \\
H28 & 1 & $-$0.19 & 16.82 & 0.298 & 4.52 & 2.76 & 0 (Healthy) \\
H39 & 1 & 0.05 & 19.41 & 0.317 & 4.08 & 2.95 & 0 (Healthy) \\
S01 & 1 & 0.73 & 28.56 & 0.389 & 2.14 & 4.63 & 1 (Schizo.) \\
S01 & 2 & 0.68 & 27.89 & 0.385 & 2.21 & 4.57 & 1 (Schizo.) \\
S12 & 1 & 0.91 & 31.42 & 0.401 & 1.98 & 5.01 & 1 (Schizo.) \\
S30 & 1 & 0.55 & 25.17 & 0.372 & 2.45 & 4.28 & 1 (Schizo.) \\
S45 & 1 & 0.82 & 29.83 & 0.395 & 2.07 & 4.81 & 1 (Schizo.) \\
\end{xltabular}}
\vspace{0.5em}\noindent\textit{Note.} Seg.\ = segment number; ZCR = zero-crossing rate; $\alpha$ Power = alpha-band (8--13~Hz) spectral power ($\mu$V$^2$/Hz); $\theta$ Power = theta-band (4--8~Hz) spectral power. Schizophrenia subjects show elevated theta power and reduced alpha power compared to healthy controls.

\subsection{EDA Summary}

Table~\ref{tab:eda_schizophrenia} reports the exploratory data analysis summary statistics for the schizophrenia dataset, including central tendency, dispersion, and distributional shape measures across all extracted features.
\needspace{3\baselineskip}
{\tablestripes\tablefontsize
\begin{xltabular}{\textwidth}{@{} L r r r r r r @{}}
\caption[Schizophrenia EDA summary]{Schizophrenia Dataset: Exploratory Data Analysis Summary}\label{tab:eda_schizophrenia} \\
\toprule
\thead{Feature} & \thead{Mean} & \thead{SD} & \thead{Median} & \thead{IQR} & \thead{Skew.} & \thead{Kurt.} \\
\midrule
\endfirsthead
\multicolumn{7}{@{}l}{\textit{(\,Tab.~\ref{tab:eda_schizophrenia} continued from previous page\,)}} \\
\toprule
\thead{Feature} & \thead{Mean} & \thead{SD} & \thead{Median} & \thead{IQR} & \thead{Skew.} & \thead{Kurt.} \\
\midrule
\endhead
\bottomrule
\endlastfoot
Mean amplitude ($\mu$V) & 0.18 & 1.24 & 0.09 & 1.52 & 0.31 & $-$0.18 \\
Variance ($\mu$V$^2$) & 23.47 & 7.82 & 22.15 & 10.34 & 0.67 & 0.12 \\
Zero-crossing rate & 0.351 & 0.042 & 0.348 & 0.056 & 0.24 & $-$0.35 \\
Alpha power (8--13~Hz) & 3.21 & 1.18 & 3.15 & 1.62 & 0.15 & $-$0.72 \\
Theta power (4--8~Hz) & 3.72 & 1.05 & 3.64 & 1.41 & 0.42 & $-$0.28 \\
Beta power (13--30~Hz) & 1.87 & 0.63 & 1.82 & 0.85 & 0.38 & $-$0.14 \\
Spectral entropy & 0.742 & 0.089 & 0.738 & 0.118 & 0.12 & $-$0.41 \\
Hjorth mobility & 0.215 & 0.048 & 0.212 & 0.064 & 0.28 & $-$0.19 \\
\end{xltabular}}
\vspace{0.5em}\noindent\textit{Note.} SD = standard deviation; IQR = interquartile range; Skew.\ = skewness (0 = symmetric); Kurt.\ = excess kurtosis (0 = mesokurtic). Moderate positive skew in variance and theta power indicates right-tailed distributions, supporting the use of non-parametric statistical tests.

\subsection{Key Observations}

Table~\ref{tab:schiz_key_obs} summarises the principal group-level differences between healthy controls and schizophrenia patients, quantifying each observation with effect direction, magnitude, and statistical significance.
\needspace{8\baselineskip}
{\tablestripes\tablefontsize
\begin{xltabular}{\textwidth}{@{} L c c c L @{}}
\caption[Schizophrenia key observations]{Schizophrenia Dataset: Key Group-Level Differences (Healthy vs.\ Schizophrenia)}\label{tab:schiz_key_obs} \\
\toprule
\thead{Observation} & \thead{Healthy} & \thead{Schizo.} & \thead{$p$-value} & \thead{Implication} \\
\midrule
\endfirsthead
\multicolumn{5}{@{}l}{\textit{(\,Tab.~\ref{tab:schiz_key_obs} continued from previous page\,)}} \\
\toprule
\thead{Observation} & \thead{Healthy} & \thead{Schizo.} & \thead{$p$-value} & \thead{Implication} \\
\midrule
\endhead
\bottomrule
\endlastfoot
Theta power ($\mu$V$^2$/Hz) & 2.93 & 4.66 & $<$.001 & Elevated theta consistent with EEG biomarker literature \\
\addlinespace
Alpha power ($\mu$V$^2$/Hz) & 4.23 & 2.17 & $<$.001 & Disrupted thalamocortical oscillations \\
\addlinespace
Signal variance & Baseline & +48\% & $<$.001 & Greater neural signal irregularity \\
\addlinespace
Class ratio & \multicolumn{2}{c}{0.87:1} & --- & No oversampling needed; stratified CV preserves ratios \\
\end{xltabular}}
\vspace{0.5em}\noindent\textit{Note.} CV = cross-validation. All $p$-values from Wilcoxon signed-rank test (non-normal data).

%% ═══════════════════════════════════════════════════════════════════════
%% F.3 — EPILEPSY DATASET
%% ═══════════════════════════════════════════════════════════════════════
% [SPACE-FIX] clearpage removed to eliminate empty page
\section{Dataset 2: Epilepsy (CHB-MIT, PhysioNet)}
\label{app:data_epilepsy}

\subsection{Dataset Description}

The epilepsy dataset originates from the CHB-MIT Scalp EEG Database hosted on PhysioNet, equally divided between non-seizure (inter-ictal) and seizure (pre-ictal/ictal) epochs. Table~\ref{tab:app_f_epilepsy_params} summarises the key acquisition parameters.
\needspace{8\baselineskip}
{\tablestripes\tablefontsize
\begin{xltabular}{\textwidth}{@{} L R @{}}
\caption[Epilepsy dataset parameters]{Epilepsy Dataset: Key Acquisition and Segmentation Parameters}\label{tab:app_f_epilepsy_params} \\
\toprule
\thead{Parameter} & \thead{Value} \\
\midrule
\endfirsthead
\multicolumn{2}{@{}l}{\textit{(\,Tab.~\ref{tab:app_f_epilepsy_params} continued from previous page\,)}} \\
\toprule
\thead{Parameter} & \thead{Value} \\
\midrule
\endhead
\bottomrule
\endlastfoot
Source repository & CHB-MIT (PhysioNet) \\
EEG channels & 23 (10--20 montage) \\
Recording sessions & 102 \\
Sampling rate & 256~Hz \\
Window size & 1,024 samples ($\approx$4~s) \\
Total segments & $\approx$5,100 \\
Class ratio (non-seizure:seizure) & 1:1 \\
\end{xltabular}}
\vspace{0.5em}\noindent\textit{Note.} CHB-MIT = Children's Hospital Boston--Massachusetts Institute of Technology. Seizure onset and offset timestamps are clinician-annotated.

Epileptic seizures produce three distinctive EEG patterns:

\begin{itemize}[leftmargin=*, itemsep=3pt]
\item \textbf{High-amplitude rhythmic discharges:} Sustained 5--8$\times$ amplitude increases during ictal epochs, driven by hypersynchronous neuronal firing in the seizure focus.
\item \textbf{Spike-and-wave complexes:} Characteristic sharp transients (spikes $<$70~ms) followed by slow waves ($\approx$200--500~ms), the hallmark electrographic signature of epileptiform activity.
\item \textbf{Abrupt frequency shifts:} Sudden transitions in the dominant frequency band, typically from background alpha/beta activity to high-amplitude delta/theta rhythms at seizure onset.
\end{itemize}

\noindent The CHB-MIT dataset is a gold-standard benchmark in the seizure detection literature, providing clinician-annotated seizure onset and offset timestamps that enable precise labelling. The dataset was selected for its multi-channel coverage (23 electrodes in the 10--20 montage), clean binary labelling, and perfect class balance.

\subsection{Sample Records (10 Rows)}

Table~\ref{tab:sample_epilepsy} documents a representative 10-record sample from the epilepsy dataset, contrasting multi-channel EEG feature values between non-seizure and seizure segments.
\needspace{3\baselineskip}
{\tablestripes\tablefontsize
\begin{xltabular}{\textwidth}{@{} C C R R R R R C @{}}
\caption[Epilepsy dataset: 10-record sample]{Epilepsy Dataset: Sample of 10 Records Showing Multi-Channel EEG Features}\label{tab:sample_epilepsy} \\
\toprule
\thead{Subj.} & \thead{Seg.} & \thead{$\bar{x}$ ($\mu$V)} & \thead{Var} & \thead{Line Len.} & \thead{$\delta$ Pwr} & \thead{$\beta$ Pwr} & \thead{Label} \\
\midrule
\endfirsthead
\multicolumn{8}{@{}l}{\textit{(\,Tab.~\ref{tab:sample_epilepsy} continued from previous page\,)}} \\
\toprule
\thead{Subj.} & \thead{Seg.} & \thead{$\bar{x}$ ($\mu$V)} & \thead{Var} & \thead{Line Len.} & \thead{$\delta$ Pwr} & \thead{$\beta$ Pwr} & \thead{Label} \\
\midrule
\endhead
\bottomrule
\endlastfoot
chb01 & 12 & 2.14 & 85.3 & 312.7 & 12.41 & 3.28 & 0 (Non-seizure) \\
chb01 & 45 & 1.98 & 79.6 & 298.4 & 11.87 & 3.15 & 0 (Non-seizure) \\
chb05 & 8 & 3.07 & 92.1 & 341.2 & 13.52 & 2.94 & 0 (Non-seizure) \\
chb12 & 22 & 2.45 & 88.7 & 325.8 & 12.98 & 3.07 & 0 (Non-seizure) \\
chb18 & 3 & 1.76 & 74.2 & 287.1 & 11.23 & 3.41 & 0 (Non-seizure) \\
chb01 & 67 & 18.42 & 542.8 & 1,247.3 & 45.71 & 8.92 & 1 (Seizure) \\
chb05 & 31 & 22.87 & 678.4 & 1,534.6 & 52.38 & 10.17 & 1 (Seizure) \\
chb07 & 15 & 15.63 & 489.2 & 1,102.8 & 41.29 & 7.84 & 1 (Seizure) \\
chb12 & 41 & 20.15 & 612.7 & 1,389.4 & 48.56 & 9.43 & 1 (Seizure) \\
chb18 & 28 & 17.94 & 528.1 & 1,198.5 & 44.12 & 8.56 & 1 (Seizure) \\
\end{xltabular}}
\vspace{0.5em}\noindent\textit{Note.} $\bar{x}$ = mean amplitude averaged across 23 channels; Var = mean variance; Line Len.\ = line length (sum of absolute consecutive differences); $\delta$ Pwr = delta-band power (0.5--4~Hz); $\beta$ Pwr = beta-band power (13--30~Hz). Seizure segments show 5--8$\times$ higher amplitude and delta power.

\subsection{EDA Summary}

Table~\ref{tab:eda_epilepsy} summarises the exploratory data analysis results for the epilepsy dataset, revealing strong positive skew and heavy tails driven by the bimodal seizure/non-seizure distribution.
\needspace{3\baselineskip}
{\tablestripes\tablefontsize
\begin{xltabular}{\textwidth}{@{} L r r r r r r @{}}
\caption[Epilepsy EDA summary]{Epilepsy Dataset: Exploratory Data Analysis Summary}\label{tab:eda_epilepsy} \\
\toprule
\thead{Feature} & \thead{Mean} & \thead{SD} & \thead{Median} & \thead{IQR} & \thead{Skew.} & \thead{Kurt.} \\
\midrule
\endfirsthead
\multicolumn{7}{@{}l}{\textit{(\,Tab.~\ref{tab:eda_epilepsy} continued from previous page\,)}} \\
\toprule
\thead{Feature} & \thead{Mean} & \thead{SD} & \thead{Median} & \thead{IQR} & \thead{Skew.} & \thead{Kurt.} \\
\midrule
\endhead
\bottomrule
\endlastfoot
Mean amplitude ($\mu$V) & 10.64 & 8.92 & 5.82 & 14.37 & 1.42 & 1.18 \\
Variance ($\mu$V$^2$) & 312.4 & 245.8 & 189.7 & 398.2 & 1.87 & 3.24 \\
Line length & 718.5 & 482.3 & 498.2 & 745.1 & 1.53 & 1.67 \\
Delta power (0.5--4~Hz) & 28.92 & 17.84 & 21.45 & 28.73 & 1.38 & 1.05 \\
Theta power (4--8~Hz) & 8.47 & 5.12 & 6.82 & 7.94 & 1.21 & 0.82 \\
Beta power (13--30~Hz) & 5.87 & 3.24 & 4.62 & 4.87 & 1.15 & 0.74 \\
Spectral entropy & 0.681 & 0.124 & 0.695 & 0.168 & $-$0.32 & $-$0.58 \\
Hjorth complexity & 1.482 & 0.312 & 1.425 & 0.418 & 0.67 & 0.23 \\
\end{xltabular}}
\vspace{0.5em}\noindent\textit{Note.} Strong positive skew ($>$1.0) in amplitude, variance, and delta power reflects the bimodal nature of seizure vs.\ non-seizure data. High kurtosis in variance (3.24) indicates heavy tails from seizure episodes. These distributional properties justify non-parametric tests (Wilcoxon signed-rank, Mann--Whitney U).

\subsection{Key Observations}

Table~\ref{tab:epilepsy_key_obs} quantifies the principal seizure-versus-non-seizure contrasts, distribution characteristics, and design advantages of the CHB-MIT dataset.
\needspace{8\baselineskip}
{\tablestripes\tablefontsize
\begin{xltabular}{\textwidth}{@{} L c L @{}}
\caption[Epilepsy key observations]{Epilepsy Dataset: Key Observations on Seizure vs.\ Non-Seizure Segments}\label{tab:epilepsy_key_obs} \\
\toprule
\thead{Observation} & \thead{Magnitude} & \thead{Implication} \\
\midrule
\endfirsthead
\multicolumn{3}{@{}l}{\textit{(\,Tab.~\ref{tab:epilepsy_key_obs} continued from previous page\,)}} \\
\toprule
\thead{Observation} & \thead{Magnitude} & \thead{Implication} \\
\midrule
\endhead
\bottomrule
\endlastfoot
Mean amplitude (seizure vs.\ non-seizure) & 5--8$\times$ higher & Strong discriminative signal for classifier \\
\addlinespace
Delta power (seizure vs.\ non-seizure) & 3--4$\times$ higher & Primary frequency-domain biomarker \\
\addlinespace
Feature skewness (time \& freq.\ domains) & $>$1.0 (positive) & Non-Gaussian; justifies Wilcoxon over paired $t$-test \\
\addlinespace
23-channel montage & Spatial coverage & Enables seizure focus localisation via CNNs \\
\addlinespace
Class balance & 1:1 & No oversampling required \\
\end{xltabular}}

%% ═══════════════════════════════════════════════════════════════════════
%% F.4 — STRESS DATASET
%% ═══════════════════════════════════════════════════════════════════════
% [SPACE-FIX] clearpage removed to eliminate empty page
\section{Dataset 3: Stress (DEAP + SAM40 + Stress-RAG)}
\label{app:data_stress}

\subsection{Dataset Description}

The stress detection pipeline uses three complementary data sources, compared in Table~\ref{tab:app_f_stress_sources}.
\needspace{8\baselineskip}
{\tablestripes\tablefontsize
\begin{xltabular}{\textwidth}{@{} l c c c L @{}}
\caption[Stress dataset source comparison]{Stress Dataset: Comparison of Three Complementary Data Sources}\label{tab:app_f_stress_sources} \\
\toprule
\thead{Source} & \thead{Channels} & \thead{Sample Rate} & \thead{Subjects} & \thead{Task / Stimuli} \\
\midrule
\endfirsthead
\multicolumn{5}{@{}l}{\textit{(\,Tab.~\ref{tab:app_f_stress_sources} continued from previous page\,)}} \\
\toprule
\thead{Source} & \thead{Channels} & \thead{Sample Rate} & \thead{Subjects} & \thead{Task / Stimuli} \\
\midrule
\endhead
\bottomrule
\endlastfoot
DEAP & 32 EEG & 512~Hz & 32 & Emotional video stimuli; SAM arousal/valence ratings derive stress labels \\
SAM40 & 32 EEG & 128~Hz & 40 & Cognitive stress tasks (arithmetic, Stroop); ground-truth task labels \\
Stress-RAG & Multi-modal & Varies & 48 & Real-time wearable sensors (ECG, EDA, PPG, temperature, accelerometer) via Emotiv EPOC X \\
\end{xltabular}}
\vspace{0.5em}\noindent\textit{Note.} DEAP = Database for Emotion Analysis using Physiological Signals; SAM = self-assessment manikin; Stress-RAG = the real-time wearable module in the agenticfinder system using the Emotiv EPOC X and the neuro-wearable multi-sensor platform.

Combined, the stress pipeline processes 120 subjects with approximately 6,000 segments after augmentation. Labels are binary: relaxed (baseline resting state, label~0) vs.\ stressed (task-induced cognitive stress, label~1).

\subsection{Sample Records (10 Rows)}

Table~\ref{tab:sample_stress} presents a representative 10-record sample from the stress dataset, illustrating EEG and physiological sensor features under both relaxed and stressed conditions.
\needspace{3\baselineskip}
{\tablestripes\tablefontsize
\begin{xltabular}{\textwidth}{@{} C C R R R R R C @{}}
\caption[Stress dataset: 10-record sample]{Stress Dataset: Sample of 10 Records Showing EEG and Physiological Features}\label{tab:sample_stress} \\
\toprule
\thead{Src.} & \thead{Subj.} & \thead{$\alpha$ Asym.} & \thead{$\beta$ Pwr} & \thead{HRV} & \thead{EDA} & \thead{Temp.} & \thead{Label} \\
\midrule
\endfirsthead
\multicolumn{8}{@{}l}{\textit{(\,Tab.~\ref{tab:sample_stress} continued from previous page\,)}} \\
\toprule
\thead{Src.} & \thead{Subj.} & \thead{$\alpha$ Asym.} & \thead{$\beta$ Pwr} & \thead{HRV} & \thead{EDA} & \thead{Temp.} & \thead{Label} \\
\midrule
\endhead
\bottomrule
\endlastfoot
DEAP & S01 & 0.124 & 2.87 & 72.4 & 1.23 & 36.2 & 0 (Relaxed) \\
DEAP & S08 & 0.098 & 2.91 & 68.7 & 1.45 & 36.1 & 0 (Relaxed) \\
SAM40 & P05 & 0.112 & 3.02 & 75.1 & 1.18 & 36.3 & 0 (Relaxed) \\
SAM40 & P19 & 0.135 & 2.78 & 70.2 & 1.31 & 36.2 & 0 (Relaxed) \\
RAG & W03 & 0.108 & 2.95 & 71.8 & 1.27 & 36.4 & 0 (Relaxed) \\
DEAP & S01 & $-$0.187 & 5.42 & 48.3 & 4.87 & 37.1 & 1 (Stressed) \\
DEAP & S15 & $-$0.215 & 5.78 & 45.1 & 5.23 & 37.3 & 1 (Stressed) \\
SAM40 & P05 & $-$0.198 & 5.15 & 51.2 & 4.56 & 36.9 & 1 (Stressed) \\
SAM40 & P27 & $-$0.172 & 4.98 & 53.7 & 4.12 & 36.8 & 1 (Stressed) \\
RAG & W03 & $-$0.205 & 5.31 & 46.8 & 5.01 & 37.2 & 1 (Stressed) \\
\end{xltabular}}
\vspace{0.5em}\noindent\textit{Note.} Src.\ = data source; $\alpha$ Asym.\ = frontal alpha asymmetry (positive = left-dominant/relaxed, negative = right-dominant/stressed); HRV = heart rate variability (ms); EDA = electrodermal activity ($\mu$S); Temp.\ = skin temperature (\textdegree{}C). Stressed subjects show reversed alpha asymmetry, elevated beta power, reduced HRV, and increased EDA.

\subsection{EDA Summary}

Table~\ref{tab:eda_stress} reports the exploratory data analysis summary for the stress dataset, covering distributional properties of both EEG-derived and wearable physiological features.
\needspace{3\baselineskip}
{\tablestripes\tablefontsize
\begin{xltabular}{\textwidth}{@{} L r r r r r r @{}}
\caption[Stress EDA summary]{Stress Dataset: Exploratory Data Analysis Summary}\label{tab:eda_stress} \\
\toprule
\thead{Feature} & \thead{Mean} & \thead{SD} & \thead{Median} & \thead{IQR} & \thead{Skew.} & \thead{Kurt.} \\
\midrule
\endfirsthead
\multicolumn{7}{@{}l}{\textit{(\,Tab.~\ref{tab:eda_stress} continued from previous page\,)}} \\
\toprule
\thead{Feature} & \thead{Mean} & \thead{SD} & \thead{Median} & \thead{IQR} & \thead{Skew.} & \thead{Kurt.} \\
\midrule
\endhead
\bottomrule
\endlastfoot
Frontal alpha asymmetry & $-$0.031 & 0.178 & $-$0.024 & 0.241 & $-$0.18 & $-$0.42 \\
Beta power ($\mu$V$^2$/Hz) & 3.98 & 1.42 & 3.72 & 1.95 & 0.54 & $-$0.31 \\
HRV (RMSSD, ms) & 59.8 & 14.2 & 61.3 & 19.7 & $-$0.28 & $-$0.45 \\
EDA ($\mu$S) & 3.05 & 1.87 & 2.64 & 2.72 & 0.72 & $-$0.18 \\
Skin temperature (\textdegree{}C) & 36.6 & 0.48 & 36.5 & 0.62 & 0.34 & $-$0.52 \\
Theta/beta ratio & 1.87 & 0.52 & 1.82 & 0.71 & 0.41 & $-$0.15 \\
Spectral entropy & 0.718 & 0.095 & 0.712 & 0.127 & 0.15 & $-$0.38 \\
\end{xltabular}}
\vspace{0.5em}\noindent\textit{Note.} Alpha asymmetry distribution is near-symmetric (skew $= -0.18$), suitable for parametric tests. EDA shows moderate positive skew (0.72) due to phasic skin conductance responses during stress tasks. HRV shows mild negative skew reflecting reduced variability under stress.

\subsection{Key Observations}

Table~\ref{tab:stress_key_obs} summarises the principal physiological contrasts between relaxed and stressed states, highlighting the multi-modal advantage of the combined EEG and wearable sensor pipeline.
\needspace{8\baselineskip}
{\tablestripes\tablefontsize
\begin{xltabular}{\textwidth}{@{} L c L @{}}
\caption[Stress key observations]{Stress Dataset: Key Observations (Relaxed vs.\ Stressed)}\label{tab:stress_key_obs} \\
\toprule
\thead{Observation} & \thead{Change} & \thead{Implication} \\
\midrule
\endfirsthead
\multicolumn{3}{@{}l}{\textit{(\,Tab.~\ref{tab:stress_key_obs} continued from previous page\,)}} \\
\toprule
\thead{Observation} & \thead{Change} & \thead{Implication} \\
\midrule
\endhead
\bottomrule
\endlastfoot
Frontal alpha asymmetry & Positive $\to$ negative & Consistent with Davidson's approach--withdrawal model \\
\addlinespace
EDA amplitude & 3--4$\times$ increase & Sympathetic nervous system activation \\
\addlinespace
Multi-modal fusion accuracy & 94.17\% & Exceeds single-modality EEG-only approaches \\
\addlinespace
Class balance & 1:1 & Task-based labelling ensures balance \\
\end{xltabular}}
\vspace{0.5em}\noindent\textit{Note.} EDA = electrodermal activity.

%% ═══════════════════════════════════════════════════════════════════════
%% F.5 — AUTISM (ASD) DATASET
%% ═══════════════════════════════════════════════════════════════════════
% [SPACE-FIX] clearpage removed to eliminate empty page
\section{Dataset 4: Autism Spectrum Disorder (KAU/Kaggle)}
\label{app:data_autism}

\subsection{Dataset Description}

The autism dataset contains pre-extracted EEG features from the King Abdulaziz University (KAU) repository and supplementary Kaggle research datasets. Children aged 6--18 years were recruited with parental consent under IRB-approved protocols at the originating institutions. Table~\ref{tab:app_f_asd_params} summarises the key acquisition and cohort parameters.
\needspace{8\baselineskip}
{\tablestripes\tablefontsize
\begin{xltabular}{\textwidth}{@{} L R @{}}
\caption[Autism dataset parameters]{Autism Spectrum Disorder Dataset: Key Acquisition and Cohort Parameters}\label{tab:app_f_asd_params} \\
\toprule
\thead{Parameter} & \thead{Value} \\
\midrule
\endfirsthead
\multicolumn{2}{@{}l}{\textit{(\,Tab.~\ref{tab:app_f_asd_params} continued from previous page\,)}} \\
\toprule
\thead{Parameter} & \thead{Value} \\
\midrule
\endhead
\bottomrule
\endlastfoot
Source repository & KAU / Kaggle \\
Total subjects & 300 \\
Neurotypical (TD) controls & 150 \\
Clinically diagnosed ASD & 150 \\
EEG channels & 19 (10--20 montage) \\
Sampling rate & 256~Hz \\
Total segments & $\approx$3,000 \\
Class ratio (TD:ASD) & 1:1 \\
\end{xltabular}}
\vspace{0.5em}\noindent\textit{Note.} TD = typically developing; ASD = Autism Spectrum Disorder; KAU = King Abdulaziz University. The 10--20 montage provides standard international electrode placement.

ASD is characterised by three altered functional connectivity patterns:

\begin{itemize}[leftmargin=*, itemsep=3pt]
\item \textbf{Reduced long-range coherence:} Diminished inter-hemispheric and fronto-parietal coherence reflecting disrupted white-matter tract integrity, consistent with the ``disconnection hypothesis'' of autism.
\item \textbf{Increased short-range connectivity:} Elevated local connectivity in frontal and temporal regions, suggesting cortical over-connectivity that may underlie restricted and repetitive behaviours.
\item \textbf{Atypical alpha/theta power distributions:} Reduced alpha-band (8--13~Hz) power and elevated theta-band (4--8~Hz) power, yielding a theta/alpha ratio nearly double that of neurotypical controls.
\end{itemize}

\noindent The agenticfinder pipeline achieves 97.67\% accuracy on this dataset, making it the second-highest-performing domain.

\subsection{Sample Records (10 Rows)}

Table~\ref{tab:sample_autism} documents a representative 10-record sample from the autism dataset, contrasting pre-extracted EEG connectivity and spectral features between neurotypical and ASD subjects.
\needspace{3\baselineskip}
{\tablestripes\tablefontsize
\begin{xltabular}{\textwidth}{@{} C R R R R R R C @{}}
\caption[Autism dataset: 10-record sample]{Autism Dataset: Sample of 10 Records Showing Pre-Extracted EEG Features}\label{tab:sample_autism} \\
\toprule
\thead{ID} & \thead{$\alpha$ Pwr} & \thead{$\theta$ Pwr} & \thead{Coher.} & \thead{PLV} & \thead{Entropy} & \thead{Hjorth} & \thead{Label} \\
\midrule
\endfirsthead
\multicolumn{8}{@{}l}{\textit{(\,Tab.~\ref{tab:sample_autism} continued from previous page\,)}} \\
\toprule
\thead{ID} & \thead{$\alpha$ Pwr} & \thead{$\theta$ Pwr} & \thead{Coher.} & \thead{PLV} & \thead{Entropy} & \thead{Hjorth} & \thead{Label} \\
\midrule
\endhead
\bottomrule
\endlastfoot
TD001 & 5.82 & 2.41 & 0.724 & 0.681 & 0.812 & 0.198 & 0 (TD) \\
TD025 & 5.47 & 2.58 & 0.718 & 0.672 & 0.798 & 0.205 & 0 (TD) \\
TD078 & 6.01 & 2.34 & 0.731 & 0.695 & 0.824 & 0.192 & 0 (TD) \\
TD112 & 5.63 & 2.49 & 0.709 & 0.668 & 0.805 & 0.201 & 0 (TD) \\
TD150 & 5.91 & 2.37 & 0.728 & 0.689 & 0.818 & 0.195 & 0 (TD) \\
ASD001 & 3.14 & 4.28 & 0.521 & 0.487 & 0.692 & 0.267 & 1 (ASD) \\
ASD034 & 2.98 & 4.51 & 0.498 & 0.462 & 0.678 & 0.278 & 1 (ASD) \\
ASD089 & 3.27 & 4.12 & 0.534 & 0.501 & 0.704 & 0.258 & 1 (ASD) \\
ASD120 & 3.05 & 4.42 & 0.512 & 0.475 & 0.685 & 0.271 & 1 (ASD) \\
ASD150 & 3.21 & 4.19 & 0.528 & 0.493 & 0.698 & 0.263 & 1 (ASD) \\
\end{xltabular}}
\vspace{0.5em}\noindent\textit{Note.} Coher.\ = inter-channel coherence (mean across all channel pairs); PLV = phase-locking value; Hjorth = Hjorth mobility parameter. ASD subjects show reduced alpha power, elevated theta power, lower coherence, and lower phase-locking values, reflecting disrupted functional connectivity.

\subsection{EDA Summary}

Table~\ref{tab:eda_autism} summarises the exploratory data analysis results for the autism dataset, confirming near-symmetric feature distributions that support parametric testing.
\needspace{3\baselineskip}
{\tablestripes\tablefontsize
\begin{xltabular}{\textwidth}{@{} L r r r r r r @{}}
\caption[Autism EDA summary]{Autism Dataset: Exploratory Data Analysis Summary}\label{tab:eda_autism} \\
\toprule
\thead{Feature} & \thead{Mean} & \thead{SD} & \thead{Median} & \thead{IQR} & \thead{Skew.} & \thead{Kurt.} \\
\midrule
\endfirsthead
\multicolumn{7}{@{}l}{\textit{(\,Tab.~\ref{tab:eda_autism} continued from previous page\,)}} \\
\toprule
\thead{Feature} & \thead{Mean} & \thead{SD} & \thead{Median} & \thead{IQR} & \thead{Skew.} & \thead{Kurt.} \\
\midrule
\endhead
\bottomrule
\endlastfoot
Alpha power ($\mu$V$^2$/Hz) & 4.45 & 1.48 & 4.38 & 2.01 & 0.12 & $-$0.64 \\
Theta power ($\mu$V$^2$/Hz) & 3.38 & 1.12 & 3.31 & 1.52 & 0.24 & $-$0.48 \\
Coherence & 0.619 & 0.112 & 0.621 & 0.152 & $-$0.05 & $-$0.71 \\
Phase-locking value & 0.582 & 0.108 & 0.585 & 0.147 & $-$0.08 & $-$0.65 \\
Spectral entropy & 0.751 & 0.072 & 0.748 & 0.098 & 0.09 & $-$0.42 \\
Hjorth mobility & 0.233 & 0.042 & 0.231 & 0.057 & 0.15 & $-$0.35 \\
Frontal $\theta$/$\alpha$ ratio & 0.82 & 0.28 & 0.79 & 0.38 & 0.35 & $-$0.22 \\
\end{xltabular}}
\vspace{0.5em}\noindent\textit{Note.} Near-symmetric distributions (skewness $<$0.35 for all features) support parametric testing. Low kurtosis ($<$0) indicates light-tailed distributions with few extreme outliers. The theta/alpha ratio is the strongest single discriminator between ASD and TD groups.

\subsection{Key Observations}

Table~\ref{tab:asd_key_obs} summarises the principal group-level differences between neurotypical controls and ASD subjects, with percentage changes and their connectivity implications.
\needspace{8\baselineskip}
{\tablestripes\tablefontsize
\begin{xltabular}{\textwidth}{@{} L c L @{}}
\caption[ASD key observations]{Autism Dataset: Key Group-Level Differences (TD vs.\ ASD)}\label{tab:asd_key_obs} \\
\toprule
\thead{Observation} & \thead{Change (ASD vs.\ TD)} & \thead{Implication} \\
\midrule
\endfirsthead
\multicolumn{3}{@{}l}{\textit{(\,Tab.~\ref{tab:asd_key_obs} continued from previous page\,)}} \\
\toprule
\thead{Observation} & \thead{Change (ASD vs.\ TD)} & \thead{Implication} \\
\midrule
\endhead
\bottomrule
\endlastfoot
Alpha power & $-$45\% & Reduced thalamocortical oscillations \\
\addlinespace
Theta power & +75\% & $\theta$/$\alpha$ ratio nearly double that of TD controls \\
\addlinespace
Coherence and PLV & $-$28--30\% & Supports ``disconnection hypothesis'' of autism \\
\addlinespace
Feature distributions & Near-symmetric (skew $<$0.35) & Parametric paired $t$-tests appropriate \\
\addlinespace
Class balance & 150:150 (1:1) & Unbiased training; no oversampling needed \\
\end{xltabular}}
\vspace{0.5em}\noindent\textit{Note.} TD = typically developing; PLV = phase-locking value.

%% ═══════════════════════════════════════════════════════════════════════
%% F.6 — PARKINSON DATASET
%% ═══════════════════════════════════════════════════════════════════════
% [SPACE-FIX] clearpage removed to eliminate empty page
\section{Dataset 5: Parkinson's Disease (PhysioNet)}
\label{app:data_parkinson}

\subsection{Dataset Description}

The Parkinson's disease  dataset contains pre-extracted EEG features from the PhysioNet repository. Despite the small sample size, the dataset yields 100\% classification accuracy due to highly discriminative EEG biomarkers. Table~\ref{tab:app_f_pd_params} summarises the key acquisition and cohort parameters.
\needspace{8\baselineskip}
{\tablestripes\tablefontsize
\begin{xltabular}{\textwidth}{@{} L R @{}}
\caption[Parkinson dataset parameters]{Parkinson's Disease Dataset: Key Acquisition and Cohort Parameters}\label{tab:app_f_pd_params} \\
\toprule
\thead{Parameter} & \thead{Value} \\
\midrule
\endfirsthead
\multicolumn{2}{@{}l}{\textit{(\,Tab.~\ref{tab:app_f_pd_params} continued from previous page\,)}} \\
\toprule
\thead{Parameter} & \thead{Value} \\
\midrule
\endhead
\bottomrule
\endlastfoot
Source repository & PhysioNet \\
Total subjects & 50 \\
Healthy controls (age-matched) & 25 \\
Clinically confirmed PD patients & 25 \\
EEG channels & 32 \\
Sampling rate & 256~Hz \\
Total segments & $\approx$500 \\
Class ratio (healthy:PD) & 1:1 \\
\end{xltabular}}
\vspace{0.5em}\noindent\textit{Note.} PD = Parkinson's disease. Age-matched controls minimise confounding from age-related EEG changes.

PD patients exhibit three characteristic EEG biomarker alterations:

\begin{itemize}[leftmargin=*, itemsep=3pt]
\item \textbf{Beta-band desynchronisation:} Reduced beta-band (13--30~Hz) power in sensorimotor cortex, reflecting impaired basal ganglia--cortical loop function and dopaminergic depletion.
\item \textbf{Elevated delta power:} Increased slow-wave (0.5--4~Hz) activity indicating cortical slowing and subcortical dysfunction characteristic of neurodegeneration.
\item \textbf{Reduced alpha peak frequency:} Alpha peak shifts from $\approx$10.3~Hz (healthy) to $\approx$7.7~Hz , a hallmark pre-motor biomarker detectable 3--7 years before motor symptom onset.
\end{itemize}

\noindent These EEG abnormalities make EEG-based screening a promising early-detection tool. The agenticfinder pipeline exploits these pre-motor biomarkers, achieving perfect discrimination with Cohen's $d > 3.0$ (very large effect size) for the top discriminative features.

\subsection{Sample Records (10 Rows)}

Table~\ref{tab:sample_parkinson} presents a representative 10-record sample from the Parkinson's disease dataset, highlighting the large group-level differences in delta power, alpha peak frequency, and beta power.
\needspace{3\baselineskip}
{\tablestripes\tablefontsize
\begin{xltabular}{\textwidth}{@{} C R R R R R R C @{}}
\caption[Parkinson dataset: 10-record sample]{Parkinson's Disease Dataset: Sample of 10 Records Showing EEG Features}\label{tab:sample_parkinson} \\
\toprule
\thead{ID} & \thead{$\delta$ Pwr} & \thead{$\alpha$ Peak} & \thead{$\beta$ Pwr} & \thead{Entropy} & \thead{Coher.} & \thead{Hjorth} & \thead{Label} \\
\midrule
\endfirsthead
\multicolumn{8}{@{}l}{\textit{(\,Tab.~\ref{tab:sample_parkinson} continued from previous page\,)}} \\
\toprule
\thead{ID} & \thead{$\delta$ Pwr} & \thead{$\alpha$ Peak} & \thead{$\beta$ Pwr} & \thead{Entropy} & \thead{Coher.} & \thead{Hjorth} & \thead{Label} \\
\midrule
\endhead
\bottomrule
\endlastfoot
HC01 & 3.42 & 10.2 & 4.87 & 0.842 & 0.751 & 0.187 & 0 (Healthy) \\
HC07 & 3.28 & 10.5 & 4.92 & 0.838 & 0.745 & 0.191 & 0 (Healthy) \\
HC12 & 3.51 & 10.1 & 5.01 & 0.851 & 0.758 & 0.184 & 0 (Healthy) \\
HC18 & 3.35 & 10.4 & 4.78 & 0.845 & 0.742 & 0.189 & 0 (Healthy) \\
HC25 & 3.47 & 10.3 & 4.95 & 0.848 & 0.754 & 0.186 & 0 (Healthy) \\
PD01 & 8.92 & 7.8 & 2.14 & 0.621 & 0.528 & 0.312 & 1  \\
PD06 & 9.15 & 7.5 & 1.98 & 0.612 & 0.515 & 0.324 & 1  \\
PD12 & 8.78 & 8.1 & 2.27 & 0.634 & 0.541 & 0.305 & 1  \\
PD19 & 9.28 & 7.3 & 1.87 & 0.605 & 0.502 & 0.331 & 1  \\
PD25 & 8.95 & 7.6 & 2.08 & 0.618 & 0.521 & 0.318 & 1  \\
\end{xltabular}}
\vspace{0.5em}\noindent\textit{Note.} $\delta$ Pwr = delta-band power (0.5--4~Hz); $\alpha$ Peak = alpha peak frequency (Hz); $\beta$ Pwr = beta-band power (13--30~Hz); Coher.\ = mean coherence. PD subjects show 2.6$\times$ higher delta power, 2.5~Hz lower alpha peak frequency, and 57\% reduced beta power. These large effect sizes explain the 100\% classification accuracy.

\subsection{EDA Summary}

Table~\ref{tab:eda_parkinson} reports the exploratory data analysis summary for the Parkinson's disease dataset, showing strongly platykurtic distributions that reflect the bimodal separation between healthy and PD groups.
\needspace{3\baselineskip}
{\tablestripes\tablefontsize
\begin{xltabular}{\textwidth}{@{} L r r r r r r @{}}
\caption[Parkinson EDA summary]{Parkinson's Disease Dataset: Exploratory Data Analysis Summary}\label{tab:eda_parkinson} \\
\toprule
\thead{Feature} & \thead{Mean} & \thead{SD} & \thead{Median} & \thead{IQR} & \thead{Skew.} & \thead{Kurt.} \\
\midrule
\endfirsthead
\multicolumn{7}{@{}l}{\textit{(\,Tab.~\ref{tab:eda_parkinson} continued from previous page\,)}} \\
\toprule
\thead{Feature} & \thead{Mean} & \thead{SD} & \thead{Median} & \thead{IQR} & \thead{Skew.} & \thead{Kurt.} \\
\midrule
\endhead
\bottomrule
\endlastfoot
Delta power ($\mu$V$^2$/Hz) & 6.17 & 3.12 & 5.84 & 5.28 & 0.42 & $-$1.12 \\
Alpha peak freq.\ (Hz) & 8.92 & 1.42 & 9.15 & 2.38 & $-$0.35 & $-$1.28 \\
Beta power ($\mu$V$^2$/Hz) & 3.51 & 1.58 & 3.42 & 2.65 & 0.18 & $-$1.15 \\
Spectral entropy & 0.731 & 0.124 & 0.738 & 0.198 & $-$0.21 & $-$1.08 \\
Coherence & 0.637 & 0.128 & 0.641 & 0.205 & $-$0.12 & $-$1.21 \\
Hjorth mobility & 0.249 & 0.072 & 0.245 & 0.118 & 0.28 & $-$0.95 \\
Delta/alpha ratio & 1.24 & 0.85 & 0.98 & 1.42 & 0.78 & $-$0.52 \\
\end{xltabular}}
\vspace{0.5em}\noindent\textit{Note.} Strongly platykurtic distributions (kurtosis $< -1.0$ for most features) reflect the bimodal separation between healthy and PD groups. The delta/alpha ratio shows the highest skewness (0.78), driven by the PD group's elevated delta and reduced alpha power.

\subsection{Key Observations}

Table~\ref{tab:pd_key_obs} quantifies the principal findings from the Parkinson's disease dataset, highlighting the bimodal feature separation and the methodological trade-offs imposed by the small sample size.
\needspace{8\baselineskip}
{\tablestripes\tablefontsize
\begin{xltabular}{\textwidth}{@{} L c L @{}}
\caption[PD key observations]{Parkinson's Disease Dataset: Key Observations}\label{tab:pd_key_obs} \\
\toprule
\thead{Observation} & \thead{Value} & \thead{Implication} \\
\midrule
\endfirsthead
\multicolumn{3}{@{}l}{\textit{(\,Tab.~\ref{tab:pd_key_obs} continued from previous page\,)}} \\
\toprule
\thead{Observation} & \thead{Value} & \thead{Implication} \\
\midrule
\endhead
\bottomrule
\endlastfoot
Cohen's $d$ for top-5 features & $>$3.0 & Largest effect sizes across all six EEG domains \\
\addlinespace
Alpha peak frequency shift & 10.3~Hz $\to$ 7.7~Hz & Hallmark pre-motor biomarker  \\
\addlinespace
Bootstrap CI width (10-fold CV) & $<$1\% & Robust generalisation despite $N = 50$ \\
\addlinespace
Clinical trade-off & Perfect accuracy & Strong biomarkers; requires multi-site validation \\
\end{xltabular}}
\vspace{0.5em}\noindent\textit{Note.} CI = confidence interval (1,000 bootstrap resamples, 95\%); CV = cross-validation.

%% ═══════════════════════════════════════════════════════════════════════
%% F.7 — DEPRESSION DATASET
%% ═══════════════════════════════════════════════════════════════════════
% [SPACE-FIX] clearpage removed to eliminate empty page
\section{Dataset 6: Depression (OpenNeuro ds003478)}
\label{app:data_depression}

\subsection{Dataset Description}

The depression dataset originates from OpenNeuro (ds003478), containing multi-channel EEG recordings from 122 subjects. After excluding ambiguous cases, 112 subjects remained for analysis. Table~\ref{tab:app_f_depression_params} summarises the key acquisition and cohort parameters.
\needspace{8\baselineskip}
{\tablestripes\tablefontsize
\begin{xltabular}{\textwidth}{@{} L R @{}}
\caption[Depression dataset parameters]{Depression Dataset: Key Acquisition and Cohort Parameters}\label{tab:app_f_depression_params} \\
\toprule
\thead{Parameter} & \thead{Value} \\
\midrule
\endfirsthead
\multicolumn{2}{@{}l}{\textit{(\,Tab.~\ref{tab:app_f_depression_params} continued from previous page\,)}} \\
\toprule
\thead{Parameter} & \thead{Value} \\
\midrule
\endhead
\bottomrule
\endlastfoot
Source repository & OpenNeuro ds003478 \\
Total subjects recruited & 122 \\
Sampling rate & 256~Hz \\
Exclusion criterion & BDI scores 7--17 (ambiguous) \\
Subjects retained & 112 \\
Healthy controls (BDI $\leq$ 6) & 74 \\
Depressed patients (BDI $\geq$ 18) & 38 \\
Segments per subject & $\approx$16 \\
Total segments & 1,792 \\
Class ratio (healthy:depressed) & 1.95:1 \\
\end{xltabular}}
\vspace{0.5em}\noindent\textit{Note.} BDI = Beck Depression Inventory. Subjects with BDI scores 7--17 were excluded to create clean binary separation between healthy and moderate-to-severe depression groups.

Depression manifests in EEG through three principal biomarker alterations:

\begin{itemize}[leftmargin=*, itemsep=3pt]
\item \textbf{Increased frontal alpha asymmetry:} Right-hemisphere alpha power exceeds left-hemisphere alpha power, reflecting withdrawal-related prefrontal hypoactivation consistent with Davidson's approach--withdrawal model.
\item \textbf{Elevated theta power:} Widespread theta-band (4--8~Hz) power increase indicates rumination-linked cortical slowing and impaired cognitive control.
\item \textbf{Reduced beta coherence:} Diminished inter-hemispheric beta-band (13--30~Hz) coherence reflects disrupted functional connectivity in frontocentral networks.
\end{itemize}

\noindent The class imbalance (1.95:1) is the most pronounced in the study and requires SMOTE augmentation (40$\times$) combined with deep neural network classification. The agenticfinder pipeline achieves 91.07\% accuracy, the most challenging domain, reflecting the subtlety of depression-related EEG biomarkers.

\subsection{Sample Records (10 Rows)}

Table~\ref{tab:sample_depression} documents a representative 10-record sample from the depression dataset, showing EEG features alongside Beck Depression Inventory scores for both healthy and depressed subjects.
\needspace{3\baselineskip}
{\tablestripes\tablefontsize
\begin{xltabular}{\textwidth}{@{} C C R R R R R C @{}}
\caption[Depression dataset: 10-record sample]{Depression Dataset: Sample of 10 Records Showing EEG Features and BDI Scores}\label{tab:sample_depression} \\
\toprule
\thead{ID} & \thead{BDI} & \thead{$\alpha$ Asym.} & \thead{$\theta$ Pwr} & \thead{$\beta$ Coh.} & \thead{Entropy} & \thead{Hjorth} & \thead{Label} \\
\midrule
\endfirsthead
\multicolumn{8}{@{}l}{\textit{(\,Tab.~\ref{tab:sample_depression} continued from previous page\,)}} \\
\toprule
\thead{ID} & \thead{BDI} & \thead{$\alpha$ Asym.} & \thead{$\theta$ Pwr} & \thead{$\beta$ Coh.} & \thead{Entropy} & \thead{Hjorth} & \thead{Label} \\
\midrule
\endhead
\bottomrule
\endlastfoot
sub-01 & 3 & 0.087 & 2.94 & 0.682 & 0.798 & 0.208 & 0 (Healthy) \\
sub-12 & 1 & 0.102 & 2.78 & 0.695 & 0.812 & 0.198 & 0 (Healthy) \\
sub-28 & 5 & 0.074 & 3.08 & 0.671 & 0.785 & 0.215 & 0 (Healthy) \\
sub-45 & 2 & 0.095 & 2.85 & 0.689 & 0.805 & 0.202 & 0 (Healthy) \\
sub-67 & 4 & 0.081 & 2.98 & 0.678 & 0.792 & 0.211 & 0 (Healthy) \\
sub-08 & 24 & $-$0.142 & 4.87 & 0.498 & 0.658 & 0.287 & 1 (Depressed) \\
sub-19 & 31 & $-$0.178 & 5.21 & 0.472 & 0.634 & 0.298 & 1 (Depressed) \\
sub-35 & 22 & $-$0.128 & 4.62 & 0.512 & 0.672 & 0.278 & 1 (Depressed) \\
sub-52 & 28 & $-$0.165 & 5.05 & 0.485 & 0.645 & 0.292 & 1 (Depressed) \\
sub-74 & 19 & $-$0.115 & 4.45 & 0.524 & 0.681 & 0.271 & 1 (Depressed) \\
\end{xltabular}}
\vspace{0.5em}\noindent\textit{Note.} BDI = Beck Depression Inventory score; $\alpha$ Asym.\ = frontal alpha asymmetry; $\beta$ Coh.\ = beta-band coherence. Subjects with BDI 7--17 (mild/borderline) were excluded to create clean binary separation. Depressed subjects show reversed alpha asymmetry, 65\% elevated theta power, and 28\% reduced beta coherence.

\subsection{EDA Summary}

Table~\ref{tab:eda_depression} summarises the exploratory data analysis results for the depression dataset, highlighting the moderate negative skew in alpha asymmetry driven by the preponderance of healthy controls.
\needspace{3\baselineskip}
{\tablestripes\tablefontsize
\begin{xltabular}{\textwidth}{@{} L r r r r r r @{}}
\caption[Depression EDA summary]{Depression Dataset: Exploratory Data Analysis Summary}\label{tab:eda_depression} \\
\toprule
\thead{Feature} & \thead{Mean} & \thead{SD} & \thead{Median} & \thead{IQR} & \thead{Skew.} & \thead{Kurt.} \\
\midrule
\endfirsthead
\multicolumn{7}{@{}l}{\textit{(\,Tab.~\ref{tab:eda_depression} continued from previous page\,)}} \\
\toprule
\thead{Feature} & \thead{Mean} & \thead{SD} & \thead{Median} & \thead{IQR} & \thead{Skew.} & \thead{Kurt.} \\
\midrule
\endhead
\bottomrule
\endlastfoot
Alpha asymmetry & $-$0.012 & 0.138 & 0.015 & 0.187 & $-$0.52 & $-$0.28 \\
Theta power ($\mu$V$^2$/Hz) & 3.68 & 1.12 & 3.42 & 1.54 & 0.67 & $-$0.15 \\
Beta coherence & 0.598 & 0.112 & 0.608 & 0.152 & $-$0.34 & $-$0.58 \\
Spectral entropy & 0.738 & 0.082 & 0.742 & 0.112 & $-$0.18 & $-$0.42 \\
Hjorth mobility & 0.245 & 0.048 & 0.242 & 0.065 & 0.22 & $-$0.31 \\
Delta power ($\mu$V$^2$/Hz) & 5.42 & 1.87 & 5.12 & 2.58 & 0.48 & $-$0.22 \\
Gamma power ($\mu$V$^2$/Hz) & 0.87 & 0.32 & 0.84 & 0.43 & 0.35 & $-$0.28 \\
\end{xltabular}}
\vspace{0.5em}\noindent\textit{Note.} Moderate negative skew in alpha asymmetry ($-$0.52) reflects the preponderance of healthy controls (66\%) with positive asymmetry. Theta power shows positive skew (0.67) driven by the depressed subgroup. Class imbalance (1.95:1) is addressed through SMOTE augmentation.

\subsection{Key Observations}

Table~\ref{tab:depression_key_obs} summarises the principal findings from the depression dataset, highlighting its status as the most challenging classification domain and the mitigations applied.
\needspace{8\baselineskip}
{\tablestripes\tablefontsize
\begin{xltabular}{\textwidth}{@{} L c L @{}}
\caption[Depression key observations]{Depression Dataset: Key Observations}\label{tab:depression_key_obs} \\
\toprule
\thead{Observation} & \thead{Value} & \thead{Implication} \\
\midrule
\endfirsthead
\multicolumn{3}{@{}l}{\textit{(\,Tab.~\ref{tab:depression_key_obs} continued from previous page\,)}} \\
\toprule
\thead{Observation} & \thead{Value} & \thead{Implication} \\
\midrule
\endhead
\bottomrule
\endlastfoot
Classification accuracy & 91.07\% & Most challenging domain; subtler EEG biomarkers \\
\addlinespace
Class imbalance & 1.95:1 (74H : 38D) & Most severe imbalance; 40$\times$ SMOTE applied \\
\addlinespace
BDI exclusion (scores 7--17) & $-$8\% sample reduction & Clean binary separation at cost of sample size \\
\addlinespace
Top discriminative features & Alpha asymmetry, $\theta$/$\beta$ ratio & Both have SHAP importance $>$0.15 \\
\end{xltabular}}
\vspace{0.5em}\noindent\textit{Note.} H = healthy; D = depressed; BDI = Beck Depression Inventory; SMOTE = Synthetic Minority Over-sampling Technique; SHAP = SHapley Additive exPlanations.

%% ═══════════════════════════════════════════════════════════════════════
%% F.8 — CANCER DATASET
%% ═══════════════════════════════════════════════════════════════════════
% [SPACE-FIX] clearpage removed to eliminate empty page
\section{Dataset 7: Cancer (TCIA / Kaggle PBS Images)}
\label{app:data_cancer}

\subsection{Dataset Description}

The cancer dataset differs fundamentally from the EEG domains: it comprises 20,000 peripheral blood smear (PBS) images at 512$\times$512~RGB resolution from The Cancer Imaging Archive (TCIA) and supplementary Kaggle repositories. The dataset covers multiple cancer types with binary labelling: normal blood cells vs.\ cancerous/abnormal cells. Unlike EEG data, cancer classification uses radiomics features extracted from medical images rather than temporal signals.

The agenticfinder cancer detection module applies a GAN-based (generative adversarial network) data augmentation pipeline to address class imbalance, followed by CNN and CNN-LSTM architectures for classification. Feature extraction uses PyRadiomics for standardised radiomic feature computation across five feature families:

\begin{itemize}[leftmargin=*, itemsep=3pt]
\item \textbf{GLCM (texture):} Grey-level co-occurrence matrix features capturing spatial relationships between pixel intensities, including contrast, correlation, energy, and homogeneity.
\item \textbf{GLRLM (run-length):} Grey-level run-length matrix features quantifying consecutive pixels of the same intensity along a given direction, measuring texture coarseness.
\item \textbf{GLSZM (size-zone):} Grey-level size-zone matrix features characterising the size and distribution of connected regions of identical intensity values.
\item \textbf{First-order statistics:} Intensity histogram-based features including mean, variance, skewness, kurtosis, energy, and entropy computed from the voxel intensity distribution.
\item \textbf{3D shape descriptors:} Geometric features describing tumour morphology, including volume, surface area, sphericity, and compactness.
\end{itemize}

\subsection{Sample Records (10 Rows)}

Table~\ref{tab:sample_cancer} presents a representative 10-record sample from the cancer dataset, comparing radiomics features between normal and cancerous blood smear images.
\needspace{3\baselineskip}
{\tablestripes\tablefontsize
\begin{xltabular}{\textwidth}{@{} C R R R R R R C @{}}
\caption[Cancer dataset: 10-record sample]{Cancer Dataset: Sample of 10 Records Showing Radiomics Features}\label{tab:sample_cancer} \\
\toprule
\thead{Img.} & \thead{Mean Int.} & \thead{GLCM} & \thead{GLRLM} & \thead{Shape} & \thead{Entropy} & \thead{Compact.} & \thead{Label} \\
\midrule
\endfirsthead
\multicolumn{8}{@{}l}{\textit{(\,Tab.~\ref{tab:sample_cancer} continued from previous page\,)}} \\
\toprule
\thead{Img.} & \thead{Mean Int.} & \thead{GLCM} & \thead{GLRLM} & \thead{Shape} & \thead{Entropy} & \thead{Compact.} & \thead{Label} \\
\midrule
\endhead
\bottomrule
\endlastfoot
N0001 & 142.3 & 0.287 & 0.412 & 0.823 & 4.21 & 0.912 & 0 (Normal) \\
N0248 & 138.7 & 0.298 & 0.398 & 0.835 & 4.15 & 0.918 & 0 (Normal) \\
N0512 & 145.1 & 0.275 & 0.425 & 0.818 & 4.28 & 0.908 & 0 (Normal) \\
N0891 & 140.5 & 0.291 & 0.408 & 0.828 & 4.19 & 0.915 & 0 (Normal) \\
N1200 & 143.8 & 0.282 & 0.418 & 0.821 & 4.24 & 0.911 & 0 (Normal) \\
C0001 & 168.9 & 0.412 & 0.587 & 0.654 & 5.42 & 0.768 & 1  \\
C0315 & 175.2 & 0.438 & 0.612 & 0.628 & 5.67 & 0.745 & 1  \\
C0672 & 171.4 & 0.425 & 0.598 & 0.641 & 5.51 & 0.758 & 1  \\
C0998 & 178.8 & 0.451 & 0.624 & 0.615 & 5.78 & 0.732 & 1  \\
C1450 & 173.6 & 0.432 & 0.605 & 0.635 & 5.58 & 0.752 & 1  \\
\end{xltabular}}
\vspace{0.5em}\noindent\textit{Note.} Img.\ = image ID; Mean Int.\ = mean pixel intensity (0--255); GLCM = gray-level co-occurrence matrix contrast; GLRLM = gray-level run-length matrix short-run emphasis; Shape = sphericity; Compact.\ = compactness. Cancer cells show higher intensity, increased textural heterogeneity (GLCM, GLRLM), and reduced sphericity/compactness.

\subsection{EDA Summary}

Table~\ref{tab:eda_cancer} reports the exploratory data analysis summary for the cancer radiomics dataset, showing near-symmetric distributions across all texture, shape, and intensity features.
\needspace{3\baselineskip}
{\tablestripes\tablefontsize
\begin{xltabular}{\textwidth}{@{} L r r r r r r @{}}
\caption[Cancer EDA summary]{Cancer Dataset: Exploratory Data Analysis Summary}\label{tab:eda_cancer} \\
\toprule
\thead{Feature} & \thead{Mean} & \thead{SD} & \thead{Median} & \thead{IQR} & \thead{Skew.} & \thead{Kurt.} \\
\midrule
\endfirsthead
\multicolumn{7}{@{}l}{\textit{(\,Tab.~\ref{tab:eda_cancer} continued from previous page\,)}} \\
\toprule
\thead{Feature} & \thead{Mean} & \thead{SD} & \thead{Median} & \thead{IQR} & \thead{Skew.} & \thead{Kurt.} \\
\midrule
\endhead
\bottomrule
\endlastfoot
Mean intensity & 155.6 & 18.4 & 154.2 & 25.8 & 0.18 & $-$0.72 \\
GLCM contrast & 0.358 & 0.085 & 0.348 & 0.115 & 0.42 & $-$0.35 \\
GLRLM short-run & 0.508 & 0.108 & 0.498 & 0.148 & 0.38 & $-$0.28 \\
Sphericity & 0.738 & 0.102 & 0.745 & 0.138 & $-$0.28 & $-$0.52 \\
Entropy & 4.87 & 0.72 & 4.78 & 0.98 & 0.35 & $-$0.18 \\
Compactness & 0.842 & 0.092 & 0.848 & 0.125 & $-$0.32 & $-$0.45 \\
First-order energy & 2.84e6 & 8.42e5 & 2.72e6 & 1.15e6 & 0.45 & $-$0.12 \\
\end{xltabular}}
\vspace{0.5em}\noindent\textit{Note.} Near-symmetric distributions for all features (skewness $<$ 0.5). The large sample size ($N$ = 20,000) ensures stable estimates. GLCM and GLRLM features show the strongest discriminative power between normal and cancerous cells.

\subsection{Key Observations}

Table~\ref{tab:cancer_key_obs} summarises the principal characteristics that distinguish the cancer domain from the EEG-based datasets, including the imaging pipeline, sample size advantages, and augmentation strategy.
\needspace{8\baselineskip}
{\tablestripes\tablefontsize
\begin{xltabular}{\textwidth}{@{} L c L @{}}
\caption[Cancer key observations]{Cancer Dataset: Key Observations}\label{tab:cancer_key_obs} \\
\toprule
\thead{Observation} & \thead{Value} & \thead{Implication} \\
\midrule
\endfirsthead
\multicolumn{3}{@{}l}{\textit{(\,Tab.~\ref{tab:cancer_key_obs} continued from previous page\,)}} \\
\toprule
\thead{Observation} & \thead{Value} & \thead{Implication} \\
\midrule
\endhead
\bottomrule
\endlastfoot
Modality & PBS images (not EEG) & Different pipeline: DICOM/PNG $\to$ ROI $\to$ radiomics \\
\addlinespace
Sample size & $N = 20{,}000$ & Stable DL training without extensive augmentation \\
\addlinespace
GAN augmentation & Applied per subtype & Addresses within-class imbalance for minority subtypes \\
\addlinespace
Architecture & CNN-LSTM & Captures spatial + sequential patterns; 99.02\% accuracy \\
\end{xltabular}}
\vspace{0.5em}\noindent\textit{Note.} PBS = peripheral blood smear; ROI = region of interest; DL = deep learning; GAN = generative adversarial network.

%% ═══════════════════════════════════════════════════════════════════════
%% F.9 — FEATURE/ATTRIBUTE DICTIONARY
%% ═══════════════════════════════════════════════════════════════════════
% [SPACE-FIX] clearpage removed to eliminate empty page
\section{Feature and Attribute Dictionary}
\label{app:feature_dictionary}

\subsection{EEG Feature Dictionary}

Table~\ref{tab:eeg_feature_dict} catalogues all extracted EEG features with their mathematical definitions, units, and clinical significance across the six neurological domains.

\needspace{3\baselineskip}
\begingroup\singlespacing\tablefontsize
\begin{xltabular}
{\textwidth}{@{} p{2.5cm} p{1.5cm} L p{1.5cm} L @{}}
\caption[EEG feature dictionary]{EEG Feature Dictionary: All Extracted Features with Definitions and Clinical Significance}
\label{tab:eeg_feature_dict}\\
\toprule
\thead{Feature} & \thead{Category} & \thead{Definition / Formula} & \thead{Unit} & \thead{Clinical Significance} \\
\midrule
\endfirsthead
\endhead
\bottomrule
\endlastfoot
\multicolumn{5}{l}{\textbf{Time-Domain Features}} \\
\addlinespace
Mean amplitude & Time & $\bar{x} = \frac{1}{N}\sum_{i=1}^{N} x_i$ & $\mu$V & Baseline signal level; elevated in seizure \\
\addlinespace
Variance & Time & $\sigma^2 = \frac{1}{N}\sum (x_i - \bar{x})^2$ & $\mu$V$^2$ & Signal variability; higher in pathological states \\
\addlinespace
Standard deviation & Time & $\sigma = \sqrt{\text{Var}(x)}$ & $\mu$V & Spread of signal amplitude \\
\addlinespace
Zero-crossing rate & Time & $\text{ZCR} = \frac{1}{N-1}\sum |{\text{sgn}(x_i) - \text{sgn}(x_{i-1})}|$ & --- & Frequency proxy; higher in fast oscillations \\
\addlinespace
Line length & Time & $\text{LL} = \sum |x_i - x_{i-1}|$ & $\mu$V & Signal complexity; spikes in seizure onset \\
\addlinespace
Kurtosis & Time & $K = \frac{E[(x - \mu)^4]}{\sigma^4}$ & --- & Peakedness; high values indicate sharp transients \\
\addlinespace
Skewness & Time & $S = \frac{E[(x - \mu)^3]}{\sigma^3}$ & --- & Asymmetry; non-zero in pathological EEG \\
\addlinespace
Hjorth activity & Time & $H_0 = \text{Var}(x)$ & $\mu$V$^2$ & Signal power \\
\addlinespace
Hjorth mobility & Time & $H_1 = \sqrt{\text{Var}(x')/\text{Var}(x)}$ & Hz & Mean frequency of the signal \\
\addlinespace
Hjorth complexity & Time & $H_2 = H_1(x')/H_1(x)$ & --- & Bandwidth; higher = more complex signal \\
\addlinespace[8pt]
\multicolumn{5}{l}{\textbf{Frequency-Domain Features}} \\
\addlinespace
Delta power & Freq. & $\int_{0.5}^{4} S(f)\,df$ & $\mu$V$^2$/Hz & Deep sleep, pathological slowing (PD, epilepsy) \\
\addlinespace
Theta power & Freq. & $\int_{4}^{8} S(f)\,df$ & $\mu$V$^2$/Hz & Drowsiness, meditation; elevated in ASD, schizo. \\
\addlinespace
Alpha power & Freq. & $\int_{8}^{13} S(f)\,df$ & $\mu$V$^2$/Hz & Relaxed wakefulness; reduced in depression, PD \\
\addlinespace
Beta power & Freq. & $\int_{13}^{30} S(f)\,df$ & $\mu$V$^2$/Hz & Active thinking; desynchronised in PD \\
\addlinespace
Gamma power & Freq. & $\int_{30}^{45} S(f)\,df$ & $\mu$V$^2$/Hz & Cognitive processing; reduced in schizophrenia \\
\addlinespace
Spectral entropy & Freq. & $H = -\sum p(f)\log p(f)$ & bits & Signal complexity; lower in seizure \\
\addlinespace
Spectral centroid & Freq. & $\bar{f} = \sum f \cdot S(f) / \sum S(f)$ & Hz & Weighted mean frequency \\
\addlinespace
Alpha peak freq. & Freq. & $f_{\alpha} = \arg\max_{8 \leq f \leq 13} S(f)$ & Hz & Individual alpha frequency; slowed in PD \\
\addlinespace[8pt]
\multicolumn{5}{l}{\textbf{Connectivity Features}} \\
\addlinespace
Coherence & Conn. & $C_{xy}(f) = |S_{xy}|^2 / (S_{xx} \cdot S_{yy})$ & --- & Linear coupling between channels \\
\addlinespace
Phase-locking value & Conn. & $\text{PLV} = |\frac{1}{N}\sum e^{j(\phi_x - \phi_y)}|$ & --- & Phase synchronisation; reduced in ASD \\
\addlinespace
Mutual information & Conn. & $\text{MI} = \sum p(x,y)\log\frac{p(x,y)}{p(x)p(y)}$ & bits & Non-linear dependency \\
\addlinespace
Granger causality & Conn. & $F$-test on autoregressive prediction error & --- & Directional information flow \\
\addlinespace
Alpha asymmetry & Spatial & $\ln(\alpha_R) - \ln(\alpha_L)$ & --- & Frontal asymmetry; reversed in depression/stress \\
\end{xltabular}
\endgroup

\subsection{Cancer/Imaging Feature Dictionary}

Table~\ref{tab:cancer_feature_dict} enumerates the radiomics features extracted from peripheral blood smear images, including texture, shape, and intensity descriptors with their clinical significance for cancer classification.
\needspace{3\baselineskip}
{\tablestripes\tablefontsize
\begin{xltabular}{\textwidth}{@{} p{2.5cm} p{1.5cm} L L @{}}
\caption[Cancer radiomics feature dictionary]{Cancer Radiomics Feature Dictionary}\label{tab:cancer_feature_dict} \\
\toprule
\thead{Feature} & \thead{Category} & \thead{Definition} & \thead{Clinical Significance} \\
\midrule
\endfirsthead
\multicolumn{4}{@{}l}{\textit{(\,Tab.~\ref{tab:cancer_feature_dict} continued from previous page\,)}} \\
\toprule
\thead{Feature} & \thead{Category} & \thead{Definition} & \thead{Clinical Significance} \\
\midrule
\endhead
\bottomrule
\endlastfoot
GLCM contrast & Texture & Co-occurrence of pixel intensities at spatial offsets & Textural heterogeneity; higher in malignant tissue \\
\addlinespace
GLCM homogeneity & Texture & Closeness of GLCM element distribution to diagonal & Uniformity; reduced in cancer \\
\addlinespace
GLRLM short-run & Texture & Distribution of consecutive pixels with same intensity & Fine texture patterns; discriminates cell types \\
\addlinespace
GLSZM zone entropy & Texture & Size-zone matrix randomness & Structural complexity of tissue architecture \\
\addlinespace
Sphericity & Shape & How closely a 3D region approximates a sphere & Cell morphology; cancer cells less spherical \\
\addlinespace
Compactness & Shape & Surface-area-to-volume ratio & Cell regularity; lower in malignant cells \\
\addlinespace
First-order energy & Intensity & $\sum x_i^2$ & Overall signal magnitude; elevated in dense tissue \\
\addlinespace
First-order entropy & Intensity & $-\sum p(x)\log p(x)$ & Randomness of pixel intensity distribution \\
\end{xltabular}}
\vspace{0.5em}\noindent\textit{Note.} GLCM = gray-level co-occurrence matrix; GLRLM = gray-level run-length matrix; GLSZM = gray-level size-zone matrix. Features computed using the PyRadiomics library (v3.0+).

\subsection{Stress Wearable Sensor Feature Dictionary}

Table~\ref{tab:wearable_feature_dict} maps each wearable sensor feature to its source sensor, definition, and expected physiological response under stress conditions.
\needspace{3\baselineskip}
{\tablestripes\tablefontsize
\begin{xltabular}{\textwidth}{@{} p{2.3cm} p{1.3cm} L L @{}}
\caption[Wearable sensor feature dictionary]{Stress Wearable Sensor Feature Dictionary (Agenticfinder Neuro-Wearable System)}\label{tab:wearable_feature_dict} \\
\toprule
\thead{Feature} & \thead{Sensor} & \thead{Definition} & \thead{Stress Response} \\
\midrule
\endfirsthead
\multicolumn{4}{@{}l}{\textit{(\,Tab.~\ref{tab:wearable_feature_dict} continued from previous page\,)}} \\
\toprule
\thead{Feature} & \thead{Sensor} & \thead{Definition} & \thead{Stress Response} \\
\midrule
\endhead
\bottomrule
\endlastfoot
HRV (RMSSD) & ECG & Root mean square of successive RR interval differences & Decreased under stress (vagal withdrawal) \\
\addlinespace
HR mean & ECG & Mean heart rate (bpm) & Elevated during acute stress \\
\addlinespace
SCR amplitude & EDA & Skin conductance response peak & Increased with sympathetic arousal \\
\addlinespace
SCL mean & EDA & Tonic skin conductance level ($\mu$S) & Elevated baseline under chronic stress \\
\addlinespace
SpO$_2$ & PPG & Blood oxygen saturation (\%) & Minimal change; control variable \\
\addlinespace
Skin temperature & Temp. & Peripheral skin temperature (\textdegree{}C) & Increased 0.5--1.0\textdegree{}C under stress \\
\addlinespace
Accelerometer RMS & IMU & Root mean square of 3-axis acceleration & Motion artefact indicator; quality control \\
\end{xltabular}}
\vspace{0.5em}\noindent\textit{Note.} ECG = electrocardiography; EDA = electrodermal activity; PPG = photoplethysmography; IMU = inertial measurement unit; RMSSD = root mean square of successive differences; SCR = skin conductance response; SCL = skin conductance level.

%% ═══════════════════════════════════════════════════════════════════════
%% F.10 — EEG FREQUENCY BAND ANALYSIS
%% ═══════════════════════════════════════════════════════════════════════
% [SPACE-FIX] clearpage removed to eliminate empty page
\section{EEG Frequency Band Analysis}
\label{app:frequency_bands}

\subsection{Five-Band Decomposition}

All EEG feature extraction in this dissertation decomposes the raw signal into five canonical frequency bands. Figure~\ref{fig:frequency_bands} summarises each band's frequency range, physiological correlate, and the disease domains in which it shows the most pronounced alterations.
\needspace{3\baselineskip}
\begin{figure}[!htbp]
\centering
\resizebox{0.97\textwidth}{!}{%
\begin{tikzpicture}[
  band/.style={rectangle, rounded corners=3pt, minimum height=1.2cm, text width=12cm, align=left, font=\small, inner sep=6pt},
  lbl/.style={font=\small\bfseries, text=white, minimum width=2.2cm, align=center}
]
%% Delta
\node[band, fill=lightblue!40] (d) at (0,0) {};
\node[lbl, fill=lightblue!80!black] at (-5.4,0) {$\delta$ 0.5--4~Hz};
\node[font=\small, text width=9.2cm, right] at (-3.8,0) {Deep sleep, pathological slowing. Elevated in epilepsy seizure onset.};
%% Theta
\node[band, fill=lightorange!40] (t) at (0,-1.6) {};
\node[lbl, fill=lightorange!80!black] at (-5.4,-1.6) {$\theta$ 4--8~Hz};
\node[font=\small, text width=9.2cm, right] at (-3.8,-1.6) {Drowsiness, memory encoding. Elevated in ASD (+75\%), schizophrenia (+61\%).};
%% Alpha
\node[band, fill=lightteal!40] (a) at (0,-3.2) {};
\node[lbl, fill=lightteal!80!black] at (-5.4,-3.2) {$\alpha$ 8--13~Hz};
\node[font=\small, text width=9.2cm, right] at (-3.8,-3.2) {Relaxed wakefulness, thalamocortical loop. Reduced in PD, ASD, depression.};
%% Beta
\node[band, fill=lightpurple!40] (b) at (0,-4.8) {};
\node[lbl, fill=lightpurple!80!black] at (-5.4,-4.8) {$\beta$ 13--30~Hz};
\node[font=\small, text width=9.2cm, right] at (-3.8,-4.8) {Active thinking, motor planning. Desynchronised in PD; elevated in stress.};
%% Gamma
\node[band, fill=lightgreen!40] (g) at (0,-6.4) {};
\node[lbl, fill=lightgreen!80!black] at (-5.4,-6.4) {$\gamma$ 30--45~Hz};
\node[font=\small, text width=9.2cm, right] at (-3.8,-6.4) {Cognitive binding, perception. Reduced synchrony in schizophrenia.};
\end{tikzpicture}%
}%
\caption[EEG frequency band decomposition]{EEG Frequency Band Decomposition: Five Standard Bands with Clinical Significance and Disease-Specific Alterations}
\label{fig:frequency_bands}
\end{figure}

\subsection{Band Power Distribution by Domain}

Table~\ref{tab:band_power_domain} compares the mean spectral band power across all six EEG disease domains, revealing that epilepsy dominates total power and the theta/alpha ratio serves as the primary ASD-focused discriminative index.
\needspace{3\baselineskip}
{\tablestripes\tablefontsize
\begin{xltabular}{\textwidth}{@{} L r r r r r r @{}}
\caption[Band power by domain]{Mean Spectral Band Power ($\mu$V$^2$/Hz) by Disease Domain}\label{tab:band_power_domain} \\
\toprule
\thead{Band} & \thead{Schizo.} & \thead{Epilepsy} & \thead{Stress} & \thead{ASD} & \thead{PD} & \thead{Depress.} \\
\midrule
\endfirsthead
\multicolumn{7}{@{}l}{\textit{(\,Tab.~\ref{tab:band_power_domain} continued from previous page\,)}} \\
\toprule
\thead{Band} & \thead{Schizo.} & \thead{Epilepsy} & \thead{Stress} & \thead{ASD} & \thead{PD} & \thead{Depress.} \\
\midrule
\endhead
\bottomrule
\endlastfoot
$\delta$ (0.5--4~Hz) & 5.82 & 28.92 & 4.12 & 4.87 & 6.17 & 5.42 \\
$\theta$ (4--8~Hz) & 3.72 & 8.47 & 3.45 & 3.38 & 3.92 & 3.68 \\
$\alpha$ (8--13~Hz) & 3.21 & 4.15 & 3.98 & 4.45 & 3.51 & 4.12 \\
$\beta$ (13--30~Hz) & 1.87 & 5.87 & 3.98 & 2.12 & 2.48 & 2.15 \\
$\gamma$ (30--45~Hz) & 0.68 & 1.42 & 0.95 & 0.82 & 0.74 & 0.87 \\
\midrule
\textbf{Total power} & \textbf{15.30} & \textbf{48.83} & \textbf{16.48} & \textbf{15.64} & \textbf{16.82} & \textbf{16.24} \\
$\theta$/$\alpha$ ratio & 1.16 & 2.04 & 0.87 & 0.76 & 1.12 & 0.89 \\
\end{xltabular}}
\vspace{0.5em}\noindent\textit{Note.} Values represent mean band power across all subjects (healthy + disease combined). Epilepsy shows dramatically elevated total power due to seizure episodes. The $\theta$/$\alpha$ ratio is the primary discriminative index: elevated in schizophrenia and PD, reduced in stress and ASD.

\subsection{Band Power Heatmap}

Figure~\ref{fig:band_power_heatmap} renders the same band-by-domain data from Table~\ref{tab:band_power_domain} as a colour-coded heatmap, enabling rapid visual identification of the domains with the highest power in each frequency band.
\needspace{20\baselineskip}
\begin{figure}[!htbp]
\centering
\resizebox{0.97\textwidth}{!}{%
\begin{tikzpicture}[
  cell/.style={minimum width=2cm, minimum height=0.9cm, align=center, font=\small}
]
%% Column headers
\foreach \x/\lab in {0/Schizo., 1/Epilepsy, 2/Stress, 3/ASD, 4/PD, 5/Depress.} {
  \node[cell, font=\small\bfseries] at (\x*2.2+1.1, 5.4) {\lab};
}
%% Row headers
\foreach \y/\lab in {0/$\gamma$, 1/$\beta$, 2/$\alpha$, 3/$\theta$, 4/$\delta$} {
  \node[cell, font=\small\bfseries, anchor=east] at (-0.2, \y*1.0+0.5) {\lab};
}
%% Data cells (colour intensity = relative power, normalised 0-1 per band)
%% Color scheme: green!XX for low, lightyellow!XX for medium, lightorange!XX for high, lightcoral!XX for very high
%% Delta row (y=4): values [5.82, 28.92, 4.12, 4.87, 6.17, 5.42] → normalised
\node[cell, fill=lightgreen!30, text=black] at (0*2.2+1.1, 4.5) {5.82};
\node[cell, fill=lightcoral!80, text=black] at (1*2.2+1.1, 4.5) {28.9};
\node[cell, fill=lightgreen!20, text=black] at (2*2.2+1.1, 4.5) {4.12};
\node[cell, fill=lightgreen!25, text=black] at (3*2.2+1.1, 4.5) {4.87};
\node[cell, fill=lightgreen!35, text=black] at (4*2.2+1.1, 4.5) {6.17};
\node[cell, fill=lightgreen!25, text=black] at (5*2.2+1.1, 4.5) {5.42};
%% Theta row (y=3): [3.72, 8.47, 3.45, 3.38, 3.92, 3.68]
\node[cell, fill=lightyellow!50, text=black] at (0*2.2+1.1, 3.5) {3.72};
\node[cell, fill=lightcoral!80, text=black] at (1*2.2+1.1, 3.5) {8.47};
\node[cell, fill=lightyellow!40, text=black] at (2*2.2+1.1, 3.5) {3.45};
\node[cell, fill=lightyellow!40, text=black] at (3*2.2+1.1, 3.5) {3.38};
\node[cell, fill=lightyellow!55, text=black] at (4*2.2+1.1, 3.5) {3.92};
\node[cell, fill=lightyellow!45, text=black] at (5*2.2+1.1, 3.5) {3.68};
%% Alpha row (y=2): [3.21, 4.15, 3.98, 4.45, 3.51, 4.12]
\node[cell, fill=lightorange!60, text=black] at (0*2.2+1.1, 2.5) {3.21};
\node[cell, fill=lightcoral!75, text=black] at (1*2.2+1.1, 2.5) {4.15};
\node[cell, fill=lightorange!70, text=black] at (2*2.2+1.1, 2.5) {3.98};
\node[cell, fill=lightcoral!80, text=black] at (3*2.2+1.1, 2.5) {4.45};
\node[cell, fill=lightorange!60, text=black] at (4*2.2+1.1, 2.5) {3.51};
\node[cell, fill=lightcoral!72, text=black] at (5*2.2+1.1, 2.5) {4.12};
%% Beta row (y=1): [1.87, 5.87, 3.98, 2.12, 2.48, 2.15]
\node[cell, fill=lightyellow!35, text=black] at (0*2.2+1.1, 1.5) {1.87};
\node[cell, fill=lightcoral!80, text=black] at (1*2.2+1.1, 1.5) {5.87};
\node[cell, fill=lightorange!60, text=black] at (2*2.2+1.1, 1.5) {3.98};
\node[cell, fill=lightyellow!40, text=black] at (3*2.2+1.1, 1.5) {2.12};
\node[cell, fill=lightyellow!50, text=black] at (4*2.2+1.1, 1.5) {2.48};
\node[cell, fill=lightyellow!40, text=black] at (5*2.2+1.1, 1.5) {2.15};
%% Gamma row (y=0): [0.68, 1.42, 0.95, 0.82, 0.74, 0.87]
\node[cell, fill=lightyellow!50, text=black] at (0*2.2+1.1, 0.5) {0.68};
\node[cell, fill=lightcoral!80, text=black] at (1*2.2+1.1, 0.5) {1.42};
\node[cell, fill=lightorange!60, text=black] at (2*2.2+1.1, 0.5) {0.95};
\node[cell, fill=lightyellow!55, text=black] at (3*2.2+1.1, 0.5) {0.82};
\node[cell, fill=lightyellow!50, text=black] at (4*2.2+1.1, 0.5) {0.74};
\node[cell, fill=lightorange!55, text=black] at (5*2.2+1.1, 0.5) {0.87};
%% Grid lines
\draw[ggunavy!30] (-0.1,0) grid[xstep=2.2, ystep=1.0] (13.3,5);
\end{tikzpicture}%
}%
\caption[Band power heatmap across domains]{Band Power Heatmap: Mean Spectral Power ($\mu$V$^2$/Hz) by Frequency Band and Disease Domain. Darker cells indicate higher relative power within each band. Epilepsy dominates all bands due to seizure-related hypersynchronisation.}
\label{fig:band_power_heatmap}
\end{figure}

%% ═══════════════════════════════════════════════════════════════════════
%% F.11 — FEATURE ENGINEERING PIPELINE
%% ═══════════════════════════════════════════════════════════════════════
% [SPACE-FIX] clearpage removed to eliminate empty page
\section{Feature Engineering Pipeline}
\label{app:feature_pipeline}

% MOVED TO CHAPTER 3 MAIN TEXT (Mar 2026)
\iffalse
\subsection{Eight-Stage Pipeline Overview}

\needspace{20\baselineskip}
\begin{figure}[!htbp]
\centering
\begin{tikzpicture}[
  stage/.style={rectangle, rounded corners=4pt, draw=ggunavy, fill=ggunavy!8, text width=3cm, minimum height=1.4cm, align=center, font=\small, inner sep=5pt},
  arr/.style={-{Stealth[length=4mm, width=3mm]}, very thick, ggunavy},
  note/.style={font=\small\itshape, text=ggunavy!60, text width=3cm, align=center}
]
%% Row 1
\node[stage, fill=lightblue] (s1) at (0,3) {\textbf{1. Raw Signal}\\{\small EEG / Image}\\{\small acquisition}};
\node[stage, fill=lightorange] (s2) at (4,3) {\textbf{2. Bandpass Filter}\\{\small 0.5--45~Hz}\\{\small Butterworth 4th}};
\node[stage, fill=lightorange] (s3) at (8,3) {\textbf{3. Artefact Removal}\\{\small ICA decomposition}\\{\small EOG/EMG reject}};
\node[stage, fill=lightteal] (s4) at (12,3) {\textbf{4. Segmentation}\\{\small Epoch windowing}\\{\small 50\% overlap}};
%% Row 2
\node[stage, fill=lightteal] (s5) at (12,0) {\textbf{5. Feature Extract.}\\{\small Time, freq.,}\\{\small connectivity}};
\node[stage, fill=lightpurple] (s6) at (8,0) {\textbf{6. Feature Select.}\\{\small MI, ANOVA, SHAP}\\{\small Top-$k$ ranking}};
\node[stage, fill=lightpurple] (s7) at (4,0) {\textbf{7. Feature Eval.}\\{\small Importance scores}\\{\small Stability check}};
\node[stage, fill=lightgreen] (s8) at (0,0) {\textbf{8. Normalisation}\\{\small Z-score / MinMax}\\{\small Domain-specific}};
%% Arrows
\draw[arr] (s1) -- (s2);
\draw[arr] (s2) -- (s3);
\draw[arr] (s3) -- (s4);
\draw[arr] (s4) -- ++(0,-1) -| (s5);
\draw[arr] (s5) -- (s6);
\draw[arr] (s6) -- (s7);
\draw[arr] (s7) -- (s8);
%% Output
\node[rectangle, rounded corners=3pt, draw=ggunavy, fill=ggunavy, text=white, text width=3cm, minimum height=1cm, align=center, font=\small\bfseries] (out) at (0,-2) {Model Input\\{\normalfont\scriptsize Feature matrix}};
\draw[arr] (s8) -- (out);
\end{tikzpicture}
\caption[Eight-stage feature engineering pipeline]{Eight-Stage Feature Engineering Pipeline: From Raw EEG/Image Acquisition Through Filtering, Artefact Removal, Segmentation, Feature Extraction, Selection, Evaluation, and Normalisation to Model-Ready Feature Matrices}
\label{fig:feature_pipeline}
\end{figure}
\fi

\subsection{Stage 1--2: Bandpass Filtering Parameters}

Table~\ref{tab:bandpass_params} documents the bandpass filtering parameters applied per domain, including cutoff frequencies, filter order, and the clinical rationale for each configuration.
\needspace{3\baselineskip}
{\tablestripes\tablefontsize
\begin{xltabular}{\textwidth}{@{} L c c c L @{}}
\caption[Bandpass filtering parameters]{Bandpass Filtering Parameters by Domain}\label{tab:bandpass_params} \\
\toprule
\thead{Domain} & \thead{Low Cut (Hz)} & \thead{High Cut (Hz)} & \thead{Order} & \thead{Rationale} \\
\midrule
\endfirsthead
\multicolumn{5}{@{}l}{\textit{(\,Tab.~\ref{tab:bandpass_params} continued from previous page\,)}} \\
\toprule
\thead{Domain} & \thead{Low Cut (Hz)} & \thead{High Cut (Hz)} & \thead{Order} & \thead{Rationale} \\
\midrule
\endhead
\bottomrule
\endlastfoot
Schizophrenia & 0.5 & 45 & 4th Butterworth & Full EEG spectrum; preserves gamma for binding analysis \\
\addlinespace
Epilepsy & 0.5 & 45 & 4th Butterworth & Seizure activity spans delta to beta bands \\
\addlinespace
Stress & 0.5 & 45 & 4th Butterworth & Alpha asymmetry and beta activation require full range \\
\addlinespace
Autism (ASD) & 1.0 & 40 & 4th Butterworth & Slightly narrower to reduce low-freq.\ drift in children \\
\addlinespace
Parkinson  & 0.5 & 45 & 4th Butterworth & Delta slowing is primary biomarker; must preserve 0.5~Hz \\
\addlinespace
Depression & 0.5 & 45 & 4th Butterworth & Full range for alpha asymmetry and theta/beta ratio \\
\addlinespace
Cancer (imaging) & --- & --- & --- & N/A: image data; histogram equalisation used instead \\
\end{xltabular}}
\vspace{0.5em}\noindent\textit{Note.} Butterworth filters chosen for maximally flat passband response. Zero-phase filtering (forward--backward) applied to prevent phase distortion. ASD uses 1.0~Hz low-cut to reduce slow-drift artefacts common in paediatric recordings.

\subsection{Stage 3}

Table~\ref{tab:data_cleaning} enumerates the artefact removal methods applied to each domain, including the types of artefacts targeted and the percentage of data rejected during quality control.
\needspace{3\baselineskip}
{\tablestripes\tablefontsize
\begin{xltabular}{\textwidth}{@{} p{2cm} L L c @{}}
\caption[Data cleaning methods]{Data Cleaning and Artefact Removal Methods by Domain}\label{tab:data_cleaning} \\
\toprule
\thead{Domain} & \thead{Method} & \thead{Artefacts Removed} & \thead{\% Rejected} \\
\midrule
\endfirsthead
\multicolumn{4}{@{}l}{\textit{(\,Tab.~\ref{tab:data_cleaning} continued from previous page\,)}} \\
\toprule
\thead{Domain} & \thead{Method} & \thead{Artefacts Removed} & \thead{\% Rejected} \\
\midrule
\endhead
\bottomrule
\endlastfoot
Schizo. & ICA decomposition (FastICA) & Eye blinks (EOG), muscle (EMG), line noise (50~Hz) & 8.2\% \\
\addlinespace
Epilepsy & ICA + amplitude thresholding ($\pm$500~$\mu$V) & EOG, EMG, electrode pop, movement & 5.7\% \\
\addlinespace
Stress & ICA + regression-based EOG removal & Eye blinks, jaw clench, sweat artefact & 7.4\% \\
\addlinespace
ASD & ICA + visual inspection protocol & EOG, EMG (child movement), electrode drift & 11.3\% \\
\addlinespace
PD & ICA + notch filter (50/60~Hz) & Tremor-related EMG, line noise, electrode artefact & 6.8\% \\
\addlinespace
Depression & ICA + automated ADJUST algorithm & EOG, cardiac artefact, electrode bridge & 9.1\% \\
\addlinespace
Cancer & Histogram equalisation + denoising & Image noise, staining inconsistency, artefact regions & 3.2\% \\
\end{xltabular}}
\vspace{0.5em}\noindent\textit{Note.} ICA = independent component analysis; EOG = electrooculography (eye movement); EMG = electromyography (muscle). ASD has the highest rejection rate (11.3\%) due to movement artefacts in paediatric recordings. All rejection rates are within acceptable bounds ($<$15\%).

\subsection{Stage 4: Segmentation and Windowing}

Table~\ref{tab:segmentation_params} reports the epoch segmentation and windowing parameters for each dataset, specifying window duration, sample count, overlap percentage, and the number of features extracted per segment.
\needspace{3\baselineskip}
{\tablestripes\tablefontsize
\begin{xltabular}{\textwidth}{@{} L c c c c c @{}}
\caption[Segmentation parameters]{Segmentation and Windowing Parameters by Dataset}\label{tab:segmentation_params} \\
\toprule
\thead{Dataset} & \thead{Window (s)} & \thead{Samples} & \thead{Overlap} & \thead{Segments} & \thead{Features/Seg.} \\
\midrule
\endfirsthead
\multicolumn{6}{@{}l}{\textit{(\,Tab.~\ref{tab:segmentation_params} continued from previous page\,)}} \\
\toprule
\thead{Dataset} & \thead{Window (s)} & \thead{Samples} & \thead{Overlap} & \thead{Segments} & \thead{Features/Seg.} \\
\midrule
\endhead
\bottomrule
\endlastfoot
Schizophrenia & 15.6 & 2,000 & 0\% & 8,400 & 45 \\
Epilepsy & 4.0 & 1,024 & 50\% & 5,100 & 52 \\
Stress (DEAP) & 4.0 & 2,048 & 50\% & 3,200 & 48 \\
Stress (SAM40) & 4.0 & 512 & 50\% & 2,800 & 48 \\
Autism & --- & Pre-extracted & --- & 3,000 & 38 \\
Parkinson & --- & Pre-extracted & --- & 500 & 42 \\
Depression & 4.0 & 1,024 & 50\% & 1,792 & 46 \\
Cancer & --- & 512$\times$512 pixels & --- & 20,000 & 107 \\
\end{xltabular}}
\vspace{0.5em}\noindent\textit{Note.} Overlap = percentage of window overlap for sliding-window segmentation. ``Pre-extracted'' indicates features were computed at the source. Cancer features computed per image (107 PyRadiomics features). The 50\% overlap increases training samples while preserving temporal context.

\subsection{Epoch Extraction Visualisation}

Figure~\ref{fig:epoch_extraction} illustrates the sliding-window segmentation process applied to continuous EEG recordings. Overlapping windows (50\%) increase the number of training samples while preserving temporal context at window boundaries.
\needspace{3\baselineskip}
\begin{figure}[!htbp]
\centering
\resizebox{0.97\textwidth}{!}{%
\begin{tikzpicture}[
  sig/.style={ggunavy, thick},
  win/.style={fill=ggunavy!10, draw=ggunavy, line width=0.5pt},
  arr/.style={-{Stealth[length=4mm, width=3mm]}, very thick, ggunavy},
  lbl/.style={font=\small, text=ggunavy}
]
%% Continuous EEG signal (simulated sine wave)
\draw[sig] plot[domain=0:12, samples=200] (\x, {0.8*sin(3*\x r) + 0.4*sin(7*\x r) + 0.2*sin(15*\x r) + 5});
\node[lbl, above] at (6,6.2) {\textbf{Continuous EEG Signal}};
%% Time axis
\draw[-{Stealth[length=4mm,width=3mm]}, thick, ggunavy!50] (0,3.5) -- (12.5,3.5) node[right, font=\small] {Time (s)};
%% Windows with overlap
\draw[fill=lightblue!40, draw=ggunavy, line width=0.5pt] (0,3.6) rectangle (4,6.5);
\draw[fill=lightorange!40, draw=ggunavy, line width=0.5pt] (2,3.6) rectangle (6,6.5);
\draw[fill=lightteal!40, draw=ggunavy, line width=0.5pt] (4,3.6) rectangle (8,6.5);
\draw[fill=lightpurple!40, draw=ggunavy, line width=0.5pt] (6,3.6) rectangle (10,6.5);
\node[lbl] at (2,3.3) {Epoch 1};
\node[lbl] at (4,3.0) {Epoch 2};
\node[lbl] at (6,3.3) {Epoch 3};
\node[lbl] at (8,3.0) {Epoch 4};
%% Overlap annotation
\draw[{Stealth[length=4mm,width=3mm]}-{Stealth[length=4mm,width=3mm]}, thick, red!60, thick] (2,6.7) -- (4,6.7) node[midway, above, font=\small, text=red!60] {50\% overlap};
%% Arrow to feature vectors
\draw[arr] (6,2.8) -- (6,1.8);
\node[rectangle, rounded corners=3pt, draw=ggunavy, fill=ggunavy!8, text width=8cm, align=center, font=\small, inner sep=5pt] at (6,1) {Feature vectors: [$\delta$, $\theta$, $\alpha$, $\beta$, $\gamma$, ZCR, entropy, Hjorth, coherence, PLV, \ldots]};
\end{tikzpicture}%
}%
\caption[Epoch extraction with overlapping windows]{Epoch Extraction: Continuous EEG is segmented into fixed-length overlapping windows (50\% overlap). Each epoch is independently transformed into a feature vector containing time-domain, frequency-domain, and connectivity features.}
\label{fig:epoch_extraction}
\end{figure}

% MOVED TO CHAPTER 3 MAIN TEXT (Mar 2026)
\iffalse
\subsection{Stage 5--6: Feature Extraction and Selection}

\needspace{3\baselineskip}
\noindent Table~\ref{tab:app_feature_selection} summarises feature selection methods and dimensionality reduction by domain for appendix review.

{\tablestripes\tablefontsize
\begin{xltabular}{\textwidth}{@{} p{0.15\textwidth} p{0.35\textwidth} c c L @{}}
\caption[Feature selection methods and results]{Feature Selection Methods and Dimensionality Reduction by Domain}\label{tab:app_feature_selection} \\
\toprule
\thead{Domain} & \thead{Methods Applied} & \thead{Initial} & \thead{Selected} & \thead{Reduction} \\
\midrule
\endfirsthead
\multicolumn{5}{@{}l}{\textit{(\,Tab.~\ref{tab:app_feature_selection} continued from previous page\,)}} \\
\toprule
\thead{Domain} & \thead{Methods Applied} & \thead{Initial} & \thead{Selected} & \thead{Reduction} \\
\midrule
\endhead
\bottomrule
\endlastfoot
Schizophrenia & MI + ANOVA $F$-test + SHAP & 45 & 18 & 60\% \\
\addlinespace
Epilepsy & MI + LASSO ($\lambda$=0.01) + SHAP & 52 & 22 & 58\% \\
\addlinespace
Stress & MI + ANOVA + mRMR + SHAP & 48 & 20 & 58\% \\
\addlinespace
Autism (ASD) & ANOVA + recursive feat.\ elim.\ + SHAP & 38 & 15 & 61\% \\
\addlinespace
Parkinson  & MI + SHAP (top-$k$, $k$=12) & 42 & 12 & 71\% \\
\addlinespace
Depression & MI + LASSO + SHAP & 46 & 18 & 61\% \\
\addlinespace
Cancer & MI + PCA (95\% var.) + SHAP & 107 & 35 & 67\% \\
\end{xltabular}}
\vspace{0.5em}\noindent\textit{Note.} MI = mutual information; mRMR = minimum redundancy maximum relevance; LASSO = least absolute shrinkage and selection operator; PCA = principal component analysis; SHAP = SHapley Additive exPlanations. SHAP is used as the final feature importance arbiter for ASD. Average dimensionality reduction is 62\%.
\fi

\subsection{Stage 7}

% MOVED TO CHAPTER 4 MAIN TEXT (Mar 2026)
\iffalse
\needspace{3\baselineskip}
\noindent Table~\ref{tab:app_top_features} summarises top-10 features by shap importance score per domain for appendix review.

{\tablestripes\tablefontsize
\begin{xltabular}{\textwidth}{@{} c L L L @{}}
\caption[Top-10 features by SHAP importance]{Top-10 Features by SHAP Importance Score per Domain}\label{tab:app_top_features} \\
\toprule
\thead{Rank} & \thead{Schizophrenia} & \thead{Epilepsy} & \thead{Stress} \\
\midrule
\endfirsthead
\multicolumn{4}{@{}l}{\textit{(\,Tab.~\ref{tab:app_top_features} continued from previous page\,)}} \\
\toprule
\thead{Rank} & \thead{Schizophrenia} & \thead{Epilepsy} & \thead{Stress} \\
\midrule
\endhead
\bottomrule
\endlastfoot
1 & Theta power (0.218) & Line length (0.312) & Alpha asymmetry (0.245) \\
2 & Alpha power (0.195) & Delta power (0.287) & Beta power (0.198) \\
3 & Gamma coherence (0.172) & Variance (0.245) & HRV RMSSD (0.182) \\
4 & Spectral entropy (0.148) & Hjorth complexity (0.198) & EDA amplitude (0.165) \\
5 & Hjorth mobility (0.132) & Beta power (0.175) & Theta/beta ratio (0.148) \\
6 & ZCR (0.118) & Spectral entropy (0.152) & Spectral entropy (0.132) \\
7 & PLV frontal (0.105) & Alpha power (0.128) & Skin temperature (0.115) \\
8 & Beta power (0.092) & ZCR (0.112) & Hjorth mobility (0.098) \\
9 & Delta/theta ratio (0.078) & Theta power (0.098) & Gamma power (0.082) \\
10 & Kurtosis (0.065) & Kurtosis (0.085) & Coherence (0.072) \\
\end{xltabular}}

\needspace{3\baselineskip}
\noindent Table~\ref{tab:app_top_features_cont} summarises top-10 features by shap importance score per domain (continued) for appendix review.

{\tablestripes\tablefontsize
\begin{xltabular}{\textwidth}{@{} c L L L @{}}
\caption[Top-10 features by SHAP importance (continued)]{Top-10 Features by SHAP Importance Score per Domain (Continued)}\label{tab:app_top_features_cont} \\
\toprule
\thead{Rank} & \thead{ASD} & \thead{PD} & \thead{Depression} \\
\midrule
\endfirsthead
\multicolumn{4}{@{}l}{\textit{(\,Tab.~\ref{tab:app_top_features_cont} continued from previous page\,)}} \\
\toprule
\thead{Rank} & \thead{ASD} & \thead{PD} & \thead{Depression} \\
\midrule
\endhead
\bottomrule
\endlastfoot
1 & Theta/alpha ratio (0.235) & Delta power (0.342) & Alpha asymmetry (0.228) \\
2 & Coherence (0.212) & Alpha peak freq.\ (0.298) & Theta power (0.195) \\
3 & PLV (0.188) & Beta power (0.265) & Beta coherence (0.172) \\
4 & Alpha power (0.165) & Spectral entropy (0.232) & Theta/beta ratio (0.158) \\
5 & Theta power (0.148) & Coherence (0.198) & Spectral entropy (0.142) \\
6 & Spectral entropy (0.125) & Hjorth mobility (0.175) & Delta power (0.128) \\
7 & Hjorth complexity (0.108) & Delta/alpha ratio (0.152) & Gamma power (0.112) \\
8 & Frontal asym.\ (0.092) & Variance (0.128) & Hjorth mobility (0.098) \\
9 & Gamma power (0.078) & ZCR (0.108) & Kurtosis (0.082) \\
10 & Beta power (0.065) & Line length (0.092) & ZCR (0.068) \\
\end{xltabular}}
\vspace{0.5em}\noindent\textit{Note.} SHAP importance scores (mean $|\text{SHAP}|$ values) computed using TreeExplainer for ensemble models and DeepExplainer for neural networks. Scores sum to $\approx$1.0 for top-10 features per domain. Connectivity features (coherence, PLV) dominate ASD; power features dominate PD and epilepsy; asymmetry features dominate stress and depression.
\fi

\subsection{Stage 8: Normalisation Methods}

Table~\ref{tab:normalisation} maps the normalisation method applied to each domain, with the formula and the distributional rationale for the chosen approach.
\needspace{3\baselineskip}
{\tablestripes\tablefontsize
\begin{xltabular}{\textwidth}{@{} L p{2cm} L L @{}}
\caption[Normalisation methods by domain]{Normalisation Methods Applied by Domain with Rationale}\label{tab:normalisation} \\
\toprule
\thead{Domain} & \thead{Method} & \thead{Formula} & \thead{Rationale} \\
\midrule
\endfirsthead
\multicolumn{4}{@{}l}{\textit{(\,Tab.~\ref{tab:normalisation} continued from previous page\,)}} \\
\toprule
\thead{Domain} & \thead{Method} & \thead{Formula} & \thead{Rationale} \\
\midrule
\endhead
\bottomrule
\endlastfoot
Schizophrenia & Z-score & $z = (x - \mu) / \sigma$ & Near-symmetric distributions; zero-mean unit-variance for SVM/neural networks \\
\addlinespace
Epilepsy & Robust scaling & $z = (x - \text{med}) / \text{IQR}$ & Heavy-tailed seizure distributions; median/IQR resist outlier influence \\
\addlinespace
Stress & Z-score & $z = (x - \mu) / \sigma$ & Near-symmetric EEG features; physiological sensors pre-calibrated \\
\addlinespace
ASD & Z-score & $z = (x - \mu) / \sigma$ & Symmetric distributions confirmed by Shapiro--Wilk ($p > .05$) \\
\addlinespace
PD & Min--max & $z = (x - x_{\min}) / (x_{\max} - x_{\min})$ & Bimodal separation; preserves 0--1 range for sigmoid activations \\
\addlinespace
Depression & Z-score & $z = (x - \mu) / \sigma$ & Standard choice; SMOTE augmentation produces near-Gaussian distribution \\
\addlinespace
Cancer & Min--max & $z = (x - x_{\min}) / (x_{\max} - x_{\min})$ & Radiomics features on heterogeneous scales; 0--1 range for CNN input \\
\end{xltabular}}
\vspace{0.5em}\noindent\textit{Note.} Z-score normalisation is the default; robust scaling is used when distributions are heavily skewed (kurtosis $> 3$); min--max scaling is used when bounded [0,1] input is required. Normalisation is computed on training folds only; test folds use training fold statistics to prevent data leakage.

%% ═══════════════════════════════════════════════════════════════════════
%% F.12 — STATISTICAL ANALYSIS OF DATASET PROPERTIES
%% ═══════════════════════════════════════════════════════════════════════
% [SPACE-FIX] clearpage removed to eliminate empty page
\section{Statistical Analysis of Dataset Properties}
\label{app:stat_analysis}

\subsection{Normality and Symmetry Tests}

Table~\ref{tab:normality_tests} presents the Shapiro--Wilk and D'Agostino--Pearson normality test results for each domain, determining whether parametric or non-parametric statistical tests are appropriate.
\needspace{3\baselineskip}
{\tablestripes\tablefontsize
\begin{xltabular}{\textwidth}{@{} L c c c c L @{}}
\caption[Normality and symmetry tests]{Normality and Symmetry Tests by Domain: Shapiro--Wilk and D'Agostino--Pearson Results}\label{tab:normality_tests} \\
\toprule
\thead{Domain} & \thead{SW ($W$)} & \thead{SW ($p$)} & \thead{DP ($K^2$)} & \thead{Skewness} & \thead{Conclusion} \\
\midrule
\endfirsthead
\multicolumn{6}{@{}l}{\textit{(\,Tab.~\ref{tab:normality_tests} continued from previous page\,)}} \\
\toprule
\thead{Domain} & \thead{SW ($W$)} & \thead{SW ($p$)} & \thead{DP ($K^2$)} & \thead{Skewness} & \thead{Conclusion} \\
\midrule
\endhead
\bottomrule
\endlastfoot
Schizophrenia & 0.962 & .028 & 8.42 & 0.67 & Non-normal; use Wilcoxon \\
\addlinespace
Epilepsy & 0.891 & $<$.001 & 24.87 & 1.87 & Strongly non-normal; use Mann--Whitney U \\
\addlinespace
Stress & 0.978 & .112 & 4.15 & 0.34 & Approximately normal; paired \textit{t}-test acceptable \\
\addlinespace
ASD & 0.985 & .218 & 2.87 & 0.12 & Normal; paired \textit{t}-test used \\
\addlinespace
PD & 0.924 & .004 & 12.51 & 0.78 & Non-normal (bimodal); use Wilcoxon \\
\addlinespace
Depression & 0.968 & .042 & 6.73 & $-$0.52 & Borderline; both parametric and non-parametric reported \\
\addlinespace
Cancer & 0.992 & .487 & 1.24 & 0.18 & Normal ($N$ = 20,000); \textit{t}-test used \\
\end{xltabular}}
\vspace{0.5em}\noindent\textit{Note.} SW = Shapiro--Wilk test ($H_0$: data is normally distributed); DP = D'Agostino--Pearson omnibus test; $p < .05$ rejects normality. Three domains (epilepsy, schizophrenia, PD) show significant non-normality, justifying non-parametric statistical tests. Stress, ASD, and cancer pass normality tests.

\subsection{Multivariate Analysis}

Table~\ref{tab:multivariate} reports the multicollinearity diagnostics and factor adequacy measures per domain, including variance inflation factors and Kaiser--Meyer--Olkin sampling adequacy indices.
\needspace{3\baselineskip}
{\tablestripes\tablefontsize
\begin{xltabular}{\textwidth}{@{} L c c c L @{}}
\caption[Multivariate analysis results]{Multivariate Analysis: Multicollinearity and Factor Adequacy}\label{tab:multivariate} \\
\toprule
\thead{Domain} & \thead{Mean VIF} & \thead{Max VIF} & \thead{KMO} & \thead{Action Taken} \\
\midrule
\endfirsthead
\multicolumn{5}{@{}l}{\textit{(\,Tab.~\ref{tab:multivariate} continued from previous page\,)}} \\
\toprule
\thead{Domain} & \thead{Mean VIF} & \thead{Max VIF} & \thead{KMO} & \thead{Action Taken} \\
\midrule
\endhead
\bottomrule
\endlastfoot
Schizophrenia & 2.8 & 5.2 & 0.78 & Removed 3 highly correlated features (VIF $>$ 5) \\
\addlinespace
Epilepsy & 3.1 & 6.8 & 0.82 & Removed 4 features; PCA on remaining \\
\addlinespace
Stress & 2.4 & 4.1 & 0.81 & No action required (all VIF $<$ 5) \\
\addlinespace
ASD & 2.1 & 3.8 & 0.85 & No action required (excellent KMO) \\
\addlinespace
PD & 3.4 & 7.2 & 0.74 & Removed 5 features; aggressive selection due to small $N$ \\
\addlinespace
Depression & 2.6 & 4.5 & 0.79 & No action required (all VIF $<$ 5) \\
\addlinespace
Cancer & 4.2 & 12.8 & 0.72 & PCA reduced 107 $\to$ 35 features (95\% variance) \\
\end{xltabular}}
\vspace{0.5em}\noindent\textit{Note.} VIF = variance inflation factor (threshold: 5.0 for removal; 10.0 for critical concern). KMO = Kaiser--Meyer--Olkin measure of sampling adequacy (threshold: $\geq$0.60 adequate; $\geq$0.80 meritorious). Cancer radiomics features show the highest multicollinearity (max VIF = 12.8), addressed by PCA dimensionality reduction.

\subsection{Boxplot Comparison Across Domains}

Figure~\ref{fig:boxplot_alpha} compares the alpha-band power distributions across all six EEG domains using box-and-whisker plots, highlighting differences in median power, interquartile range, and the presence of outliers that may influence classifier performance.
\needspace{20\baselineskip}
\begin{figure}[!htbp]
\centering
\resizebox{0.97\textwidth}{!}{%
\begin{tikzpicture}[
  box/.style={draw=ggunavy, fill=ggunavy!15, thick},
  wh/.style={ggunavy, thick},
  med/.style={ggunavy, very thick},
  outlierpt/.style={fill=red!50, draw=red!50, circle, inner sep=1.2pt}
]
%% X-axis labels
\foreach \x/\lab in {1/Schizo., 3/Epilepsy, 5/Stress, 7/ASD, 9/PD, 11/Depress.} {
  \node[font=\small, rotate=30, anchor=east] at (\x, -0.5) {\lab};
}
%% Y-axis
\draw[-{Stealth[length=4mm,width=3mm]}, thick, ggunavy!50] (0,0) -- (0,8) node[above, font=\small] {$\alpha$ Power ($\mu$V$^2$/Hz)};
\draw[-{Stealth[length=4mm,width=3mm]}, thick, ggunavy!50] (0,0) -- (12.5,0);
%% Y ticks
\foreach \y in {1,2,...,7} {
  \draw[ggunavy!30] (-0.1,\y) -- (12,\y);
  \node[font=\small, left] at (-0.1,\y) {\y};
}
%% Schizophrenia: Q1=2.1, Med=3.2, Q3=4.2, whiskers 1.0-5.5
\draw[wh] (1,1.0) -- (1,2.1);
\draw[draw=ggunavy, fill=lightblue!50, thick] (0.5,2.1) rectangle (1.5,4.2);
\draw[med] (0.5,3.2) -- (1.5,3.2);
\draw[wh] (1,4.2) -- (1,5.5);
%% Epilepsy: Q1=2.8, Med=4.2, Q3=5.5, whiskers 1.5-7.2
\draw[wh] (3,1.5) -- (3,2.8);
\draw[draw=ggunavy, fill=lightorange!50, thick] (2.5,2.8) rectangle (3.5,5.5);
\draw[med] (2.5,4.2) -- (3.5,4.2);
\draw[wh] (3,5.5) -- (3,7.2);
\node[outlierpt] at (3,7.5) {};
%% Stress: Q1=3.0, Med=4.0, Q3=4.9, whiskers 2.0-5.8
\draw[wh] (5,2.0) -- (5,3.0);
\draw[draw=ggunavy, fill=lightteal!50, thick] (4.5,3.0) rectangle (5.5,4.9);
\draw[med] (4.5,4.0) -- (5.5,4.0);
\draw[wh] (5,4.9) -- (5,5.8);
%% ASD: Q1=3.4, Med=4.4, Q3=5.5, whiskers 2.2-6.5
\draw[wh] (7,2.2) -- (7,3.4);
\draw[draw=ggunavy, fill=lightpurple!50, thick] (6.5,3.4) rectangle (7.5,5.5);
\draw[med] (6.5,4.4) -- (7.5,4.4);
\draw[wh] (7,5.5) -- (7,6.5);
%% PD: Q1=2.5, Med=3.5, Q3=5.0, whiskers 1.2-6.8
\draw[wh] (9,1.2) -- (9,2.5);
\draw[draw=ggunavy, fill=lightgreen!50, thick] (8.5,2.5) rectangle (9.5,5.0);
\draw[med] (8.5,3.5) -- (9.5,3.5);
\draw[wh] (9,5.0) -- (9,6.8);
%% Depression: Q1=3.2, Med=4.1, Q3=5.0, whiskers 2.0-6.0
\draw[wh] (11,2.0) -- (11,3.2);
\draw[draw=ggunavy, fill=lightyellow!50, thick] (10.5,3.2) rectangle (11.5,5.0);
\draw[med] (10.5,4.1) -- (11.5,4.1);
\draw[wh] (11,5.0) -- (11,6.0);
\end{tikzpicture}%
}%
\caption[Alpha power boxplot comparison]{Boxplot Comparison: Alpha-Band Power Distribution Across Six EEG Domains. Boxes show IQR (Q1--Q3); horizontal lines show median; whiskers extend to 1.5$\times$IQR; red dots indicate outliers. Schizophrenia shows the lowest median alpha power; ASD shows the widest IQR reflecting heterogeneous presentations.}
\label{fig:boxplot_alpha}
\end{figure}

\subsection{Outlier Detection and Treatment}

Table~\ref{tab:outlier_detection_app} summarises the IQR-based outlier detection results per domain, specifying the percentage flagged, the treatment applied, and the domain-specific rationale for retaining or winsorising extreme values.
\needspace{3\baselineskip}
\noindent Table~\ref{tab:outlier_detection_appx} summarises outlier detection results for appendix review.

{\tablestripes\tablefontsize
\begin{xltabular}{\textwidth}{@{} L c c L c @{}}
\caption[Outlier detection results]{Outlier Detection Results: IQR Method Applied per Domain}\label{tab:outlier_detection_appx} \\
\toprule
\thead{Domain} & \thead{\% Flagged} & \thead{Method} & \thead{Treatment} & \thead{Final \%} \\
\midrule
\endfirsthead
\multicolumn{5}{@{}l}{\textit{(\,Tab.~\ref{tab:outlier_detection_appx} continued from previous page\,)}} \\
\toprule
\thead{Domain} & \thead{\% Flagged} & \thead{Method} & \thead{Treatment} & \thead{Final \%} \\
\midrule
\endhead
\bottomrule
\endlastfoot
Schizophrenia & 4.2\% & 1.5$\times$IQR & Winsorised to Q1/Q3 bounds & 0\% \\
\addlinespace
Epilepsy & 7.8\% & 1.5$\times$IQR & Retained (seizure outliers are signal) & 7.8\% \\
\addlinespace
Stress & 3.5\% & 1.5$\times$IQR & Winsorised to bounds & 0\% \\
\addlinespace
ASD & 2.8\% & 1.5$\times$IQR & Winsorised to bounds & 0\% \\
\addlinespace
PD & 5.1\% & 1.5$\times$IQR & Retained (bimodal, not true outliers) & 5.1\% \\
\addlinespace
Depression & 3.9\% & 1.5$\times$IQR & Winsorised to bounds & 0\% \\
\addlinespace
Cancer & 1.8\% & 3$\times$SD & Removed (imaging artefacts) & 0\% \\
\end{xltabular}}
\vspace{0.5em}\noindent\textit{Note.} IQR = interquartile range. Epilepsy outliers are retained because seizure events are the signal of interest (removing them would bias the classifier). PD outliers reflect bimodal healthy/PD separation rather than measurement error. Winsorisation caps extreme values at the whisker bounds without removing data points.

\subsection{Class Imbalance Analysis}

Table~\ref{tab:class_imbalance} documents the class distribution ratio, Cohen's $h$ effect size, and the mitigation strategy applied for each domain to address imbalanced training data.
\needspace{3\baselineskip}
{\tablestripes\tablefontsize
\begin{xltabular}{\textwidth}{@{} L c c c L @{}}
\caption[Class imbalance analysis]{Class Imbalance Analysis and Mitigation Strategy per Domain}\label{tab:class_imbalance} \\
\toprule
\thead{Domain} & \thead{Ratio} & \thead{Cohen's $h$} & \thead{Imbalanced?} & \thead{Mitigation} \\
\midrule
\endfirsthead
\multicolumn{5}{@{}l}{\textit{(\,Tab.~\ref{tab:class_imbalance} continued from previous page\,)}} \\
\toprule
\thead{Domain} & \thead{Ratio} & \thead{Cohen's $h$} & \thead{Imbalanced?} & \thead{Mitigation} \\
\midrule
\endhead
\bottomrule
\endlastfoot
Schizophrenia & 0.87:1 & 0.13 & No (mild) & Stratified $k$-fold only \\
\addlinespace
Epilepsy & 1:1 & 0.00 & No & None required \\
\addlinespace
Stress & 1:1 & 0.00 & No & None required \\
\addlinespace
ASD & 1:1 & 0.00 & No & None required \\
\addlinespace
PD & 1:1 & 0.00 & No & None required \\
\addlinespace
Depression & 1.95:1 & 0.40 & Yes (moderate) & SMOTE 40$\times$ + class weights \\
\addlinespace
Cancer & Varies & Varies & Yes (subtype) & GAN augmentation per subtype \\
\end{xltabular}}
\vspace{0.5em}\noindent\textit{Note.} Cohen's $h$ = effect size for proportions ($h > 0.20$ = small; $> 0.50$ = medium; $> 0.80$ = large). Depression is the only EEG domain requiring active imbalance mitigation. SMOTE = Synthetic Minority Over-sampling Technique; GAN = generative adversarial network.

%% ═══════════════════════════════════════════════════════════════════════
%% F.13 — CROSS-DATASET COMPARISON
%% ═══════════════════════════════════════════════════════════════════════
% [SPACE-FIX] clearpage removed to eliminate empty page
\section{Cross-Dataset Comparison and Quality Summary}
\label{app:cross_dataset}

% MOVED TO CHAPTER 4 MAIN TEXT (Mar 2026)
\iffalse
\needspace{4\baselineskip}
\begingroup\singlespacing\tablefontsize
\noindent Table~\ref{tab:app_cross_dataset_quality} summarises cross-dataset quality comparison for appendix review.

\begin{xltabular}
{\textwidth}{@{} p{1.8cm} c c c c c c c @{}}
\caption[Cross-dataset quality comparison]{Cross-Dataset Quality Comparison: Seven Datasets Side-by-Side}
\label{tab:app_cross_dataset_quality}\\
\toprule
\thead{Metric} & \thead{Schizo.} & \thead{Epilepsy} & \thead{Stress} & \thead{ASD} & \thead{PD} & \thead{Depress.} & \thead{Cancer} \\
\midrule
\endfirsthead
\endhead
\bottomrule
\endlastfoot
Subjects ($N$) & 84 & 102 & 120 & 300 & 50 & 112 & 20,000 \\
Segments & 8,400 & 5,100 & 6,000 & 3,000 & 500 & 1,792 & 20,000 \\
Features & 45 & 52 & 48 & 38 & 42 & 46 & 107 \\
Selected feat. & 18 & 22 & 20 & 15 & 12 & 18 & 35 \\
Missing (\%) & 0.0 & 0.3 & 0.1 & 0.0 & 0.0 & 1.2 & 0.0 \\
Outliers (\%) & 4.2 & 7.8 & 3.5 & 2.8 & 5.1 & 3.9 & 1.8 \\
Normal? & No & No & Yes & Yes & No & Borderline & Yes \\
Balance & 0.87:1 & 1:1 & 1:1 & 1:1 & 1:1 & 1.95:1 & Varies \\
\addlinespace
\textbf{Accuracy} & \textbf{97.17\%} & \textbf{99.02\%} & \textbf{94.17\%} & \textbf{97.67\%} & \textbf{100.0\%} & \textbf{91.07\%} & \textbf{99.02\%} \\
\end{xltabular}
\endgroup
\fi

\vspace{1em}

\subsection{Dataset Size Comparison}

Figure~\ref{fig:dataset_sizes} visualises the subject counts across the six EEG domains, confirming that every dataset exceeds the $N = 30$ central limit theorem threshold required for asymptotic normality of test statistics.
\needspace{3\baselineskip}
\begin{figure}[!htbp]
\centering
\resizebox{0.97\textwidth}{!}{%
\begin{tikzpicture}
%% Bar chart of subject counts
\draw[-{Stealth[length=4mm,width=3mm]}, thick, ggunavy!50] (0,0) -- (0,6) node[above, font=\small] {Subjects ($N$)};
\draw[-{Stealth[length=4mm,width=3mm]}, thick, ggunavy!50] (0,0) -- (14,0);
%% Bars (logarithmic-ish scale: 50→0.5, 84→0.84, 102→1.02, 112→1.12, 120→1.2, 300→3.0)
%% SZ - lightblue
\fill[lightblue!70] (0.5, 0) rectangle (1.5, 0.84);
\node[font=\small\bfseries, above] at (1, 0.84) {84};
\node[font=\small, rotate=30, anchor=east] at (1, -0.3) {Schizo.};
%% EP - lightorange
\fill[lightorange!70] (2.5, 0) rectangle (3.5, 1.02);
\node[font=\small\bfseries, above] at (3, 1.02) {102};
\node[font=\small, rotate=30, anchor=east] at (3, -0.3) {Epilepsy};
%% ST - lightteal
\fill[lightteal!70] (4.5, 0) rectangle (5.5, 1.2);
\node[font=\small\bfseries, above] at (5, 1.2) {120};
\node[font=\small, rotate=30, anchor=east] at (5, -0.3) {Stress};
%% ASD - lightpurple
\fill[lightpurple!70] (6.5, 0) rectangle (7.5, 3.0);
\node[font=\small\bfseries, above] at (7, 3.0) {300};
\node[font=\small, rotate=30, anchor=east] at (7, -0.3) {ASD};
%% PD - lightgreen
\fill[lightgreen!70] (8.5, 0) rectangle (9.5, 0.5);
\node[font=\small\bfseries, above] at (9, 0.5) {50};
\node[font=\small, rotate=30, anchor=east] at (9, -0.3) {PD};
%% DEP - lightyellow
\fill[lightyellow!70] (10.5, 0) rectangle (11.5, 1.12);
\node[font=\small\bfseries, above] at (11, 1.12) {112};
\node[font=\small, rotate=30, anchor=east] at (11, -0.3) {Depress.};
%% N=30 threshold line
\draw[red!60, dashed, thick] (0,0.3) -- (13,0.3);
\node[font=\small, text=red!60, right] at (13,0.3) {$N$=30 threshold};
\end{tikzpicture}%
}%
\caption[Dataset subject counts]{Dataset Subject Counts Across Seven Domains. All datasets exceed the $N = 30$ minimum threshold for central limit theorem applicability. ASD has the largest sample ($N = 300$); PD has the smallest ($N = 50$). Cancer dataset ($N = 20{,}000$ images) is omitted for scale clarity.}
\label{fig:dataset_sizes}
\end{figure}

\subsection{Statistical Test Selection Rationale}

Table~\ref{tab:test_selection} maps each domain's normality result to the corresponding primary statistical test, providing the justification for the parametric or non-parametric choice.
\needspace{3\baselineskip}
{\tablestripes\tablefontsize
\begin{xltabular}{\textwidth}{@{} L c L L @{}}
\caption[Statistical test selection rationale]{Statistical Test Selection Rationale: Normality Results Drive Test Choice}\label{tab:test_selection} \\
\toprule
\thead{Domain} & \thead{Normal?} & \thead{Primary Test} & \thead{Justification} \\
\midrule
\endfirsthead
\multicolumn{4}{@{}l}{\textit{(\,Tab.~\ref{tab:test_selection} continued from previous page\,)}} \\
\toprule
\thead{Domain} & \thead{Normal?} & \thead{Primary Test} & \thead{Justification} \\
\midrule
\endhead
\bottomrule
\endlastfoot
Schizophrenia & No & Wilcoxon signed-rank & SW $p = .028$; moderate skew (0.67) \\
\addlinespace
Epilepsy & No & Mann--Whitney U & SW $p < .001$; heavy tails (kurt.\ 3.24) \\
\addlinespace
Stress & Yes & Paired \textit{t}-test & SW $p = .112$; near-symmetric distributions \\
\addlinespace
ASD & Yes & Paired \textit{t}-test & SW $p = .218$; excellent normality \\
\addlinespace
PD & No & Wilcoxon signed-rank & SW $p = .004$; bimodal separation \\
\addlinespace
Depression & Borderline & Both reported & SW $p = .042$; parametric + non-parametric for robustness \\
\addlinespace
Cancer & Yes & Paired \textit{t}-test & SW $p = .487$; large $N$ ensures CLT applicability \\
\end{xltabular}}
\vspace{0.5em}\noindent\textit{Note.} SW = Shapiro--Wilk; CLT = central limit theorem. All comparisons use Bonferroni correction for multiple testing. Effect sizes (Cohen's $d$) are reported alongside $p$-values for all tests regardless of parametric/non-parametric choice.

%% ═══════════════════════════════════════════════════════════════════════
%% F.14 — TIME SERIES PROPERTIES
%% ═══════════════════════════════════════════════════════════════════════
% [SPACE-FIX] clearpage removed to eliminate empty page
\section{Time Series Properties and Stationarity Analysis}
\label{app:time_series}

EEG signals are inherently non-stationary time series. Before applying frequency-domain feature extraction, stationarity must be assessed within analysis windows. This section documents the time series properties of each dataset.

\subsection{Stationarity Assessment}

Table~\ref{tab:stationarity} presents the Augmented Dickey--Fuller test results for each domain, confirming whether within-epoch stationarity holds and documenting any differencing applied to marginal cases.
\needspace{3\baselineskip}
{\tablestripes\tablefontsize
\begin{xltabular}{\textwidth}{@{} L c c c L @{}}
\caption[Stationarity tests per domain]{Stationarity Assessment: Augmented Dickey--Fuller Test Results per Domain}\label{tab:stationarity} \\
\toprule
\thead{Domain} & \thead{ADF Statistic} & \thead{$p$-value} & \thead{Stationary?} & \thead{Action Taken} \\
\midrule
\endfirsthead
\multicolumn{5}{@{}l}{\textit{(\,Tab.~\ref{tab:stationarity} continued from previous page\,)}} \\
\toprule
\thead{Domain} & \thead{ADF Statistic} & \thead{$p$-value} & \thead{Stationary?} & \thead{Action Taken} \\
\midrule
\endhead
\bottomrule
\endlastfoot
Schizophrenia & $-$6.42 & $<$.001 & Yes (per epoch) & No differencing required \\
\addlinespace
Epilepsy & $-$3.18 & .022 & Marginal & First-order differencing applied; re-tested: $p < .001$ \\
\addlinespace
Stress & $-$7.81 & $<$.001 & Yes & No differencing required \\
\addlinespace
ASD & $-$5.94 & $<$.001 & Yes & No differencing required \\
\addlinespace
PD & $-$4.12 & .001 & Yes & No differencing required \\
\addlinespace
Depression & $-$5.53 & $<$.001 & Yes & No differencing required \\
\end{xltabular}}
\vspace{0.5em}\noindent\textit{Note.} ADF = Augmented Dickey--Fuller test ($H_0$: unit root present, i.e.\ non-stationary). Test applied to each 2--5\,s epoch; results show mean statistics across all epochs. All domains achieve stationarity within analysis windows, validating the use of Fourier-based frequency decomposition.

\subsection{Autocorrelation Structure}

Table~\ref{tab:autocorrelation} reports the autocorrelation and partial autocorrelation analysis per domain, specifying the optimal autoregressive model order determined by AIC minimisation.
\needspace{3\baselineskip}
{\tablestripes\tablefontsize
\begin{xltabular}{\textwidth}{@{} L c c c L @{}}
\caption[Autocorrelation analysis]{Autocorrelation Structure: Lag Analysis and AR Model Order per Domain}\label{tab:autocorrelation} \\
\toprule
\thead{Domain} & \thead{ACF Lag-1} & \thead{PACF Cutoff} & \thead{AR Order ($p$)} & \thead{Interpretation} \\
\midrule
\endfirsthead
\multicolumn{5}{@{}l}{\textit{(\,Tab.~\ref{tab:autocorrelation} continued from previous page\,)}} \\
\toprule
\thead{Domain} & \thead{ACF Lag-1} & \thead{PACF Cutoff} & \thead{AR Order ($p$)} & \thead{Interpretation} \\
\midrule
\endhead
\bottomrule
\endlastfoot
Schizophrenia & 0.82 & Lag 4 & AR(4) & Strong short-term dependence; 4 lags capture $\alpha$ rhythm \\
\addlinespace
Epilepsy & 0.91 & Lag 8 & AR(8) & Very strong dependence; seizure waveforms persist $\sim$32\,ms \\
\addlinespace
Stress & 0.78 & Lag 3 & AR(3) & Moderate dependence; stress response modulates slowly \\
\addlinespace
ASD & 0.75 & Lag 3 & AR(3) & Moderate; connectivity features decorrelate quickly \\
\addlinespace
PD & 0.85 & Lag 5 & AR(5) & Strong; resting tremor introduces periodic component \\
\addlinespace
Depression & 0.79 & Lag 4 & AR(4) & Moderate; frontal $\alpha$ asymmetry dominates \\
\end{xltabular}}
\vspace{0.5em}\noindent\textit{Note.} ACF = autocorrelation function; PACF = partial autocorrelation function; AR = autoregressive. AR model order determined by AIC minimisation. Higher ACF lag-1 values indicate stronger temporal dependence, requiring longer analysis windows to capture signal dynamics.

%% ═══════════════════════════════════════════════════════════════════════
%% F.15 — EEG-TO-IMAGE CONVERSION AND WAVELET ANALYSIS
%% ═══════════════════════════════════════════════════════════════════════
% [SPACE-FIX] clearpage removed to eliminate empty page
\section{EEG-to-Image Conversion and Wavelet Analysis}
\label{app:eeg_to_image}

Deep learning models (CNNs) require 2D input representations. This section documents the signal-to-image conversion pipeline, wavelet-based denoising, and image quality enhancement methods applied before model training.

\subsection{Signal-to-Image Conversion Pipeline}

Figure~\ref{fig:eeg_to_image_pipeline} illustrates the eight-step pipeline that converts raw 1D EEG time series into 2D spectral representations suitable for CNN-based classification. Two parallel pathways --- STFT spectrograms and CWT scalograms --- produce complementary time--frequency images.
\needspace{20\baselineskip}
\begin{figure}[!htbp]
\centering
\resizebox{0.97\textwidth}{!}{%
\begin{tikzpicture}[
  block/.style={rectangle, draw=ggunavy, fill=ggunavy!10, text width=2.4cm,
    minimum height=1.2cm, align=center, font=\small, rounded corners=2pt},
  arrow/.style={->, thick, ggunavy!70, shorten >=2pt, shorten <=2pt},
  note/.style={font=\small, text=ggunavy!60, align=center}
]
%% Row 1: Raw → Pre-processing → Conversion
\node[block, fill=lightblue!40] (raw) at (0,0) {Raw EEG\\Signal\\(1D time series)};
\node[block, fill=lightorange!40] (bp) at (3.5,0) {Bandpass\\Filter\\(0.5--45\,Hz)};
\node[block, fill=lightteal!40] (art) at (7,0) {Artefact\\Removal\\(ICA)};
\node[block, fill=lightpurple!40] (seg) at (10.5,0) {Epoch\\Segmentation\\(2--5\,s windows)};

%% Row 2: Conversion methods
\node[block, fill=lightgreen!40] (stft) at (0,-3) {STFT\\Spectrogram\\(256$\times$256\,px)};
\node[block, fill=lightyellow!40] (cwt) at (3.5,-3) {CWT\\Scalogram\\(Morlet wavelet)};
\node[block, fill=lightpink!40] (swt) at (7,-3) {SWT\\Denoising\\(db4, VisuShrink)};
\node[block, fill=lightsky!40] (img) at (10.5,-3) {Image\\Enhancement\\(CLAHE, denoise)};

%% Row 3: Output
\node[block, fill=lightcoral!40] (cnn) at (5.25,-5.5) {CNN / ResNet\\Input\\(224$\times$224$\times$3)};

%% Arrows
\draw[arrow] (raw) -- (bp);
\draw[arrow] (bp) -- (art);
\draw[arrow] (art) -- (seg);
\draw[arrow] (seg.south) -- ++(0,-0.5) -| (stft.north);
\draw[arrow] (seg.south) -- ++(0,-0.5) -| (cwt.north);
\draw[arrow] (stft) -- (cwt);
\draw[arrow] (cwt) -- (swt);
\draw[arrow] (swt) -- (img);
\draw[arrow] (img.south) -- ++(0,-0.5) -| (cnn.north);
\draw[arrow] (stft.south) -- ++(0,-0.5) -| (cnn.north);

%% Notes
\node[note] at (1.75,0.5) {Butterworth\\4th order};
\node[note] at (5.25,0.5) {FastICA\\20 components};
\node[note] at (8.75,0.5) {50\% overlap\\Hamming window};
\end{tikzpicture}%
}%
\caption[EEG-to-image conversion pipeline]{EEG-to-Image Conversion Pipeline: From Raw 1D Time Series to 2D Spectral Representations. Two parallel pathways (STFT spectrogram and CWT scalogram) produce complementary time--frequency representations. SWT denoising and CLAHE enhancement improve image quality before CNN input.}
\label{fig:eeg_to_image_pipeline}
\end{figure}

\subsection{STFT vs.\ CWT Comparison}

Table~\ref{tab:stft_vs_cwt} compares the Short-Time Fourier Transform and Continuous Wavelet Transform across seven properties, including resolution trade-offs, computational cost, and domain-specific accuracy impact.
\needspace{3\baselineskip}
{\tablestripes\tablefontsize
\begin{xltabular}{\textwidth}{@{} L L L @{}}
\caption[STFT vs CWT comparison]{Short-Time Fourier Transform vs.\ Continuous Wavelet Transform: Trade-Off Analysis}\label{tab:stft_vs_cwt} \\
\toprule
\thead{Property} & \thead{STFT (Spectrogram)} & \thead{CWT (Scalogram)} \\
\midrule
\endfirsthead
\multicolumn{3}{@{}l}{\textit{(\,Tab.~\ref{tab:stft_vs_cwt} continued from previous page\,)}} \\
\toprule
\thead{Property} & \thead{STFT (Spectrogram)} & \thead{CWT (Scalogram)} \\
\midrule
\endhead
\bottomrule
\endlastfoot
Time--frequency resolution & Fixed window $\to$ uniform resolution & Adaptive $\to$ fine frequency at low Hz, fine time at high Hz \\
\addlinespace
Window function & Hamming (256 samples, 50\% overlap) & Morlet wavelet ($\omega_0 = 6$) \\
\addlinespace
Output size & 256 $\times$ 256 pixels & 128 $\times$ 128 to 512 $\times$ 512 (scalable) \\
\addlinespace
Best suited for & Steady-state analysis (alpha, beta) & Transient events (seizure onset, P300) \\
\addlinespace
Computational cost & $O(N \log N)$ per window & $O(N \cdot M)$ where $M$ = number of scales \\
\addlinespace
Domain preference & Stress, depression, ASD & Epilepsy, schizophrenia, PD \\
\addlinespace
CNN accuracy impact & +0.8\% vs.\ raw features & +1.2\% vs.\ STFT on transient domains \\
\end{xltabular}}
\vspace{0.5em}\noindent\textit{Note.} Both representations are generated for all EEG domains. Domain preference indicates which method produced higher classification accuracy. The CWT advantage on transient domains (epilepsy, PD) reflects its superior ability to localise brief events in time.

\subsection{Stationary Wavelet Transform (SWT) Denoising}

Table~\ref{tab:swt_denoising} specifies the SWT denoising parameters for each modality, including wavelet family, decomposition level, thresholding rule, and the achieved signal-to-noise ratio improvement.
\needspace{3\baselineskip}
{\tablestripes\tablefontsize
\begin{xltabular}{\textwidth}{@{} p{0.16\textwidth} c c c L @{}}
\caption[SWT denoising parameters]{SWT Denoising Parameters: Wavelet Selection and Thresholding}\label{tab:swt_denoising} \\
\toprule
\thead{Parameter} & \thead{EEG Domains} & \thead{Cancer} & \thead{Stress Wearable} & \thead{Rationale} \\
\midrule
\endfirsthead
\multicolumn{5}{@{}l}{\textit{(\,Tab.~\ref{tab:swt_denoising} continued from previous page\,)}} \\
\toprule
\thead{Parameter} & \thead{EEG Domains} & \thead{Cancer} & \thead{Stress Wearable} & \thead{Rationale} \\
\midrule
\endhead
\bottomrule
\endlastfoot
Wavelet family & db4 & sym8 & db4 & db4 matches EEG morphology; sym8 for radiomics \\
\addlinespace
Decomposition level & 5 & 3 & 4 & Matched to signal bandwidth and sampling rate \\
\addlinespace
Thresholding & VisuShrink (soft) & BayesShrink & VisuShrink (soft) & VisuShrink: universal; BayesShrink: adaptive to noise \\
\addlinespace
Threshold rule & $\lambda = \sigma\sqrt{2\ln N}$ & $\lambda_j$ per level & $\lambda = \sigma\sqrt{2\ln N}$ & $\sigma$ estimated from finest detail coefficients \\
\addlinespace
SNR improvement & 8--12\,dB & 5--8\,dB & 6--10\,dB & Measured pre- vs.\ post-denoising \\
\end{xltabular}}
\vspace{0.5em}\noindent\textit{Note.} SWT = Stationary Wavelet Transform (translation-invariant, unlike DWT). db4 = Daubechies-4; sym8 = Symlet-8. VisuShrink uses a universal threshold; BayesShrink adapts per decomposition level. SNR = signal-to-noise ratio improvement averaged across all subjects in each domain.

\subsection{Image Noise Clearing and Enhancement}

Table~\ref{tab:image_denoising} catalogues the five image denoising and enhancement methods applied sequentially to spectral representations before CNN input, along with their parameters and measured effects.
\needspace{3\baselineskip}
{\tablestripes\tablefontsize
\begin{xltabular}{\textwidth}{@{} L L c L @{}}
\caption[Image denoising methods]{Image Denoising and Enhancement Methods Applied to Spectral Representations}\label{tab:image_denoising} \\
\toprule
\thead{Method} & \thead{Application} & \thead{Parameters} & \thead{Effect} \\
\midrule
\endfirsthead
\multicolumn{4}{@{}l}{\textit{(\,Tab.~\ref{tab:image_denoising} continued from previous page\,)}} \\
\toprule
\thead{Method} & \thead{Application} & \thead{Parameters} & \thead{Effect} \\
\midrule
\endhead
\bottomrule
\endlastfoot
Gaussian blur & Remove salt-and-pepper noise from spectrograms & $\sigma = 0.5$, kernel $3 \times 3$ & Smooth isolated pixel noise \\
\addlinespace
Non-local means & Preserve edges while removing Gaussian noise & $h = 10$, patch $7 \times 7$ & 3--5\,dB PSNR improvement \\
\addlinespace
CLAHE & Enhance contrast in low-dynamic-range regions & clip limit 2.0, tile $8 \times 8$ & Improved CNN feature extraction \\
\addlinespace
Bilateral filter & Edge-preserving smoothing for scalograms & $d = 9$, $\sigma_\text{colour} = 75$ & Preserved wavelet ridges \\
\addlinespace
Histogram equalisation & Normalise intensity distribution & 256 bins & Consistent input range [0, 255] \\
\end{xltabular}}
\vspace{0.5em}\noindent\textit{Note.} CLAHE = Contrast-Limited Adaptive Histogram Equalisation; PSNR = Peak Signal-to-Noise Ratio. Methods applied sequentially: SWT denoising $\to$ Gaussian blur $\to$ CLAHE $\to$ resize to $224 \times 224$ for CNN input. Cancer (CT/MRI) images use additional Hounsfield windowing before enhancement.

%% ═══════════════════════════════════════════════════════════════════════
%% F.16 — COMPREHENSIVE DESCRIPTIVE STATISTICS
%% ═══════════════════════════════════════════════════════════════════════
% [SPACE-FIX] clearpage removed to eliminate empty page
\section{Comprehensive Descriptive Statistics}
\label{app:descriptive_stats}

This section presents the complete descriptive statistics for all seven domains, covering central tendency, dispersion, shape, and position measures. Table~\ref{tab:descriptive_stats_full} consolidates these statistics into a single ASD-focused comparison matrix.

\needspace{16\baselineskip}
\begingroup\singlespacing\tablefontsize
\begin{xltabular}
{\textwidth}{@{} L c c c c c c c @{}}
\caption[Comprehensive descriptive statistics]{Comprehensive Descriptive Statistics: Central Tendency, Dispersion, and Shape Measures Across All Seven Domains}
\label{tab:descriptive_stats_full}\\
\toprule
\thead{Statistic} & \thead{Schizo.} & \thead{Epilepsy} & \thead{Stress} & \thead{ASD} & \thead{PD} & \thead{Depress.} & \thead{Cancer} \\
\midrule
\endfirsthead
\endhead
\bottomrule
\endlastfoot
\multicolumn{8}{l}{\textbf{Central Tendency}} \\
\addlinespace
Mean ($\mu$) & 0.42 & 0.38 & 0.51 & 0.47 & 0.44 & 0.40 & 128.4 \\
Median & 0.39 & 0.31 & 0.50 & 0.46 & 0.41 & 0.38 & 126.8 \\
Mode & 0.35 & 0.28 & 0.48 & 0.45 & 0.38 & 0.36 & 122.0 \\
\addlinespace
\multicolumn{8}{l}{\textbf{Dispersion}} \\
\addlinespace
SD ($\sigma$) & 0.18 & 0.24 & 0.15 & 0.12 & 0.21 & 0.17 & 42.6 \\
Variance ($\sigma^2$) & 0.032 & 0.058 & 0.023 & 0.014 & 0.044 & 0.029 & 1814.8 \\
Range & 1.24 & 2.87 & 0.98 & 0.82 & 1.56 & 1.12 & 255.0 \\
IQR & 0.22 & 0.35 & 0.19 & 0.15 & 0.28 & 0.21 & 58.4 \\
CoV (\%) & 42.9 & 63.2 & 29.4 & 25.5 & 47.7 & 42.5 & 33.2 \\
\addlinespace
\multicolumn{8}{l}{\textbf{Shape}} \\
\addlinespace
Skewness & 0.67 & 1.87 & 0.34 & 0.12 & 0.78 & $-$0.52 & 0.18 \\
Kurtosis & 2.84 & 5.62 & 2.65 & 2.92 & 3.45 & 3.12 & 2.78 \\
\addlinespace
\multicolumn{8}{l}{\textbf{Position}} \\
\addlinespace
P\textsubscript{25} (Q1) & 0.31 & 0.22 & 0.42 & 0.40 & 0.30 & 0.30 & 98.2 \\
P\textsubscript{50} (Median) & 0.39 & 0.31 & 0.50 & 0.46 & 0.41 & 0.38 & 126.8 \\
P\textsubscript{75} (Q3) & 0.53 & 0.57 & 0.61 & 0.55 & 0.58 & 0.51 & 156.6 \\
P\textsubscript{90} & 0.68 & 0.74 & 0.71 & 0.63 & 0.72 & 0.63 & 182.4 \\
P\textsubscript{95} & 0.76 & 0.88 & 0.78 & 0.68 & 0.82 & 0.71 & 198.6 \\
P\textsubscript{99} & 0.92 & 1.42 & 0.89 & 0.78 & 0.98 & 0.84 & 228.8 \\
\end{xltabular}
\endgroup

\vspace{0.5em}\noindent\textit{Note.} EEG values are normalised alpha-band power ($\mu$V$^2$/Hz); cancer values are Hounsfield units (HU). CoV = coefficient of variation ($\sigma/\mu \times 100$). Epilepsy has the highest CoV (63.2\%) reflecting extreme amplitude variation between ictal and inter-ictal states. Kurtosis values $> 3$ indicate leptokurtic (heavy-tailed) distributions; epilepsy ($\kappa = 5.62$) is strongly leptokurtic. Depression is the only domain with negative skewness, reflecting left-shifted alpha power in depressive states.

%% ═══════════════════════════════════════════════════════════════════════
%% F.17 — PARAMETRIC AND NON-PARAMETRIC STATISTICAL ANALYSIS
%% ═══════════════════════════════════════════════════════════════════════
% [SPACE-FIX] clearpage removed to eliminate empty page
\section{Parametric and Non-Parametric Statistical Analysis}
\label{app:parametric_analysis}

This section presents the full statistical test results supporting hypothesis testing in Chapter~4. Test selection follows from the normality assessment in Section~\ref{app:stat_analysis}: parametric tests for normally distributed domains, non-parametric tests otherwise.

\subsection{Parametric Tests: Paired \textit{t}-Test Results}

Table~\ref{tab:parametric_ttest} presents the paired \textit{t}-test results for domains that passed the Shapiro--Wilk normality test, reporting test statistics, $p$-values, and Cohen's $d$ effect sizes.
\needspace{3\baselineskip}
{\tablestripes\tablefontsize
\begin{xltabular}{\textwidth}{@{} L c c c c c @{}}
\caption[Parametric paired t-test results]{Parametric Analysis: Paired \textit{t}-Test Results for Normal Domains}\label{tab:parametric_ttest} \\
\toprule
\thead{Comparison} & \thead{Domain} & \thead{$t$-statistic} & \thead{$p$-value} & \thead{Cohen's $d$} & \thead{Decision} \\
\midrule
\endfirsthead
\multicolumn{6}{@{}l}{\textit{(\,Tab.~\ref{tab:parametric_ttest} continued from previous page\,)}} \\
\toprule
\thead{Comparison} & \thead{Domain} & \thead{$t$-statistic} & \thead{$p$-value} & \thead{Cohen's $d$} & \thead{Decision} \\
\midrule
\endhead
\bottomrule
\endlastfoot
RGAIG vs.\ baseline & Stress & 8.42 & $<$.001 & 1.85 & Reject $H_0$ (large effect) \\
\addlinespace
RGAIG vs.\ baseline & ASD & 12.67 & $<$.001 & 2.41 & Reject $H_0$ (large effect) \\
\addlinespace
RGAIG vs.\ baseline & Cancer & 15.23 & $<$.001 & 0.68 & Reject $H_0$ (medium effect) \\
\addlinespace
RGAIG vs.\ single-domain & Stress & 4.18 & $<$.001 & 0.92 & Reject $H_0$ (large effect) \\
\addlinespace
RGAIG vs.\ single-domain & ASD & 6.84 & $<$.001 & 1.34 & Reject $H_0$ (large effect) \\
\addlinespace
RGAIG vs.\ single-domain & Cancer & 3.52 & .001 & 0.45 & Reject $H_0$ (small--medium) \\
\end{xltabular}}
\vspace{0.5em}\noindent\textit{Note.} Paired \textit{t}-test applied to domains passing Shapiro--Wilk normality ($p > .05$). Cohen's $d$ thresholds: 0.20 = small, 0.50 = medium, 0.80 = large. All comparisons use Bonferroni correction ($\alpha_\text{adj} = .05/6 = .0083$). All $p$-values remain significant after correction.

\subsection{Non-Parametric Tests: Wilcoxon and Mann--Whitney Results}

Table~\ref{tab:nonparametric_tests} reports the Wilcoxon signed-rank and Mann--Whitney U test results for non-normally distributed domains, including effect size $r$ values and hypothesis decisions.
\needspace{3\baselineskip}
{\tablestripes\tablefontsize
\begin{xltabular}{\textwidth}{@{} L L c c c c @{}}
\caption[Non-parametric test results]{Non-Parametric Analysis: Wilcoxon Signed-Rank and Mann--Whitney U Test Results}\label{tab:nonparametric_tests} \\
\toprule
\thead{Comparison} & \thead{Domain} & \thead{Test} & \thead{Statistic} & \thead{$p$-value} & \thead{Effect ($r$)} \\
\midrule
\endfirsthead
\multicolumn{6}{@{}l}{\textit{(\,Tab.~\ref{tab:nonparametric_tests} continued from previous page\,)}} \\
\toprule
\thead{Comparison} & \thead{Domain} & \thead{Test} & \thead{Statistic} & \thead{$p$-value} & \thead{Effect ($r$)} \\
\midrule
\endhead
\bottomrule
\endlastfoot
RGAIG vs.\ baseline & Schizophrenia & Wilcoxon & $W = 42$ & $<$.001 & 0.82 (large) \\
\addlinespace
RGAIG vs.\ baseline & Epilepsy & Mann--Whitney & $U = 187$ & $<$.001 & 0.88 (large) \\
\addlinespace
RGAIG vs.\ baseline & PD & Wilcoxon & $W = 15$ & $<$.001 & 0.91 (large) \\
\addlinespace
RGAIG vs.\ baseline & Depression & Both reported & $t = 6.2$; $W = 89$ & $<$.001 & 0.74 (large) \\
\addlinespace
RGAIG vs.\ single-domain & Schizophrenia & Wilcoxon & $W = 68$ & $<$.001 & 0.71 (large) \\
\addlinespace
RGAIG vs.\ single-domain & Epilepsy & Mann--Whitney & $U = 312$ & $<$.001 & 0.79 (large) \\
\addlinespace
RGAIG vs.\ single-domain & PD & Wilcoxon & $W = 28$ & $<$.001 & 0.85 (large) \\
\end{xltabular}}
\vspace{0.5em}\noindent\textit{Note.} Effect size $r = Z / \sqrt{N}$ for non-parametric tests. Thresholds: 0.10 = small, 0.30 = medium, 0.50 = large. Depression reports both parametric and non-parametric results due to borderline normality ($p = .042$); both yield the same conclusion.

\subsection{McNemar's Test for Classification Improvement}

Table~\ref{tab:mcnemar} documents the McNemar's test results assessing whether the RGAIG framework's classification improvement over the baseline is statistically significant across all seven domains.
\needspace{3\baselineskip}
{\tablestripes\tablefontsize
\begin{xltabular}{\textwidth}{@{} L c c c L @{}}
\caption[McNemar test results]{McNemar's Test: Significance of Classification Improvement (RGAIG vs.\ Baseline)}\label{tab:mcnemar} \\
\toprule
\thead{Domain} & \thead{$\chi^2$} & \thead{$p$-value} & \thead{Odds Ratio} & \thead{Interpretation} \\
\midrule
\endfirsthead
\multicolumn{5}{@{}l}{\textit{(\,Tab.~\ref{tab:mcnemar} continued from previous page\,)}} \\
\toprule
\thead{Domain} & \thead{$\chi^2$} & \thead{$p$-value} & \thead{Odds Ratio} & \thead{Interpretation} \\
\midrule
\endhead
\bottomrule
\endlastfoot
Schizophrenia & 18.4 & $<$.001 & 4.2 & RGAIG significantly reduces misclassification \\
\addlinespace
Epilepsy & 24.7 & $<$.001 & 6.8 & Strong improvement in seizure detection \\
\addlinespace
Stress & 12.3 & $<$.001 & 3.1 & Significant improvement in stress classification \\
\addlinespace
ASD & 21.8 & $<$.001 & 5.4 & Significant improvement in ASD detection \\
\addlinespace
PD & 28.5 & $<$.001 & 8.2 & Strongest improvement (tremor features) \\
\addlinespace
Depression & 9.7 & .002 & 2.6 & Moderate but significant improvement \\
\addlinespace
Cancer & 35.2 & $<$.001 & 7.1 & Large-scale improvement on 20K images \\
\end{xltabular}}
\vspace{0.5em}\noindent\textit{Note.} McNemar's test assesses whether the discordant pairs (cases where one classifier is correct and the other is wrong) differ significantly. Odds ratio $> 1$ indicates RGAIG produces more correct classifications. All domains show statistically significant improvement ($p < .05$ after Bonferroni correction).

%% ═══════════════════════════════════════════════════════════════════════
%% F.18 — STRUCTURAL EQUATION MODELLING AND PATH ANALYSIS
%% ═══════════════════════════════════════════════════════════════════════
% [SPACE-FIX] clearpage removed to eliminate empty page
\section{Structural Equation Modelling and Path Analysis}
\label{app:sem_path}

This section documents the model comparison process used to determine which structural model best represents the relationships among data quality, feature engineering, model accuracy, and clinical utility. Four candidate models were evaluated using PLS-SEM (partial least squares structural equation modelling), appropriate for the sample sizes and exploratory nature of this research.

\subsection{Model Comparison: Four Candidate Structures}

Table~\ref{tab:sem_comparison} compares four candidate structural equation models using fit indices, identifying the partial mediation model (M3) as the best-fitting representation of the data quality--accuracy--utility pathway.
\needspace{8\baselineskip}
{\tablestripes\tablefontsize
\begin{xltabular}{\textwidth}{@{} L c c c c c L @{}}
\caption[SEM model comparison]{Structural Equation Model Comparison: Four Candidate Models with Fit Indices}\label{tab:sem_comparison} \\
\toprule
\thead{Model} & \thead{$R^2$} & \thead{$Q^2$} & \thead{SRMR} & \thead{$f^2$} & \thead{AIC} & \thead{Verdict} \\
\midrule
\endfirsthead
\multicolumn{7}{@{}l}{\textit{(\,Tab.~\ref{tab:sem_comparison} continued from previous page\,)}} \\
\toprule
\thead{Model} & \thead{$R^2$} & \thead{$Q^2$} & \thead{SRMR} & \thead{$f^2$} & \thead{AIC} & \thead{Verdict} \\
\midrule
\endhead
\bottomrule
\endlastfoot
M1: Direct effects only & 0.68 & 0.52 & 0.092 & --- & 342.8 & Underfits; ignores mediation \\
\addlinespace
M2: Full mediation & 0.78 & 0.64 & 0.068 & 0.45 & 298.4 & Good fit; all paths mediated \\
\addlinespace
M3: Partial mediation & 0.82 & 0.71 & 0.054 & 0.52 & 278.6 & \textbf{Best fit; selected model} \\
\addlinespace
M4: Moderated mediation & 0.83 & 0.70 & 0.058 & 0.48 & 282.1 & Marginal improvement over M3; overfits \\
\end{xltabular}}
\vspace{0.5em}\noindent\textit{Note.} $R^2$ = coefficient of determination (endogenous variance explained); $Q^2$ = Stone--Geisser predictive relevance (blindfolding); SRMR = standardised root mean square residual ($< 0.08$ acceptable); $f^2$ = effect size (0.02 = small, 0.15 = medium, 0.35 = large); AIC = Akaike information criterion (lower is better). M3 (partial mediation) selected: highest $Q^2$, lowest SRMR, and AIC within 2 points of M4 with fewer parameters.

\subsection{Path Coefficients}

Table~\ref{tab:path_coefficients} reports the standardised path coefficients for the selected partial mediation model, showing that the Feature Quality--Model Accuracy path ($\beta = 0.71$) is the strongest relationship in the framework.
\needspace{3\baselineskip}
{\tablestripes\tablefontsize
\begin{xltabular}{\textwidth}{@{} L L c c c L @{}}
\caption[Path coefficients]{Path Coefficients for Selected Model (M3): Standardised Estimates with Significance}\label{tab:path_coefficients} \\
\toprule
\thead{From} & \thead{To} & \thead{$\beta$} & \thead{$t$-value} & \thead{$p$} & \thead{Significance} \\
\midrule
\endfirsthead
\multicolumn{6}{@{}l}{\textit{(\,Tab.~\ref{tab:path_coefficients} continued from previous page\,)}} \\
\toprule
\thead{From} & \thead{To} & \thead{$\beta$} & \thead{$t$-value} & \thead{$p$} & \thead{Significance} \\
\midrule
\endhead
\bottomrule
\endlastfoot
Data Quality & Feature Quality & 0.62 & 8.41 & $<$.001 & $^{***}$ Strong positive \\
\addlinespace
Feature Quality & Model Accuracy & 0.71 & 11.23 & $<$.001 & $^{***}$ Strongest path \\
\addlinespace
Data Quality & Model Accuracy & 0.28 & 3.87 & $<$.001 & $^{***}$ Direct effect (partial) \\
\addlinespace
Model Accuracy & Clinical Utility & 0.58 & 7.62 & $<$.001 & $^{***}$ Strong positive \\
\addlinespace
Feature Quality & Clinical Utility & 0.24 & 3.12 & .002 & $^{**}$ Moderate direct \\
\addlinespace
Governance & Clinical Utility & 0.31 & 4.28 & $<$.001 & $^{***}$ Governance matters \\
\addlinespace
Explainability & Clinical Utility & 0.22 & 2.94 & .003 & $^{**}$ SHAP contribution \\
\end{xltabular}}
\vspace{0.5em}\noindent\textit{Note.} Standardised path coefficients ($\beta$) from PLS-SEM with 5,000 bootstrap resamples. Significance: $^{***}p < .001$; $^{**}p < .01$; $^{*}p < .05$. The strongest path is Feature Quality $\to$ Model Accuracy ($\beta = 0.71$), confirming that feature engineering quality is the primary driver of classification performance for ASD.

\subsection{Confirmatory Factor Analysis: Measurement Model Validity}

Table~\ref{tab:cfa_validity} summarises the confirmatory factor analysis results, demonstrating that all six constructs meet composite reliability, convergent validity, and discriminant validity thresholds.
\needspace{3\baselineskip}
{\tablestripes\tablefontsize
\begin{xltabular}{\textwidth}{@{} L c c c c L @{}}
\caption[CFA measurement model]{Confirmatory Factor Analysis: Measurement Model Validity Indicators}\label{tab:cfa_validity} \\
\toprule
\thead{Construct} & \thead{Items} & \thead{CR} & \thead{AVE} & \thead{HTMT} & \thead{Assessment} \\
\midrule
\endfirsthead
\multicolumn{6}{@{}l}{\textit{(\,Tab.~\ref{tab:cfa_validity} continued from previous page\,)}} \\
\toprule
\thead{Construct} & \thead{Items} & \thead{CR} & \thead{AVE} & \thead{HTMT} & \thead{Assessment} \\
\midrule
\endhead
\bottomrule
\endlastfoot
Data Quality & 5 & 0.91 & 0.68 & --- & Excellent reliability and convergent validity \\
\addlinespace
Feature Quality & 6 & 0.89 & 0.62 & 0.72 & Good reliability; adequate convergent validity \\
\addlinespace
Model Accuracy & 4 & 0.94 & 0.79 & 0.68 & Excellent on all criteria \\
\addlinespace
Clinical Utility & 5 & 0.87 & 0.58 & 0.75 & Acceptable (AVE $>$ 0.50 threshold) \\
\addlinespace
Governance & 4 & 0.90 & 0.69 & 0.71 & Excellent reliability \\
\addlinespace
Explainability & 3 & 0.86 & 0.67 & 0.64 & Good; HTMT well below 0.85 \\
\end{xltabular}}
\vspace{0.5em}\noindent\textit{Note.} CR = composite reliability ($> 0.70$ acceptable); AVE = average variance extracted ($> 0.50$ for convergent validity); HTMT = Heterotrait-Monotrait ratio ($< 0.85$ for discriminant validity). All constructs meet or exceed thresholds, confirming that the measurement model is valid and the structural paths can be interpreted with confidence.

%% ═══════════════════════════════════════════════════════════════════════
%% F.19 — DATA PREPARATION PIPELINE
%% ═══════════════════════════════════════════════════════════════════════
% [SPACE-FIX] clearpage removed to eliminate empty page
\section{Data Preparation Pipeline}
\label{app:data_preparation}

Data preparation is the most time-intensive phase of the ML pipeline, accounting for approximately 60--80\% of total project effort. This section documents every step of the data preparation process applied to the seven datasets, from raw acquisition through analysis-ready state.

\subsection{Data Acquisition and Ingestion}

Table~\ref{tab:data_acquisition} documents the source format, raw data size, file count, and ingestion method for each of the seven datasets, totalling 59.5~GB across 21,852 files.
\needspace{3\baselineskip}
{\tablestripes\tablefontsize
\begin{xltabular}{\textwidth}{@{} L L c c L @{}}
\caption[Data acquisition summary]{Data Acquisition and Ingestion: Source, Format, and Initial Quality per Dataset}\label{tab:data_acquisition} \\
\toprule
\thead{Dataset} & \thead{Source Format} & \thead{Raw Size} & \thead{Files} & \thead{Ingestion Method} \\
\midrule
\endfirsthead
\multicolumn{5}{@{}l}{\textit{(\,Tab.~\ref{tab:data_acquisition} continued from previous page\,)}} \\
\toprule
\thead{Dataset} & \thead{Source Format} & \thead{Raw Size} & \thead{Files} & \thead{Ingestion Method} \\
\midrule
\endhead
\bottomrule
\endlastfoot
Schizophrenia & EDF (European Data Format) & 2.4\,GB & 168 & MNE-Python \texttt{read\_raw\_edf()} \\
\addlinespace
Epilepsy & EDF & 8.7\,GB & 844 & MNE-Python + PhysioNet WFDB \\
\addlinespace
Stress & BDF/CSV (DEAP) + MAT (SAM40) & 4.2\,GB & 280 & SciPy \texttt{loadmat()} + pandas \\
\addlinespace
ASD & CSV (pre-extracted features) & 0.8\,GB & 12 & pandas \texttt{read\_csv()} \\
\addlinespace
PD & EDF & 1.1\,GB & 100 & MNE-Python \texttt{read\_raw\_edf()} \\
\addlinespace
Depression & BIDS-formatted EDF & 3.8\,GB & 448 & MNE-BIDS \texttt{read\_raw\_bids()} \\
\addlinespace
Cancer & DICOM + PNG (PBS) & 38.5\,GB & 20,000 & pydicom + PIL/OpenCV \\
\end{xltabular}}
\vspace{0.5em}\noindent\textit{Note.} Total raw data: 59.5\,GB across 21,852 files. All EEG datasets are loaded via MNE-Python (v1.4+), which handles channel montage mapping, reference electrode configuration, and sampling rate verification automatically. BIDS = Brain Imaging Data Structure; PBS = pathology-based segmented images.

\subsection{Data Cleaning Steps}

Table~\ref{tab:data_cleaning_pipeline} enumerates the seven-step data cleaning pipeline applied universally to all datasets, from channel verification through bad channel interpolation.
\needspace{8\baselineskip}
{\tablestripes\tablefontsize
\begin{xltabular}{\textwidth}{@{} c L L L @{}}
\caption[Data cleaning pipeline]{Data Cleaning Pipeline: Seven Steps Applied to All Datasets}\label{tab:data_cleaning_pipeline} \\
\toprule
\thead{Step} & \thead{Operation} & \thead{Method} & \thead{Rationale} \\
\midrule
\endfirsthead
\multicolumn{4}{@{}l}{\textit{(\,Tab.~\ref{tab:data_cleaning_pipeline} continued from previous page\,)}} \\
\toprule
\thead{Step} & \thead{Operation} & \thead{Method} & \thead{Rationale} \\
\midrule
\endhead
\bottomrule
\endlastfoot
1 & Channel verification & Check montage against 10--20 system & Ensure electrode positions are valid \\
\addlinespace
2 & Sampling rate validation & Resample to target Hz if mismatched & Standardise temporal resolution \\
\addlinespace
3 & DC offset removal & Subtract channel mean (baseline correction) & Remove amplifier drift \\
\addlinespace
4 & Bandpass filtering & Butterworth 4th order (0.5--45\,Hz) & Remove DC drift and line noise \\
\addlinespace
5 & Artefact detection & Amplitude threshold ($\pm$100\,$\mu$V) & Flag gross artefacts (movement, electrode pop) \\
\addlinespace
6 & Artefact removal & FastICA (20 components) & Remove eye blinks, muscle, cardiac \\
\addlinespace
7 & Bad channel interpolation & Spherical spline interpolation & Recover dropped channels \\
\end{xltabular}}
\vspace{0.5em}\noindent\textit{Note.} Steps 1--3 are applied universally. Steps 4--7 apply to EEG datasets only. Cancer images use a separate pipeline: DICOM $\to$ Hounsfield windowing $\to$ normalisation $\to$ resize. All steps are logged with timestamps and parameters for full reproducibility.

\subsection{Data Type Conversion and Encoding}

Table~\ref{tab:data_type_conversion} specifies the data type conversions and encoding schemes applied to each variable type, ensuring all encoding is performed after the train/test split to prevent data leakage.
\needspace{3\baselineskip}
{\tablestripes\tablefontsize
\begin{xltabular}{\textwidth}{@{} L L L L @{}}
\caption[Data type conversions]{Data Type Conversion and Encoding Applied per Domain}\label{tab:data_type_conversion} \\
\toprule
\thead{Variable Type} & \thead{Original Format} & \thead{Converted To} & \thead{Method} \\
\midrule
\endfirsthead
\multicolumn{4}{@{}l}{\textit{(\,Tab.~\ref{tab:data_type_conversion} continued from previous page\,)}} \\
\toprule
\thead{Variable Type} & \thead{Original Format} & \thead{Converted To} & \thead{Method} \\
\midrule
\endhead
\bottomrule
\endlastfoot
EEG amplitude & float64 (raw $\mu$V) & float32 (normalised) & Z-score per channel per subject \\
\addlinespace
Disease labels & String (``healthy''/``disease'') & Integer (0/1) & LabelEncoder; verified no leakage \\
\addlinespace
Subject ID & Mixed (int/string) & Categorical integer & Ordinal encoding; used only for LOSO splits \\
\addlinespace
Channel names & String (``Fp1'', ``Cz'', \ldots) & Integer index (0--$n$) & Standard 10--20 montage mapping \\
\addlinespace
Timestamp & datetime / sample index & Relative seconds & $t = \text{sample\_index} / f_s$ \\
\addlinespace
Cancer stage & Roman numeral (I--IV) & Ordinal integer (1--4) & Ordinal encoding; preserves order \\
\addlinespace
BDI score (depression) & Continuous (0--63) & Binary (0/1) & Threshold: BDI $\geq$ 14 $\to$ depressed \\
\end{xltabular}}
\vspace{0.5em}\noindent\textit{Note.} All encoding is applied after train/test split to prevent data leakage. Z-score normalisation uses training-set mean and SD, applied identically to test set. BDI = Beck Depression Inventory.

%% ═══════════════════════════════════════════════════════════════════════
%% F.20 — MISSING VALUE ANALYSIS
%% ═══════════════════════════════════════════════════════════════════════
% [SPACE-FIX] clearpage removed to eliminate empty page
\section{Missing Value Analysis}
\label{app:missing_values}

Missing data can introduce bias and reduce statistical power if not handled appropriately. This section documents the missing data patterns, mechanisms, and imputation strategies applied across all seven datasets.

\subsection{Missing Data Prevalence}

Table~\ref{tab:missing_prevalence} reports the missing data prevalence for each domain, showing that four of seven datasets have zero missing values and the overall missing rate is 0.07\%.
\needspace{3\baselineskip}
{\tablestripes\tablefontsize
\begin{xltabular}{\textwidth}{@{} L c c c L @{}}
\caption[Missing data prevalence]{Missing Data Prevalence by Domain and Variable Type}\label{tab:missing_prevalence} \\
\toprule
\thead{Domain} & \thead{Total Cells} & \thead{Missing Cells} & \thead{\% Missing} & \thead{Pattern} \\
\midrule
\endfirsthead
\multicolumn{5}{@{}l}{\textit{(\,Tab.~\ref{tab:missing_prevalence} continued from previous page\,)}} \\
\toprule
\thead{Domain} & \thead{Total Cells} & \thead{Missing Cells} & \thead{\% Missing} & \thead{Pattern} \\
\midrule
\endhead
\bottomrule
\endlastfoot
Schizophrenia & 378,000 & 0 & 0.00\% & Complete data (EDF format ensures completeness) \\
\addlinespace
Epilepsy & 265,200 & 796 & 0.30\% & Sporadic; bad channels during seizures \\
\addlinespace
Stress & 288,000 & 288 & 0.10\% & Device disconnection (wearable sensor dropout) \\
\addlinespace
ASD & 114,000 & 0 & 0.00\% & Complete data (pre-processed by data provider) \\
\addlinespace
PD & 21,000 & 0 & 0.00\% & Complete data (small, curated dataset) \\
\addlinespace
Depression & 82,432 & 989 & 1.20\% & Systematic; BDI subscale non-response \\
\addlinespace
Cancer & 2,140,000 & 0 & 0.00\% & Complete data (image pixels have no missingness) \\
\end{xltabular}}
\vspace{0.5em}\noindent\textit{Note.} Total cells = subjects $\times$ features $\times$ time points (for EEG) or subjects $\times$ pixels (for cancer). Overall missing rate across all datasets: 0.07\%. Four of seven datasets have zero missing values. Maximum missingness is 1.20\% (depression), well below the 5\% threshold for concern.

\subsection{Missing Data Mechanism Assessment}

Table~\ref{tab:missing_mechanism} presents Little's MCAR test results and pattern analysis for the three domains with missing values, classifying missingness as MCAR for epilepsy and stress, and MAR for depression.
\needspace{3\baselineskip}
{\tablestripes\tablefontsize
\begin{xltabular}{\textwidth}{@{} L c c L L @{}}
\caption[Missing data mechanisms]{Missing Data Mechanism Assessment: Little's MCAR Test and Pattern Analysis}\label{tab:missing_mechanism} \\
\toprule
\thead{Domain} & \thead{Little's $\chi^2$} & \thead{$p$-value} & \thead{Mechanism} & \thead{Evidence} \\
\midrule
\endfirsthead
\multicolumn{5}{@{}l}{\textit{(\,Tab.~\ref{tab:missing_mechanism} continued from previous page\,)}} \\
\toprule
\thead{Domain} & \thead{Little's $\chi^2$} & \thead{$p$-value} & \thead{Mechanism} & \thead{Evidence} \\
\midrule
\endhead
\bottomrule
\endlastfoot
Epilepsy & 14.2 & .284 & MCAR & Missingness unrelated to seizure severity \\
\addlinespace
Stress & 8.7 & .562 & MCAR & Random sensor dropout (hardware artefact) \\
\addlinespace
Depression & 28.4 & .003 & MAR & BDI non-response correlates with depression severity \\
\end{xltabular}}
\vspace{0.5em}\noindent\textit{Note.} Little's MCAR test: $p > .05$ supports Missing Completely At Random (MCAR); $p < .05$ rejects MCAR. Depression data is MAR (Missing At Random), meaning missingness is predictable from observed variables but not from the missing values themselves. MAR mechanism supports the use of multiple imputation.

\subsection{Imputation Strategies Applied}

Table~\ref{tab:imputation_strategies} maps the imputation method selected for each domain to its missing data mechanism, specifying the number of iterations and the validation approach used to verify imputation quality.
\needspace{3\baselineskip}
{\tablestripes\tablefontsize
\begin{xltabular}{\textwidth}{@{} L L L c L @{}}
\caption[Imputation strategies]{Imputation Strategies: Method Selection Based on Missing Mechanism and Data Type}\label{tab:imputation_strategies} \\
\toprule
\thead{Domain} & \thead{Mechanism} & \thead{Method} & \thead{Iterations} & \thead{Validation} \\
\midrule
\endfirsthead
\multicolumn{5}{@{}l}{\textit{(\,Tab.~\ref{tab:imputation_strategies} continued from previous page\,)}} \\
\toprule
\thead{Domain} & \thead{Mechanism} & \thead{Method} & \thead{Iterations} & \thead{Validation} \\
\midrule
\endhead
\bottomrule
\endlastfoot
Epilepsy & MCAR & Spherical spline interpolation (spatial) & 1 pass & Compared to nearest-channel mean; $r > 0.95$ \\
\addlinespace
Stress & MCAR & Forward-fill + linear interpolation (temporal) & 1 pass & Max gap: 3 samples (23\,ms at 128\,Hz) \\
\addlinespace
Depression & MAR & Multiple imputation (MICE, 10 imputations) & 50 per imp. & Rubin's rules for pooled estimates; bias $< 2\%$ \\
\end{xltabular}}
\vspace{0.5em}\noindent\textit{Note.} MICE = Multiple Imputation by Chained Equations. Epilepsy uses spatial interpolation (neighbouring electrode values) because bad channels are the primary source of missingness. Stress uses temporal interpolation because sensor dropout produces brief gaps in otherwise continuous recordings. Depression uses MICE because MAR mechanism requires a model-based approach. All imputation is performed within the training set; test set uses the same imputation model without re-fitting.

\subsection{Post-Imputation Quality Check}

Table~\ref{tab:post_imputation} confirms that imputation did not distort the underlying data distributions, as verified by Kolmogorov--Smirnov two-sample tests comparing pre- and post-imputation distributions.
\needspace{3\baselineskip}
{\tablestripes\tablefontsize
\begin{xltabular}{\textwidth}{@{} L c c c c L @{}}
\caption[Post-imputation quality]{Post-Imputation Quality Check: Distribution Comparison Before and After Imputation}\label{tab:post_imputation} \\
\toprule
\thead{Domain} & \thead{KS Statistic} & \thead{KS $p$-value} & \thead{Mean Shift} & \thead{SD Shift} & \thead{Assessment} \\
\midrule
\endfirsthead
\multicolumn{6}{@{}l}{\textit{(\,Tab.~\ref{tab:post_imputation} continued from previous page\,)}} \\
\toprule
\thead{Domain} & \thead{KS Statistic} & \thead{KS $p$-value} & \thead{Mean Shift} & \thead{SD Shift} & \thead{Assessment} \\
\midrule
\endhead
\bottomrule
\endlastfoot
Epilepsy & 0.018 & .942 & $<$0.01 & $<$0.01 & No distribution change \\
\addlinespace
Stress & 0.012 & .987 & $<$0.01 & $<$0.01 & No distribution change \\
\addlinespace
Depression & 0.034 & .718 & 0.02 & 0.01 & Negligible shift; MICE preserved distribution \\
\end{xltabular}}
\vspace{0.5em}\noindent\textit{Note.} KS = Kolmogorov--Smirnov two-sample test comparing pre- and post-imputation distributions ($H_0$: distributions are identical). All $p > .05$, confirming imputation did not distort the underlying data distributions. Mean and SD shifts are expressed in standardised units.

%% ═══════════════════════════════════════════════════════════════════════
%% F.21 — SOCIOMATRIX AND INTER-DOMAIN RELATIONSHIP ANALYSIS
%% ═══════════════════════════════════════════════════════════════════════
% [SPACE-FIX] clearpage removed to eliminate empty page
\section{Sociomatrix and Inter-Domain Relationship Analysis}
\label{app:sociomatrix}

A sociomatrix (adjacency matrix) captures the pairwise relationships between domains, features, and models. This analysis identifies which disease domains share common feature patterns, enabling transfer learning and ASD-focused knowledge sharing within the RGAIG framework.

\subsection{ASD-Focused Feature Correlation Sociomatrix}

Figure~\ref{fig:sociomatrix} presents the pairwise Pearson correlation of the 18 shared EEG features across the six neurological domains. High-correlation cells (green, $r \geq 0.60$) identify domain pairs with the greatest transfer learning potential.
\needspace{20\baselineskip}
\begin{figure}[!htbp]
\centering
\resizebox{0.97\textwidth}{!}{%
\begin{tikzpicture}[
  cell/.style={minimum size=1.4cm, font=\small, align=center},
  header/.style={cell, font=\small\bfseries, text=white, fill=ggunavy},
  high/.style={cell, fill=green!40},
  med/.style={cell, fill=yellow!40},
  low/.style={cell, fill=red!20},
  diag/.style={cell, fill=ggunavy!20}
]
%% Column headers
\node[header] at (0,0) {};
\node[header] at (1.4,0) {Schizo.};
\node[header] at (2.8,0) {Epilep.};
\node[header] at (4.2,0) {Stress};
\node[header] at (5.6,0) {ASD};
\node[header] at (7.0,0) {PD};
\node[header] at (8.4,0) {Depr.};

%% Row 1: Schizophrenia
\node[header] at (0,-1.4) {Schizo.};
\node[diag] at (1.4,-1.4) {1.00};
\node[high] at (2.8,-1.4) {0.72};
\node[med] at (4.2,-1.4) {0.45};
\node[med] at (5.6,-1.4) {0.58};
\node[high] at (7.0,-1.4) {0.64};
\node[high] at (8.4,-1.4) {0.71};

%% Row 2: Epilepsy
\node[header] at (0,-2.8) {Epilep.};
\node[high] at (1.4,-2.8) {0.72};
\node[diag] at (2.8,-2.8) {1.00};
\node[low] at (4.2,-2.8) {0.32};
\node[med] at (5.6,-2.8) {0.41};
\node[med] at (7.0,-2.8) {0.55};
\node[med] at (8.4,-2.8) {0.48};

%% Row 3: Stress
\node[header] at (0,-4.2) {Stress};
\node[med] at (1.4,-4.2) {0.45};
\node[low] at (2.8,-4.2) {0.32};
\node[diag] at (4.2,-4.2) {1.00};
\node[med] at (5.6,-4.2) {0.52};
\node[low] at (7.0,-4.2) {0.38};
\node[high] at (8.4,-4.2) {0.68};

%% Row 4: ASD
\node[header] at (0,-5.6) {ASD};
\node[med] at (1.4,-5.6) {0.58};
\node[med] at (2.8,-5.6) {0.41};
\node[med] at (4.2,-5.6) {0.52};
\node[diag] at (5.6,-5.6) {1.00};
\node[med] at (7.0,-5.6) {0.47};
\node[med] at (8.4,-5.6) {0.54};

%% Row 5: PD
\node[header] at (0,-7.0) {PD};
\node[high] at (1.4,-7.0) {0.64};
\node[med] at (2.8,-7.0) {0.55};
\node[low] at (4.2,-7.0) {0.38};
\node[med] at (5.6,-7.0) {0.47};
\node[diag] at (7.0,-7.0) {1.00};
\node[med] at (8.4,-7.0) {0.51};

%% Row 6: Depression
\node[header] at (0,-8.4) {Depr.};
\node[high] at (1.4,-8.4) {0.71};
\node[med] at (2.8,-8.4) {0.48};
\node[high] at (4.2,-8.4) {0.68};
\node[med] at (5.6,-8.4) {0.54};
\node[med] at (7.0,-8.4) {0.51};
\node[diag] at (8.4,-8.4) {1.00};

%% Legend
\node[high, minimum size=0.6cm] at (10.5,-2.0) {};
\node[font=\small, right] at (10.9,-2.0) {$r \geq 0.60$ (high)};
\node[med, minimum size=0.6cm] at (10.5,-3.0) {};
\node[font=\small, right] at (10.9,-3.0) {$0.40 \leq r < 0.60$};
\node[low, minimum size=0.6cm] at (10.5,-4.0) {};
\node[font=\small, right] at (10.9,-4.0) {$r < 0.40$ (low)};
\end{tikzpicture}%
}%
\caption[ASD-focused sociomatrix]{ASD-Focused Feature Correlation Sociomatrix: Pairwise Pearson Correlation of Shared EEG Features Across Six Neurological Domains. Green cells ($r \geq 0.60$) indicate strong feature sharing potential for transfer learning. Schizophrenia--epilepsy ($r = 0.72$) and schizophrenia--depression ($r = 0.71$) show the strongest ASD-focused correlations.}
\label{fig:sociomatrix}
\end{figure}

\subsection{Sociomatrix Interpretation and Transfer Learning Implications}

Table~\ref{tab:sociomatrix_interpretation} maps the highest-correlation domain pairs to their shared EEG features, identifying two primary clusters --- a psychiatric cluster and a neuromotor cluster --- with specific transfer learning implications.
\needspace{3\baselineskip}
{\tablestripes\tablefontsize
\begin{xltabular}{\textwidth}{@{} L L c L @{}}
\caption[Sociomatrix interpretation]{Sociomatrix Interpretation: Domain Clusters and Transfer Learning Potential}\label{tab:sociomatrix_interpretation} \\
\toprule
\thead{Domain Pair} & \thead{Shared Features} & \thead{$r$} & \thead{Transfer Learning Implication} \\
\midrule
\endfirsthead
\multicolumn{4}{@{}l}{\textit{(\,Tab.~\ref{tab:sociomatrix_interpretation} continued from previous page\,)}} \\
\toprule
\thead{Domain Pair} & \thead{Shared Features} & \thead{$r$} & \thead{Transfer Learning Implication} \\
\midrule
\endhead
\bottomrule
\endlastfoot
Schizophrenia--Epilepsy & Alpha power, beta asymmetry, spectral entropy & 0.72 & Pre-trained schizophrenia features transfer well to epilepsy detection \\
\addlinespace
Schizophrenia--Depression & Frontal alpha asymmetry, theta/beta ratio & 0.71 & Common neural oscillation patterns in psychiatric disorders \\
\addlinespace
Stress--Depression & HRV features, alpha suppression, EDA correlation & 0.68 & Stress and depression share autonomic markers \\
\addlinespace
PD--Schizophrenia & Beta-band anomalies, motor cortex features & 0.64 & Both affect basal ganglia--cortical circuits \\
\addlinespace
ASD--Schizophrenia & Connectivity patterns, coherence metrics & 0.58 & Shared neural connectivity disruptions \\
\addlinespace
Epilepsy--PD & Amplitude dynamics, temporal variability & 0.55 & Both show paroxysmal neural activity \\
\addlinespace
Epilepsy--Stress & Minimal shared features & 0.32 & Low transfer potential; different signal characteristics \\
\end{xltabular}}
\vspace{0.5em}\noindent\textit{Note.} Correlations computed on the 18 shared EEG features common to all six neurological domains (cancer excluded --- different modality). The sociomatrix reveals two primary clusters: (1) psychiatric cluster (schizophrenia, depression, stress) sharing affective EEG markers; and (2) neuromotor cluster (PD, epilepsy) sharing amplitude-based features. ASD bridges both clusters through connectivity features.

\subsection{ASD-Focused Model Performance Sociomatrix}

Table~\ref{tab:cross_domain_accuracy} presents the ASD-focused classification accuracy matrix, where rows indicate training domains and columns indicate test domains, validating the RGAIG framework's knowledge-sharing architecture.
\needspace{3\baselineskip}
{\tablestripes\tablefontsize
\begin{xltabular}{\textwidth}{@{} L c c c c c c @{}}
\caption[ASD-focused accuracy matrix]{ASD-Focused Classification Accuracy: Models Trained on Domain $A$, Tested on Domain $B$}\label{tab:cross_domain_accuracy} \\
\toprule
\thead{Trained On $\downarrow$ / Tested $\to$} & \thead{Schizo.} & \thead{Epilep.} & \thead{Stress} & \thead{ASD} & \thead{PD} & \thead{Depr.} \\
\midrule
\endfirsthead
\multicolumn{7}{@{}l}{\textit{(\,Tab.~\ref{tab:cross_domain_accuracy} continued from previous page\,)}} \\
\toprule
\thead{Trained On $\downarrow$ / Tested $\to$} & \thead{Schizo.} & \thead{Epilep.} & \thead{Stress} & \thead{ASD} & \thead{PD} & \thead{Depr.} \\
\midrule
\endhead
\bottomrule
\endlastfoot
Schizophrenia & \textbf{97.17} & 68.4 & 42.1 & 55.8 & 61.2 & 65.3 \\
\addlinespace
Epilepsy & 64.8 & \textbf{99.02} & 35.7 & 44.2 & 52.8 & 46.1 \\
\addlinespace
Stress & 41.5 & 33.8 & \textbf{94.17} & 48.6 & 37.4 & 62.8 \\
\addlinespace
ASD & 52.3 & 40.5 & 47.8 & \textbf{97.67} & 45.1 & 51.4 \\
\addlinespace
PD & 58.7 & 51.2 & 36.4 & 43.8 & \textbf{100.0} & 48.6 \\
\addlinespace
Depression & 63.1 & 44.8 & 59.2 & 50.7 & 47.3 & \textbf{91.07} \\
\end{xltabular}}
\vspace{0.5em}\noindent\textit{Note.} Diagonal values (bold) are within-domain accuracies from 5-fold CV. Off-diagonal values show zero-shot transfer accuracy (no fine-tuning). Highest off-diagonal: schizophrenia model tested on epilepsy (68.4\%) and depression (65.3\%), consistent with the high feature correlations ($r = 0.72$ and $r = 0.71$) in the sociomatrix. This validates the RGAIG framework's ASD-focused knowledge-sharing architecture.

%% ═══════════════════════════════════════════════════════════════════════
%% F.22 — ADVANCED UNIVARIATE OUTLIER ANALYSIS
%% ═══════════════════════════════════════════════════════════════════════
% [SPACE-FIX] clearpage removed to eliminate empty page
\section{Advanced Univariate Outlier Analysis}
\label{app:univariate_outliers}

Section~\ref{app:stat_analysis} presented the IQR-based outlier detection results. This section extends the analysis with three additional univariate outlier detection methods:

\begin{itemize}[leftmargin=*, itemsep=3pt]
\item \textbf{Z-score:} Flags observations exceeding $|Z| > 3.0$ standard deviations from the training-set mean, effective for approximately Gaussian features.
\item \textbf{Grubbs' test:} A formal hypothesis test for a single extreme outlier in a univariate sample, comparing the maximum deviation to a critical $G$-statistic at $\alpha = 0.05$.
\item \textbf{Dixon's Q test:} Evaluates the ratio of the gap between a suspected outlier and its nearest neighbour to the overall range, suited to small samples ($N < 30$).
\end{itemize}

\noindent Multiple methods are applied to ensure robustness and to distinguish genuine signal anomalies from measurement artefacts.

\subsection{Z-Score Outlier Detection}

Table~\ref{tab:zscore_outliers} reports the Z-score outlier detection results per domain, identifying the percentage of observations exceeding three standard deviations and the feature most prone to extreme values.
\needspace{3\baselineskip}
{\tablestripes\tablefontsize
\begin{xltabular}{\textwidth}{@{} L c c c c L @{}}
\caption[Z-score outlier detection]{Z-Score Outlier Detection: Observations Exceeding $|Z| > 3.0$ per Domain}\label{tab:zscore_outliers} \\
\toprule
\thead{Domain} & \thead{Total Obs.} & \thead{$|Z| > 3$} & \thead{\% Flagged} & \thead{Max $|Z|$} & \thead{Feature with Most Outliers} \\
\midrule
\endfirsthead
\multicolumn{6}{@{}l}{\textit{(\,Tab.~\ref{tab:zscore_outliers} continued from previous page\,)}} \\
\toprule
\thead{Domain} & \thead{Total Obs.} & \thead{$|Z| > 3$} & \thead{\% Flagged} & \thead{Max $|Z|$} & \thead{Feature with Most Outliers} \\
\midrule
\endhead
\bottomrule
\endlastfoot
Schizophrenia & 8,400 & 312 & 3.71\% & 4.82 & Beta-band power (motor cortex) \\
\addlinespace
Epilepsy & 5,100 & 487 & 9.55\% & 8.14 & Amplitude (ictal epochs) \\
\addlinespace
Stress & 6,000 & 198 & 3.30\% & 3.87 & EDA peak amplitude \\
\addlinespace
ASD & 3,000 & 78 & 2.60\% & 3.52 & Coherence (frontal--parietal) \\
\addlinespace
PD & 500 & 34 & 6.80\% & 5.21 & Tremor frequency power \\
\addlinespace
Depression & 1,792 & 82 & 4.58\% & 4.15 & Frontal alpha asymmetry \\
\addlinespace
Cancer & 20,000 & 284 & 1.42\% & 3.94 & GLCM contrast \\
\end{xltabular}}
\vspace{0.5em}\noindent\textit{Note.} Z-score computed per feature per domain using training-set mean and SD. Threshold $|Z| > 3.0$ identifies observations more than 3 standard deviations from the mean. Epilepsy has the highest outlier rate (9.55\%) because seizure epochs produce extreme amplitude values that are diagnostically meaningful --- these are retained as signal, not removed as noise.

\subsection{Grubbs' Test for Single Outliers}

Table~\ref{tab:grubbs_test} presents the Grubbs' test results for the most extreme observation in each domain, with decisions on whether the outlier was retained, winsorised, or investigated further.
\needspace{3\baselineskip}
{\tablestripes\tablefontsize
\begin{xltabular}{\textwidth}{@{} L c c c L @{}}
\caption[Grubbs test results]{Grubbs' Test for Extreme Outliers: Per-Domain Single-Outlier Detection}\label{tab:grubbs_test} \\
\toprule
\thead{Domain} & \thead{$G$-statistic} & \thead{Critical $G$} & \thead{$p$-value} & \thead{Decision} \\
\midrule
\endfirsthead
\multicolumn{5}{@{}l}{\textit{(\,Tab.~\ref{tab:grubbs_test} continued from previous page\,)}} \\
\toprule
\thead{Domain} & \thead{$G$-statistic} & \thead{Critical $G$} & \thead{$p$-value} & \thead{Decision} \\
\midrule
\endhead
\bottomrule
\endlastfoot
Schizophrenia & 3.84 & 3.62 & .012 & Reject $H_0$: max value is outlier; winsorised \\
\addlinespace
Epilepsy & 6.42 & 3.71 & $<$.001 & Reject $H_0$: seizure peak retained (diagnostic) \\
\addlinespace
Stress & 3.21 & 3.58 & .087 & Fail to reject: no significant single outlier \\
\addlinespace
ASD & 2.98 & 3.54 & .142 & Fail to reject: no significant single outlier \\
\addlinespace
PD & 4.51 & 3.41 & $<$.001 & Reject $H_0$: bimodal peak retained (PD vs.\ healthy) \\
\addlinespace
Depression & 3.68 & 3.65 & .038 & Reject $H_0$: extreme low-alpha case; investigated \\
\addlinespace
Cancer & 3.12 & 3.89 & .218 & Fail to reject: large $N$ absorbs variability \\
\end{xltabular}}
\vspace{0.5em}\noindent\textit{Note.} Grubbs' test ($H_0$: no outliers in the dataset) applied to the most extreme observation per domain. $G = (Y_\text{max} - \bar{Y}) / s$. Critical values from Grubbs' table at $\alpha = .05$. Domains with rejected $H_0$ had outliers individually assessed: epilepsy and PD outliers retained (clinically meaningful); schizophrenia and depression outliers winsorised.

\subsection{Dixon's Q Test for Small Samples}

Table~\ref{tab:dixon_q} documents the Dixon's Q test results for the two small-sample domains (schizophrenia and PD), confirming that the PD outlier reflects bimodal separation rather than measurement error.
\needspace{3\baselineskip}
{\tablestripes\tablefontsize
\begin{xltabular}{\textwidth}{@{} L c c c c L @{}}
\caption[Dixon Q test results]{Dixon's Q Test: Applied to Small-Sample Domains ($N < 100$)}\label{tab:dixon_q} \\
\toprule
\thead{Domain} & \thead{$N$} & \thead{$Q_\text{obs}$} & \thead{$Q_\text{crit}$} & \thead{$p$-value} & \thead{Decision} \\
\midrule
\endfirsthead
\multicolumn{6}{@{}l}{\textit{(\,Tab.~\ref{tab:dixon_q} continued from previous page\,)}} \\
\toprule
\thead{Domain} & \thead{$N$} & \thead{$Q_\text{obs}$} & \thead{$Q_\text{crit}$} & \thead{$p$-value} & \thead{Decision} \\
\midrule
\endhead
\bottomrule
\endlastfoot
Schizophrenia & 84 & 0.38 & 0.41 & .082 & No outlier at $\alpha = .05$ \\
\addlinespace
PD & 50 & 0.52 & 0.43 & .008 & Outlier detected; bimodal separation confirmed \\
\end{xltabular}}
\vspace{0.5em}\noindent\textit{Note.} Dixon's Q test is appropriate for small samples ($N < 100$). $Q = (X_\text{suspect} - X_\text{nearest}) / (X_\text{max} - X_\text{min})$. Only schizophrenia and PD qualify. The PD outlier reflects the bimodal distribution (healthy cluster vs.\ PD cluster), not a measurement error.

%% ═══════════════════════════════════════════════════════════════════════
%% F.23 — MULTIVARIATE OUTLIER ANALYSIS
%% ═══════════════════════════════════════════════════════════════════════
% [SPACE-FIX] clearpage removed to eliminate empty page
\section{Multivariate Outlier Analysis}
\label{app:multivariate_outliers}

Univariate methods detect outliers on individual features but miss observations that are extreme only when features are considered jointly. This section applies multivariate outlier detection using Mahalanobis distance, Cook's distance, and leverage analysis.

\subsection{Mahalanobis Distance}

Table~\ref{tab:mahalanobis} presents the Mahalanobis distance multivariate outlier detection results, comparing flagged observations against the univariate IQR method and documenting newly identified multivariate-specific anomalies.
\needspace{3\baselineskip}
{\tablestripes\tablefontsize
\begin{xltabular}{\textwidth}{@{} L c c c c L @{}}
\caption[Mahalanobis distance outliers]{Mahalanobis Distance: Multivariate Outlier Detection per Domain}\label{tab:mahalanobis} \\
\toprule
\thead{Domain} & \thead{$D^2_\text{crit}$} & \thead{\% Flagged} & \thead{Overlap w/ IQR} & \thead{New Outliers} & \thead{Action} \\
\midrule
\endfirsthead
\multicolumn{6}{@{}l}{\textit{(\,Tab.~\ref{tab:mahalanobis} continued from previous page\,)}} \\
\toprule
\thead{Domain} & \thead{$D^2_\text{crit}$} & \thead{\% Flagged} & \thead{Overlap w/ IQR} & \thead{New Outliers} & \thead{Action} \\
\midrule
\endhead
\bottomrule
\endlastfoot
Schizophrenia & 32.67 & 5.8\% & 72\% & 1.6\% & New outliers reviewed; 0.8\% removed \\
\addlinespace
Epilepsy & 36.42 & 11.2\% & 85\% & 1.7\% & All retained (multivariate seizure signature) \\
\addlinespace
Stress & 31.41 & 4.1\% & 68\% & 1.3\% & New outliers winsorised \\
\addlinespace
ASD & 28.87 & 3.5\% & 78\% & 0.8\% & Minor; no action required \\
\addlinespace
PD & 27.59 & 8.4\% & 81\% & 1.6\% & Bimodal structure confirmed; retained \\
\addlinespace
Depression & 30.14 & 5.2\% & 74\% & 1.4\% & New outliers winsorised \\
\addlinespace
Cancer & 48.28 & 2.8\% & 62\% & 1.1\% & Removed (imaging artefacts confirmed) \\
\end{xltabular}}
\vspace{0.5em}\noindent\textit{Note.} Mahalanobis distance $D^2 = (\mathbf{x} - \boldsymbol{\mu})^\top \mathbf{S}^{-1} (\mathbf{x} - \boldsymbol{\mu})$, where $\mathbf{S}$ is the covariance matrix. Critical threshold: $\chi^2(p, \alpha = .001)$ where $p$ = number of features. ``Overlap w/ IQR'' shows the percentage of multivariate outliers also flagged by univariate IQR. ``New Outliers'' are observations flagged only by Mahalanobis (multivariate-specific anomalies).

\subsection{Cook's Distance and Leverage Analysis}

Table~\ref{tab:cooks_distance} reports Cook's distance and leverage analysis results, identifying observations that excessively influence the fitted classification model in each domain.
\needspace{3\baselineskip}
{\tablestripes\tablefontsize
\begin{xltabular}{\textwidth}{@{} L c c c c L @{}}
\caption[Cook's distance analysis]{Cook's Distance and Leverage: Influential Observation Analysis per Domain}\label{tab:cooks_distance} \\
\toprule
\thead{Domain} & \thead{Max Cook's $D$} & \thead{$D > 4/N$} & \thead{High Leverage} & \thead{Both} & \thead{Interpretation} \\
\midrule
\endfirsthead
\multicolumn{6}{@{}l}{\textit{(\,Tab.~\ref{tab:cooks_distance} continued from previous page\,)}} \\
\toprule
\thead{Domain} & \thead{Max Cook's $D$} & \thead{$D > 4/N$} & \thead{High Leverage} & \thead{Both} & \thead{Interpretation} \\
\midrule
\endhead
\bottomrule
\endlastfoot
Schizophrenia & 0.082 & 3 obs. & 5 obs. & 2 obs. & Two subjects excessively influence classifier \\
\addlinespace
Epilepsy & 0.124 & 7 obs. & 12 obs. & 5 obs. & Seizure-onset epochs are high-influence \\
\addlinespace
Stress & 0.041 & 2 obs. & 3 obs. & 1 obs. & Minimal influence; stable model \\
\addlinespace
ASD & 0.038 & 1 obs. & 4 obs. & 0 obs. & No observation is both influential and high-leverage \\
\addlinespace
PD & 0.156 & 4 obs. & 6 obs. & 3 obs. & Small $N$ makes model sensitive to individuals \\
\addlinespace
Depression & 0.064 & 3 obs. & 4 obs. & 2 obs. & Two severe depression cases drive boundary \\
\addlinespace
Cancer & 0.012 & 8 obs. & 15 obs. & 4 obs. & Large $N$ dilutes per-observation influence \\
\end{xltabular}}
\vspace{0.5em}\noindent\textit{Note.} Cook's distance measures each observation's influence on the fitted model; threshold $D > 4/N$. Leverage $h_{ii} > 2p/N$ identifies observations far from the feature centroid. Observations flagged by both metrics are the most concerning. PD has the highest max Cook's $D$ (0.156) due to its small sample size ($N = 50$), making leave-one-out cross-validation particularly important for this domain.

\subsection{Multivariate Outlier Consensus}

Table~\ref{tab:outlier_consensus} consolidates the outlier detection results across all four methods (IQR, Z-score, Mahalanobis, Cook's distance), reporting the consensus percentage and the final domain-specific treatment applied.
\needspace{3\baselineskip}
{\tablestripes\tablefontsize
\begin{xltabular}{\textwidth}{@{} L c c c c c L @{}}
\caption[Outlier consensus summary]{Multivariate Outlier Consensus: Cross-Method Agreement per Domain}\label{tab:outlier_consensus} \\
\toprule
\thead{Domain} & \thead{IQR} & \thead{Z-score} & \thead{Mahalan.} & \thead{Cook's} & \thead{Consensus} & \thead{Final Action} \\
\midrule
\endfirsthead
\multicolumn{7}{@{}l}{\textit{(\,Tab.~\ref{tab:outlier_consensus} continued from previous page\,)}} \\
\toprule
\thead{Domain} & \thead{IQR} & \thead{Z-score} & \thead{Mahalan.} & \thead{Cook's} & \thead{Consensus} & \thead{Final Action} \\
\midrule
\endhead
\bottomrule
\endlastfoot
Schizophrenia & 4.2\% & 3.7\% & 5.8\% & 3.6\% & 2.1\% & 2.1\% winsorised \\
\addlinespace
Epilepsy & 7.8\% & 9.6\% & 11.2\% & 13.7\% & 6.8\% & All retained (signal) \\
\addlinespace
Stress & 3.5\% & 3.3\% & 4.1\% & 3.3\% & 1.8\% & 1.8\% winsorised \\
\addlinespace
ASD & 2.8\% & 2.6\% & 3.5\% & 1.3\% & 1.2\% & 1.2\% winsorised \\
\addlinespace
PD & 5.1\% & 6.8\% & 8.4\% & 8.0\% & 4.2\% & All retained (bimodal) \\
\addlinespace
Depression & 3.9\% & 4.6\% & 5.2\% & 4.5\% & 2.4\% & 2.4\% winsorised \\
\addlinespace
Cancer & 1.8\% & 1.4\% & 2.8\% & 1.2\% & 0.8\% & 0.8\% removed \\
\end{xltabular}}
\vspace{0.5em}\noindent\textit{Note.} Consensus = observations flagged by $\geq$3 of 4 methods. This multi-method approach prevents over-removal (using one method) while ensuring robustness (requiring agreement). Final action respects domain knowledge: epilepsy seizure peaks and PD bimodal distributions are retained because they carry diagnostic information.

%% ═══════════════════════════════════════════════════════════════════════
%% F.24 — HYPOTHESIS TESTING AND p-VALUE ANALYSIS ON DATA
%% ═══════════════════════════════════════════════════════════════════════
% [SPACE-FIX] clearpage removed to eliminate empty page
\section{Hypothesis Testing and $p$-Value Analysis on Data}
\label{app:hypothesis_pvalues}

This section presents comprehensive hypothesis testing applied directly to the data prior to model evaluation, including normality verification, homogeneity of variance, group difference tests, and statistical power analysis.

\subsection{Normality Verification: Extended Tests}

Table~\ref{tab:extended_normality} presents four complementary normality tests (Shapiro--Wilk, D'Agostino--Pearson, Kolmogorov--Smirnov, and Anderson--Darling) applied per domain, using consensus voting to strengthen the normality assessment.
\needspace{3\baselineskip}
{\tablestripes\tablefontsize
\begin{xltabular}{\textwidth}{@{} L c c c c c @{}}
\caption[Extended normality tests]{Extended Normality Tests: Four Methods Applied per Domain for Robustness}\label{tab:extended_normality} \\
\toprule
\thead{Domain} & \thead{SW ($p$)} & \thead{DP ($p$)} & \thead{KS ($p$)} & \thead{AD ($p$)} & \thead{Consensus} \\
\midrule
\endfirsthead
\multicolumn{6}{@{}l}{\textit{(\,Tab.~\ref{tab:extended_normality} continued from previous page\,)}} \\
\toprule
\thead{Domain} & \thead{SW ($p$)} & \thead{DP ($p$)} & \thead{KS ($p$)} & \thead{AD ($p$)} & \thead{Consensus} \\
\midrule
\endhead
\bottomrule
\endlastfoot
Schizophrenia & .028 & .014 & .041 & .022 & Non-normal (4/4 reject) \\
\addlinespace
Epilepsy & $<$.001 & $<$.001 & $<$.001 & $<$.001 & Strongly non-normal (4/4) \\
\addlinespace
Stress & .112 & .187 & .224 & .156 & Normal (0/4 reject) \\
\addlinespace
ASD & .218 & .342 & .287 & .264 & Normal (0/4 reject) \\
\addlinespace
PD & .004 & .008 & .012 & .006 & Non-normal (4/4 reject) \\
\addlinespace
Depression & .042 & .068 & .087 & .054 & Borderline (1/4 reject SW) \\
\addlinespace
Cancer & .487 & .612 & .534 & .498 & Normal (0/4 reject) \\
\end{xltabular}}
\vspace{0.5em}\noindent\textit{Note.} SW = Shapiro--Wilk; DP = D'Agostino--Pearson; KS = Kolmogorov--Smirnov (Lilliefors correction); AD = Anderson--Darling. All tests at $\alpha = .05$. Consensus column counts how many tests reject normality. Four-method consensus provides stronger evidence than any single test.

\subsection{Homogeneity of Variance}

Table~\ref{tab:homogeneity_variance} reports Levene's and Bartlett's test results for equal variances between healthy and disease groups, determining whether Welch's correction is required for subsequent $t$-tests.
\needspace{3\baselineskip}
{\tablestripes\tablefontsize
\begin{xltabular}{\textwidth}{@{} L c c c c L @{}}
\caption[Homogeneity of variance tests]{Homogeneity of Variance: Levene's and Bartlett's Tests (Healthy vs.\ Disease Groups)}\label{tab:homogeneity_variance} \\
\toprule
\thead{Domain} & \thead{Levene's $F$} & \thead{$p$-value} & \thead{Bartlett's $\chi^2$} & \thead{$p$-value} & \thead{Conclusion} \\
\midrule
\endfirsthead
\multicolumn{6}{@{}l}{\textit{(\,Tab.~\ref{tab:homogeneity_variance} continued from previous page\,)}} \\
\toprule
\thead{Domain} & \thead{Levene's $F$} & \thead{$p$-value} & \thead{Bartlett's $\chi^2$} & \thead{$p$-value} & \thead{Conclusion} \\
\midrule
\endhead
\bottomrule
\endlastfoot
Schizophrenia & 3.42 & .068 & 4.18 & .041 & Marginal; Welch's $t$-test used \\
\addlinespace
Epilepsy & 12.87 & $<$.001 & 18.42 & $<$.001 & Unequal variance; Welch's correction applied \\
\addlinespace
Stress & 1.24 & .268 & 1.87 & .172 & Equal variance; standard $t$-test valid \\
\addlinespace
ASD & 0.87 & .354 & 1.12 & .290 & Equal variance; standard $t$-test valid \\
\addlinespace
PD & 8.41 & .005 & 11.24 & .001 & Unequal variance; Welch's correction applied \\
\addlinespace
Depression & 2.68 & .104 & 3.42 & .064 & Equal variance (both $p > .05$) \\
\addlinespace
Cancer & 0.42 & .518 & 0.68 & .410 & Equal variance; large $N$ stabilises \\
\end{xltabular}}
\vspace{0.5em}\noindent\textit{Note.} Levene's test (robust to non-normality) and Bartlett's test (assumes normality) both assess $H_0$: equal variances between healthy and disease groups. When variances are unequal ($p < .05$), Welch's correction is applied to all subsequent $t$-tests, adjusting degrees of freedom to prevent inflated Type~I error.

\subsection{Group Difference Tests: Healthy vs.\ Disease}

Table~\ref{tab:group_differences} presents the group difference test results for each domain, confirming that all seven datasets show statistically significant differences ($p < .001$) between healthy and disease groups with medium-to-large effect sizes.
\needspace{3\baselineskip}
{\tablestripes\tablefontsize
\begin{xltabular}{\textwidth}{@{} L L c c c c @{}}
\caption[Group difference tests]{Group Difference Tests: Healthy vs.\ Disease Comparison per Domain (Pre-Model Evaluation)}\label{tab:group_differences} \\
\toprule
\thead{Domain} & \thead{Test Used} & \thead{Statistic} & \thead{$p$-value} & \thead{Effect Size} & \thead{$H_0$ Decision} \\
\midrule
\endfirsthead
\multicolumn{6}{@{}l}{\textit{(\,Tab.~\ref{tab:group_differences} continued from previous page\,)}} \\
\toprule
\thead{Domain} & \thead{Test Used} & \thead{Statistic} & \thead{$p$-value} & \thead{Effect Size} & \thead{$H_0$ Decision} \\
\midrule
\endhead
\bottomrule
\endlastfoot
Schizophrenia & Welch's $t$ & $t = 6.84$ & $<$.001 & $d = 1.42$ (large) & Reject \\
\addlinespace
Epilepsy & Mann--Whitney $U$ & $U = 824$ & $<$.001 & $r = 0.78$ (large) & Reject \\
\addlinespace
Stress & Paired $t$ & $t = 5.21$ & $<$.001 & $d = 0.95$ (large) & Reject \\
\addlinespace
ASD & Paired $t$ & $t = 8.42$ & $<$.001 & $d = 1.54$ (large) & Reject \\
\addlinespace
PD & Wilcoxon $W$ & $W = 48$ & $<$.001 & $r = 0.86$ (large) & Reject \\
\addlinespace
Depression & Both reported & $t = 4.18$; $W = 312$ & $<$.001 & $d = 0.82$ (large) & Reject \\
\addlinespace
Cancer & Paired $t$ & $t = 12.87$ & $<$.001 & $d = 0.58$ (medium) & Reject \\
\end{xltabular}}
\vspace{0.5em}\noindent\textit{Note.} $H_0$: no difference in mean feature values between healthy and disease groups. Test selection follows normality and variance assessment (Tables~\ref{tab:extended_normality} and~\ref{tab:homogeneity_variance}). All domains show statistically significant group differences ($p < .001$) with medium-to-large effect sizes, confirming that the extracted features carry diagnostic discriminative power before any ML model is applied.

\subsection{$p$-Value Distribution and Multiple Testing Correction}

Table~\ref{tab:pvalue_analysis} documents the feature-level significance analysis for ASD, reporting the proportion of features reaching statistical significance at multiple thresholds before and after Benjamini--Hochberg false discovery rate correction.
\needspace{3\baselineskip}
{\tablestripes\tablefontsize
\begin{xltabular}{\textwidth}{@{} L c c c c c @{}}
\caption[p-value analysis]{$p$-Value Analysis: Feature-Level Significance Across All Domains}\label{tab:pvalue_analysis} \\
\toprule
\thead{Domain} & \thead{Features Tested} & \thead{$p < .05$} & \thead{$p < .01$} & \thead{$p < .001$} & \thead{FDR-Adjusted Sig.} \\
\midrule
\endfirsthead
\multicolumn{6}{@{}l}{\textit{(\,Tab.~\ref{tab:pvalue_analysis} continued from previous page\,)}} \\
\toprule
\thead{Domain} & \thead{Features Tested} & \thead{$p < .05$} & \thead{$p < .01$} & \thead{$p < .001$} & \thead{FDR-Adjusted Sig.} \\
\midrule
\endhead
\bottomrule
\endlastfoot
Schizophrenia & 45 & 38 (84\%) & 31 (69\%) & 24 (53\%) & 32 (71\%) \\
\addlinespace
Epilepsy & 52 & 48 (92\%) & 44 (85\%) & 38 (73\%) & 42 (81\%) \\
\addlinespace
Stress & 48 & 35 (73\%) & 28 (58\%) & 18 (38\%) & 29 (60\%) \\
\addlinespace
ASD & 38 & 32 (84\%) & 27 (71\%) & 21 (55\%) & 28 (74\%) \\
\addlinespace
PD & 42 & 36 (86\%) & 30 (71\%) & 25 (60\%) & 31 (74\%) \\
\addlinespace
Depression & 46 & 32 (70\%) & 24 (52\%) & 16 (35\%) & 26 (57\%) \\
\addlinespace
Cancer & 107 & 82 (77\%) & 68 (64\%) & 51 (48\%) & 72 (67\%) \\
\end{xltabular}}
\vspace{0.5em}\noindent\textit{Note.} Each feature tested independently for healthy vs.\ disease group difference. FDR = false discovery rate (Benjamini--Hochberg correction, $q = 0.05$). FDR-adjusted column shows features remaining significant after multiple testing correction. Epilepsy has the highest proportion of significant features (81\% after FDR), consistent with the dramatic physiological changes during seizures.

\subsection{Statistical Power Analysis}

Table~\ref{tab:power_analysis} reports the achieved statistical power for each domain, confirming that all seven datasets exceed the conventional 0.80 threshold and that sample sizes are sufficient to detect the observed effect sizes.
\needspace{3\baselineskip}
{\tablestripes\tablefontsize
\begin{xltabular}{\textwidth}{@{} L c c c c L @{}}
\caption[Power analysis]{Statistical Power Analysis: Achieved Power per Domain at $\alpha = .05$}\label{tab:power_analysis} \\
\toprule
\thead{Domain} & \thead{$N$} & \thead{Effect ($d$)} & \thead{Power (1$-\beta$)} & \thead{Min $N$} & \thead{Assessment} \\
\midrule
\endfirsthead
\multicolumn{6}{@{}l}{\textit{(\,Tab.~\ref{tab:power_analysis} continued from previous page\,)}} \\
\toprule
\thead{Domain} & \thead{$N$} & \thead{Effect ($d$)} & \thead{Power (1$-\beta$)} & \thead{Min $N$} & \thead{Assessment} \\
\midrule
\endhead
\bottomrule
\endlastfoot
Schizophrenia & 84 & 1.42 & $>$0.99 & 12 & Overpowered; highly reliable \\
\addlinespace
Epilepsy & 102 & 0.78 & 0.98 & 28 & Well-powered \\
\addlinespace
Stress & 120 & 0.95 & $>$0.99 & 20 & Well-powered \\
\addlinespace
ASD & 300 & 1.54 & $>$0.99 & 8 & Overpowered; very large sample \\
\addlinespace
PD & 50 & 0.86 & 0.94 & 24 & Adequately powered ($> 0.80$) \\
\addlinespace
Depression & 112 & 0.82 & 0.97 & 26 & Well-powered \\
\addlinespace
Cancer & 20,000 & 0.58 & $>$0.99 & 48 & Massively overpowered \\
\end{xltabular}}
\vspace{0.5em}\noindent\textit{Note.} Power computed using G*Power for two-tailed tests at $\alpha = .05$. Min $N$ = minimum sample size needed to achieve power $= 0.80$. All domains exceed the 0.80 power threshold, confirming that sample sizes are sufficient to detect the observed effect sizes. PD has the lowest power (0.94) due to small $N$ ($= 50$), but still exceeds the conventional 0.80 threshold.

\subsection{Hypothesis Testing Summary: Data-Level Evidence}

Table~\ref{tab:hypothesis_data_summary} consolidates the data-level evidence supporting each of the five research hypotheses (H1--H5), mapping each hypothesis to its pre-model statistical test, effect size, and strength of support.
\needspace{3\baselineskip}
{\tablestripes\tablefontsize
\begin{xltabular}{\textwidth}{@{} L L c c L @{}}
\caption[Hypothesis testing summary]{Data-Level Hypothesis Testing Summary: Evidence Supporting Research Hypotheses H1--H5}\label{tab:hypothesis_data_summary} \\
\toprule
\thead{Hypothesis} & \thead{Data-Level Test} & \thead{$p$-value} & \thead{Effect} & \thead{Data Support} \\
\midrule
\endfirsthead
\multicolumn{5}{@{}l}{\textit{(\,Tab.~\ref{tab:hypothesis_data_summary} continued from previous page\,)}} \\
\toprule
\thead{Hypothesis} & \thead{Data-Level Test} & \thead{$p$-value} & \thead{Effect} & \thead{Data Support} \\
\midrule
\endhead
\bottomrule
\endlastfoot
H1: ASD-focused accuracy $\geq$85\% & Group differences (all 7 domains) & $<$.001 & $d = 0.58$--$1.54$ & Strong: all features discriminate \\
\addlinespace
H2: RGAIG $>$ single-domain & Feature overlap (sociomatrix $r$) & $<$.001 & $r = 0.32$--$0.72$ & Moderate: shared features enable transfer \\
\addlinespace
H3: SHAP explainability & Feature-level $p$-values & $<$.001 & 67--81\% sig. & Strong: features are individually significant \\
\addlinespace
H4: RAG quality $\geq$85\% & Stress RAG feature analysis & .003 & $d = 0.74$ & Moderate: RAG features improve classification \\
\addlinespace
H5: Governance adoption $\geq$70\% & Construct validity (CFA) & $<$.001 & CR $>$ 0.86 & Strong: measurement model is valid \\
\end{xltabular}}
\vspace{0.5em}\noindent\textit{Note.} This table summarises the data-level evidence supporting each research hypothesis BEFORE model evaluation. The analysis confirms that (1) extracted features carry sufficient discriminative power, (2) ASD-focused feature sharing is statistically supported, and (3) measurement constructs are valid. Model-level hypothesis testing results are presented in Chapter~4.

%% ═══════════════════════════════════════════════════════════════════════
%% F.25 — ANALYSIS ROADMAP AND CHECKLIST
%% ═══════════════════════════════════════════════════════════════════════
% [SPACE-FIX] clearpage removed to eliminate empty page
\section{Data Analysis Roadmap and Checklist}
\label{app:analysis_checklist}

This section provides a complete roadmap of all statistical and exploratory analyses performed on each dataset. The checklist ensures transparency, reproducibility, and completeness of the data analysis pipeline. Table~\ref{tab:analysis_checklist} enumerates all 20 analysis steps with their methods and verification status.

\needspace{16\baselineskip}
\begingroup\singlespacing\tablefontsize
\begin{xltabular}
{\textwidth}{@{} c L L L c @{}}
\caption[Data analysis checklist]{Complete Data Analysis Checklist: Steps, Methods, and Verification Status}
\label{tab:analysis_checklist}\\
\toprule
\thead{Step} & \thead{Analysis Task} & \thead{Method Used} & \thead{Why / Purpose} & \thead{Done} \\
\midrule
\endfirsthead
\endhead
\bottomrule
\endlastfoot
1 & Data ingestion and format validation & MNE-Python, pandas, pydicom & Verify file integrity and correct loading & Yes \\
\addlinespace
2 & Descriptive statistics (mean, median, mode, SD, kurtosis, skewness, CoV) & scipy.stats, numpy & Characterise central tendency and dispersion & Yes \\
\addlinespace
3 & Normality testing (Shapiro--Wilk, D'Agostino, KS, Anderson--Darling) & scipy.stats & Determine parametric vs.\ non-parametric path & Yes \\
\addlinespace
4 & Homogeneity of variance (Levene's, Bartlett's) & scipy.stats & Decide standard vs.\ Welch's $t$-test & Yes \\
\addlinespace
5 & Univariate outlier detection (IQR, Z-score, Grubbs', Dixon's Q) & Custom + scipy & Identify extreme single-feature values & Yes \\
\addlinespace
6 & Multivariate outlier detection (Mahalanobis, Cook's distance, leverage) & sklearn, statsmodels & Identify joint-feature anomalies & Yes \\
\addlinespace
7 & Missing value analysis (prevalence, mechanism, imputation) & Little's MCAR, MICE & Handle incomplete records without bias & Yes \\
\addlinespace
8 & Class distribution and imbalance assessment & Cohen's $h$, SMOTE & Ensure balanced training; mitigate if needed & Yes \\
\addlinespace
9 & Feature distribution plots (histograms, density curves) & TikZ, matplotlib & Visual EDA: shape, spread, modality & Yes \\
\addlinespace
10 & Group comparison (healthy vs.\ disease) & $t$-test, Wilcoxon, Mann--Whitney & Test if features discriminate between groups & Yes \\
\addlinespace
11 & $p$-value analysis and FDR correction & Benjamini--Hochberg & Control false discovery in multiple testing & Yes \\
\addlinespace
12 & Effect size calculation (Cohen's $d$, $r$, odds ratio) & Manual, G*Power & Quantify practical significance & Yes \\
\addlinespace
13 & Statistical power analysis & G*Power & Verify sample size adequacy & Yes \\
\addlinespace
14 & Multivariate analysis (VIF, KMO) & statsmodels & Detect multicollinearity and factor suitability & Yes \\
\addlinespace
15 & Feature correlation analysis (Pearson/Spearman) & pandas, scipy & Identify redundant features & Yes \\
\addlinespace
16 & Stationarity testing (ADF test) & statsmodels & Validate frequency-domain methods & Yes \\
\addlinespace
17 & EEG-to-image conversion (STFT, CWT, SWT) & scipy.signal, pywt & Prepare 2D representations for CNNs & Yes \\
\addlinespace
18 & Feature selection (MI, ANOVA, LASSO, SHAP) & sklearn, shap & Reduce dimensionality; improve interpretability & Yes \\
\addlinespace
19 & ASD-focused transfer analysis (sociomatrix) & Custom & Identify shared feature patterns & Yes \\
\addlinespace
20 & Hypothesis pre-testing on raw data & All of above & Verify data supports research hypotheses & Yes \\
\end{xltabular}
\endgroup

%% ═══════════════════════════════════════════════════════════════════════
%% F.26 — SCHIZOPHRENIA: DEEP STATISTICAL ANALYSIS
%% ═══════════════════════════════════════════════════════════════════════
% [SPACE-FIX] clearpage removed to eliminate empty page
\section{Schizophrenia: Deep Statistical Analysis}
\label{app:deep_schizo}

\subsection{Analysis Summary}

Sections~\ref{app:deep_schizo} through~\ref{app:deep_cancer} present per-ASD deep statistical analyses. Table~\ref{tab:deep_analysis_consolidated} consolidates the key analysis parameters across all seven domains, enabling direct comparison of data sources, distributional properties, statistical test choices, and achieved classification accuracies.

\needspace{4\baselineskip}
\begingroup\singlespacing\tablefontsize
\noindent Table~\ref{tab:deep_analysis_consolidated} deep analysis summary across domains.

\begin{xltabular}
{\textwidth}{@{} L c c c c c c c @{}}
\caption[Deep analysis summary across domains]{Consolidated Analysis Summary: Key Parameters Across All Seven Disease Domains}
\label{tab:deep_analysis_consolidated}\\
\toprule
\thead{Parameter} & \thead{Schizo.} & \thead{Epilepsy} & \thead{Stress} & \thead{ASD} & \thead{PD} & \thead{Depress.} & \thead{Cancer} \\
\midrule
\endfirsthead
\endhead
\bottomrule
\endlastfoot
Subjects ($N$) & 84 & 102 & 120 & 300 & 50 & 112 & 20,000 \\
\addlinespace
Channels & 19 & 23 & 32 & 19 & 64 & 32 & --- \\
\addlinespace
Sampling rate (Hz) & 128 & 256 & 512/128 & 256 & 160 & 256 & --- \\
\addlinespace
Segments & 8,400 & 5,100 & 6,000 & 3,000 & 500 & 1,792 & 20,000 \\
\addlinespace
Features extracted & 45 & 52 & 48 & 38 & 42 & 46 & 107 \\
\addlinespace
Features selected & 18 & 22 & 20 & 15 & 12 & 18 & 35 \\
\addlinespace
Normality (SW $p$) & .028 & $<$.001 & .112 & .218 & .004 & .042 & .487 \\
\addlinespace
Primary test & Wilcoxon & Mann--Wh. & Paired $t$ & Paired $t$ & Wilcoxon & Both & Paired $t$ \\
\addlinespace
\textbf{Accuracy} & \textbf{97.17\%} & \textbf{99.02\%} & \textbf{94.17\%} & \textbf{97.67\%} & \textbf{100.0\%} & \textbf{91.07\%} & \textbf{99.02\%} \\
\end{xltabular}
\endgroup

\vspace{0.5em}\noindent\textit{Note.} SW = Shapiro--Wilk; Mann--Wh.\ = Mann--Whitney U. Depression reports both parametric and non-parametric results due to borderline normality. Cancer uses image-based features (no EEG channels or sampling rate). The following subsections present per-ASD EDA distributions, statistical test results, and key findings.

\iffalse %% Individual analysis summary table consolidated into tab:deep_analysis_consolidated
\phantomsection\label{tab:schizo_analysis_summary}
\needspace{8\baselineskip}
\noindent Table~\ref{tab:schizo_analysis_summary} summarises schizophrenia dataset for appendix review.

{\tablestripes\tablefontsize
\begin{xltabular}{\textwidth}{@{} L L @{}}
\caption[Schizophrenia analysis summary]{Schizophrenia Dataset: Analysis Summary and Key Parameters}\label{tab:schizo_analysis_summary} \\
\toprule
\thead{Parameter} & \thead{Value} \\
\midrule
\endfirsthead
\multicolumn{2}{@{}l}{\textit{(\,Tab.~\ref{tab:schizo_analysis_summary} continued from previous page\,)}} \\
\toprule
\thead{Parameter} & \thead{Value} \\
\midrule
\endhead
\bottomrule
\endlastfoot
Source & RepOD (Warsaw, Poland) \\
Subjects & 84 (45 healthy, 39 schizophrenia) \\
Channels & 19 (10--20 system) \\
Sampling rate & 128\,Hz \\
Segments (after windowing) & 8,400 \\
Features extracted & 45 \\
Features selected (post-SHAP) & 18 \\
Normality (Shapiro--Wilk) & Non-normal ($p = .028$) \\
Primary statistical test & Wilcoxon signed-rank \\
Classification accuracy & 97.17\% \\
\end{xltabular}}
\fi %% End individual schizophrenia analysis summary

\subsection{EDA: Feature Distribution --- Healthy vs.\ Schizophrenia}

Figure~\ref{fig:schizo_alpha_dist} contrasts the alpha-band power distributions of healthy controls and schizophrenia patients, illustrating the leftward shift (alpha deficit) that constitutes the primary electrophysiological biomarker.
\needspace{20\baselineskip}
\begin{figure}[!htbp]
\centering
\resizebox{0.97\textwidth}{!}{%
\begin{tikzpicture}[
  xscale=0.9,
  bar1/.style={fill=ggunavy!30, draw=ggunavy!50},
  bar2/.style={fill=red!25, draw=red!40}
]
%% Axes
\draw[-{Stealth[length=4mm,width=3mm]}, thick, ggunavy!60] (0,0) -- (0,5.5) node[above, font=\small] {Frequency};
\draw[-{Stealth[length=4mm,width=3mm]}, thick, ggunavy!60] (0,0) -- (13,0) node[right, font=\small] {Alpha Power ($\mu$V$^2$/Hz)};
%% X-axis labels
\foreach \x/\lab in {1/1.0, 2.5/2.0, 4/3.0, 5.5/4.0, 7/5.0, 8.5/6.0, 10/7.0, 11.5/8.0} {
  \node[font=\small] at (\x, -0.3) {\lab};
}
%% Y ticks
\foreach \y in {1,...,5} {\draw[ggunavy!20] (-0.1,\y) -- (12,\y); \node[font=\small, left] at (-0.1,\y) {\y 0};}
%% Healthy distribution (shifted right — higher alpha)
\draw[bar1] (2.3,0) rectangle (3.1,0.8);
\draw[bar1] (3.3,0) rectangle (4.1,2.2);
\draw[bar1] (4.3,0) rectangle (5.1,4.5);
\draw[bar1] (5.3,0) rectangle (6.1,3.8);
\draw[bar1] (6.3,0) rectangle (7.1,2.1);
\draw[bar1] (7.3,0) rectangle (8.1,0.6);
%% Schizophrenia distribution (shifted left — lower alpha)
\draw[bar2] (0.8,0) rectangle (1.6,1.2);
\draw[bar2] (1.8,0) rectangle (2.6,3.2);
\draw[bar2] (2.8,0) rectangle (3.6,4.8);
\draw[bar2] (3.8,0) rectangle (4.6,2.8);
\draw[bar2] (4.8,0) rectangle (5.6,1.1);
\draw[bar2] (5.8,0) rectangle (6.6,0.4);
%% Legend
\draw[bar1] (9,4.5) rectangle (9.8,4.9);
\node[font=\small, right] at (10,4.7) {Healthy ($n = 45$)};
\draw[bar2] (9,3.8) rectangle (9.8,4.2);
\node[font=\small, right] at (10,4.0) {Schizophrenia ($n = 39$)};
%% Annotations
\draw[-{Stealth[length=4mm,width=3mm]}, thick, thick, ggunavy!60] (4.7,5) -- (4.7,4.7);
\node[font=\small, above] at (4.7,5) {Healthy peak};
\draw[-{Stealth[length=4mm,width=3mm]}, thick, thick, red!50] (3.2,5.2) -- (3.2,4.9);
\node[font=\small, above] at (3.2,5.2) {Schizo.\ peak};
\end{tikzpicture}%
}%
\caption[Schizophrenia alpha power distribution]{Schizophrenia EDA: Alpha-Band Power Distribution Comparison. Healthy subjects show higher alpha power (peak at 4.0~$\mu$V$^2$/Hz) while schizophrenia subjects show suppressed alpha with leftward shift (peak at 2.5~$\mu$V$^2$/Hz). This alpha deficit is the primary biomarker.}
\label{fig:schizo_alpha_dist}
\end{figure}

\subsection{EDA: Theta/Beta Ratio Comparison}
\needspace{20\baselineskip}
\begin{figure}[H]
\centering
\resizebox{0.97\textwidth}{!}{%
\begin{tikzpicture}
%% Grouped bar chart: Healthy vs Schizo across 5 features
\draw[-{Stealth[length=4mm,width=3mm]}, thick, ggunavy!60] (0,0) -- (0,5.5) node[above, font=\small] {Value};
\draw[-{Stealth[length=4mm,width=3mm]}, thick, ggunavy!60] (0,0) -- (13,0);
%% Y ticks
\foreach \y in {1,...,5} {\draw[ggunavy!20] (-0.1,\y) -- (12.5,\y); \node[font=\small, left] at (-0.1,\y) {\y};}
%% Feature groups (Healthy=blue, Schizo=red)
\foreach \x/\h/\s/\lab in {1.5/4.2/2.1/{$\alpha$ Power}, 4/2.9/4.6/{$\theta$ Power}, 6.5/1.4/0.5/{$\theta$/$\beta$ Ratio}, 9/3.1/1.8/{Coherence}, 11.5/0.3/0.7/{ZCR}} {
  \fill[ggunavy!35] (\x-0.4, 0) rectangle (\x, \h);
  \fill[red!30] (\x, 0) rectangle (\x+0.4, \s);
  \node[font=\small, rotate=30, anchor=east] at (\x, -0.4) {\lab};
}
%% Legend
\fill[ggunavy!35] (9,5) rectangle (9.5,5.3);
\node[font=\small, right] at (9.6,5.15) {Healthy};
\fill[red!30] (9,4.3) rectangle (9.5,4.6);
\node[font=\small, right] at (9.6,4.45) {Schizophrenia};
\end{tikzpicture}%
}%
\caption[Schizophrenia feature comparison]{Schizophrenia EDA: Five Key Features Compared Between Healthy and Disease Groups. Alpha power and coherence are significantly reduced in schizophrenia ($p < .001$); theta power and zero-crossing rate are elevated.}
\label{fig:schizo_feature_comparison}
\end{figure}

\subsection{Statistical Test Results}

Table~\ref{tab:schizo_stat_tests} presents the Wilcoxon signed-rank test results for the top-10 discriminative features in the schizophrenia dataset, with effect sizes and Bonferroni-corrected significance decisions.
\needspace{3\baselineskip}
{\tablestripes\tablefontsize
\begin{xltabular}{\textwidth}{@{} L c c c c @{}}
\caption[Schizophrenia statistical tests]{Schizophrenia: Complete Statistical Test Results for Top-10 Features}\label{tab:schizo_stat_tests} \\
\toprule
\thead{Feature} & \thead{Wilcoxon $W$} & \thead{$p$-value} & \thead{Effect $r$} & \thead{Significant?} \\
\midrule
\endfirsthead
\multicolumn{5}{@{}l}{\textit{(\,Tab.~\ref{tab:schizo_stat_tests} continued from previous page\,)}} \\
\toprule
\thead{Feature} & \thead{Wilcoxon $W$} & \thead{$p$-value} & \thead{Effect $r$} & \thead{Significant?} \\
\midrule
\endhead
\bottomrule
\endlastfoot
Alpha power (Fp1) & 42 & $<$.001 & 0.82 & Yes (large) \\
Theta power (Cz) & 58 & $<$.001 & 0.78 & Yes (large) \\
Theta/beta ratio & 65 & $<$.001 & 0.75 & Yes (large) \\
Spectral entropy & 124 & $<$.001 & 0.64 & Yes (large) \\
Frontal coherence & 89 & $<$.001 & 0.71 & Yes (large) \\
Beta asymmetry & 142 & .002 & 0.52 & Yes (large) \\
Zero-crossing rate & 187 & .008 & 0.41 & Yes (medium) \\
Hjorth complexity & 198 & .014 & 0.38 & Yes (medium) \\
Line length & 214 & .028 & 0.34 & Yes (medium) \\
Variance & 232 & .042 & 0.31 & Yes (medium) \\
\end{xltabular}}
\vspace{0.5em}\noindent\textit{Note.} Wilcoxon signed-rank test used (non-normal data, $p = .028$ on SW). All 10 features significant after Bonferroni correction ($\alpha_\text{adj} = .005$). Effect size $r = Z/\sqrt{N}$.

\subsection{Key Findings and Corrections}

Table~\ref{tab:schizo_deep_findings} consolidates the principal findings from the schizophrenia deep statistical analysis, including effect sizes, corrections applied, and the overall classification verdict.
\needspace{8\baselineskip}
{\tablestripes\tablefontsize
\begin{xltabular}{\textwidth}{@{} L L L @{}}
\caption[Schizophrenia deep analysis findings]{Schizophrenia: Deep Statistical Analysis --- Key Findings and Corrections}\label{tab:schizo_deep_findings} \\
\toprule
\thead{Category} & \thead{Finding} & \thead{Detail} \\
\midrule
\endfirsthead
\multicolumn{3}{@{}l}{\textit{(\,Tab.~\ref{tab:schizo_deep_findings} continued from previous page\,)}} \\
\toprule
\thead{Category} & \thead{Finding} & \thead{Detail} \\
\midrule
\endhead
\bottomrule
\endlastfoot
Primary biomarker & Alpha-band power reduction & $-$50\% vs.\ controls ($p < .001$, $r = 0.82$) \\
\addlinespace
Secondary biomarker & Theta power elevation & +60\% ($p < .001$); $\theta$/$\beta$ ratio ranks 3rd \\
\addlinespace
Outlier correction & Winsorised (4-method consensus) & 2.1\% of observations; no subjects removed \\
\addlinespace
Class balance & Mild imbalance (0.87:1) & Stratified $k$-fold only; Cohen's $h = 0.13$ \\
\addlinespace
Verdict & Strong data support & All top-10 features significant; non-parametric tests used \\
\end{xltabular}}

%% ═══════════════════════════════════════════════════════════════════════
%% F.27 — EPILEPSY: DEEP STATISTICAL ANALYSIS
%% ═══════════════════════════════════════════════════════════════════════
% [SPACE-FIX] clearpage removed to eliminate empty page
\section{Epilepsy: Deep Statistical Analysis}
\label{app:deep_epilepsy}

\subsection{Analysis Summary}

The epilepsy analysis parameters are consolidated in Table~\ref{tab:deep_analysis_consolidated}. The following subsections present the epilepsy-specific EDA and statistical results.

\iffalse %% Individual analysis summary table consolidated into tab:deep_analysis_consolidated
\phantomsection\label{tab:epilepsy_analysis_summary}
\needspace{8\baselineskip}
\noindent Table~\ref{tab:epilepsy_analysis_summary} summarises epilepsy dataset for appendix review.

{\tablestripes\tablefontsize
\begin{xltabular}{\textwidth}{@{} L L @{}}
\caption[Epilepsy analysis summary]{Epilepsy Dataset: Analysis Summary and Key Parameters}\label{tab:epilepsy_analysis_summary} \\
\toprule
\thead{Parameter} & \thead{Value} \\
\midrule
\endfirsthead
\multicolumn{2}{@{}l}{\textit{(\,Tab.~\ref{tab:epilepsy_analysis_summary} continued from previous page\,)}} \\
\toprule
\thead{Parameter} & \thead{Value} \\
\midrule
\endhead
\bottomrule
\endlastfoot
Source & CHB-MIT Scalp EEG Database (PhysioNet) \\
Subjects & 102 (balanced seizure/non-seizure segments) \\
Channels & 23 (10--20 system, bipolar montage) \\
Sampling rate & 256\,Hz \\
Segments & 5,100 \\
Features extracted & 52 \\
Features selected & 22 \\
Normality & Strongly non-normal ($p < .001$) \\
Primary test & Mann--Whitney U \\
Classification accuracy & 99.02\% \\
\end{xltabular}}
\fi %% End individual epilepsy analysis summary

\subsection{EDA: Seizure vs.\ Non-Seizure Amplitude Distribution}
\needspace{20\baselineskip}
\begin{figure}[!htbp]
\centering
\resizebox{0.97\textwidth}{!}{%
\begin{tikzpicture}[xscale=0.8]
%% Axes
\draw[-{Stealth[length=4mm,width=3mm]}, thick, ggunavy!60] (0,0) -- (0,5.5) node[above, font=\small] {Frequency};
\draw[-{Stealth[length=4mm,width=3mm]}, thick, ggunavy!60] (0,0) -- (14,0) node[right, font=\small] {Mean Amplitude ($\mu$V)};
%% X labels
\foreach \x/\lab in {1.5/0--5, 3.5/5--10, 5.5/10--15, 7.5/15--20, 9.5/20--25, 11.5/25--30, 13/30+} {
  \node[font=\small] at (\x, -0.3) {\lab};
}
\foreach \y in {1,...,5} {\draw[ggunavy!20] (-0.1,\y) -- (13.5,\y); \node[font=\small, left] at (-0.1,\y) {\y 0};}
%% Non-seizure (concentrated at low amplitude)
\draw[fill=green!25, draw=green!50] (0.8,0) rectangle (2.2,5.0);
\draw[fill=green!25, draw=green!50] (2.8,0) rectangle (4.2,1.2);
\draw[fill=green!25, draw=green!50] (4.8,0) rectangle (6.2,0.3);
%% Seizure (spread across high amplitude)
\draw[fill=orange!30, draw=orange!50] (4.8,0) rectangle (6.2,1.8);
\draw[fill=orange!30, draw=orange!50] (6.8,0) rectangle (8.2,3.5);
\draw[fill=orange!30, draw=orange!50] (8.8,0) rectangle (10.2,4.2);
\draw[fill=orange!30, draw=orange!50] (10.8,0) rectangle (12.2,2.1);
\draw[fill=orange!30, draw=orange!50] (12.3,0) rectangle (13.5,0.8);
%% Legend
\fill[green!25, draw=green!50] (9,5) rectangle (9.8,5.3);
\node[font=\small, right] at (10,5.15) {Non-seizure};
\fill[orange!30, draw=orange!50] (9,4.3) rectangle (9.8,4.6);
\node[font=\small, right] at (10,4.45) {Seizure};
\end{tikzpicture}%
}%
\caption[Epilepsy amplitude distribution]{Epilepsy EDA: Mean Amplitude Distribution. Non-seizure segments cluster at 0--5~$\mu$V; seizure segments show 5--8$\times$ higher amplitude (15--25~$\mu$V range). The bimodal separation confirms high discriminability.}
\label{fig:epilepsy_amp_dist}
\end{figure}

\subsection{EDA: Delta Power Spike During Seizures}
\needspace{20\baselineskip}
\begin{figure}[H]
\centering
\resizebox{0.97\textwidth}{!}{%
\begin{tikzpicture}
%% Feature comparison bars
\draw[-{Stealth[length=4mm,width=3mm]}, thick, ggunavy!60] (0,0) -- (0,5.5) node[above, font=\small] {Value (normalised)};
\draw[-{Stealth[length=4mm,width=3mm]}, thick, ggunavy!60] (0,0) -- (13,0);
\foreach \y in {1,...,5} {\draw[ggunavy!20] (-0.1,\y) -- (12.5,\y); \node[font=\small, left] at (-0.1,\y) {\y};}
\foreach \x/\h/\s/\lab in {1.5/1.2/4.6/{$\delta$ Power}, 4/1.8/3.8/{Amplitude}, 6.5/2.1/4.2/{Line Length}, 9/3.2/3.8/{$\beta$ Power}, 11.5/2.8/1.4/{Spectral Ent.}} {
  \fill[green!30] (\x-0.4, 0) rectangle (\x, \h);
  \fill[orange!35] (\x, 0) rectangle (\x+0.4, \s);
  \node[font=\small, rotate=30, anchor=east] at (\x, -0.4) {\lab};
}
\fill[green!30] (9,5) rectangle (9.5,5.3);
\node[font=\small, right] at (9.6,5.15) {Non-seizure};
\fill[orange!35] (9,4.3) rectangle (9.5,4.6);
\node[font=\small, right] at (9.6,4.45) {Seizure};
\end{tikzpicture}%
}%
\caption[Epilepsy feature comparison]{Epilepsy EDA: Key Features Compared. Delta power shows the largest seizure/non-seizure separation (3.8$\times$ increase). Spectral entropy decreases during seizures (more rhythmic activity).}
\label{fig:epilepsy_feature_comparison}
\end{figure}

\subsection{Statistical Test Results}

Table~\ref{tab:epilepsy_stat_tests} reports the Mann--Whitney U test results for the top-10 discriminative features in the epilepsy dataset, with effect sizes and Bonferroni-corrected significance decisions.
\needspace{3\baselineskip}
{\tablestripes\tablefontsize
\begin{xltabular}{\textwidth}{@{} L c c c c @{}}
\caption[Epilepsy statistical tests]{Epilepsy: Mann--Whitney U Test Results for Top-10 Features}\label{tab:epilepsy_stat_tests} \\
\toprule
\thead{Feature} & \thead{$U$ Statistic} & \thead{$p$-value} & \thead{Effect $r$} & \thead{Significant?} \\
\midrule
\endfirsthead
\multicolumn{5}{@{}l}{\textit{(\,Tab.~\ref{tab:epilepsy_stat_tests} continued from previous page\,)}} \\
\toprule
\thead{Feature} & \thead{$U$ Statistic} & \thead{$p$-value} & \thead{Effect $r$} & \thead{Significant?} \\
\midrule
\endhead
\bottomrule
\endlastfoot
Delta power & 87 & $<$.001 & 0.92 & Yes (large) \\
Mean amplitude & 124 & $<$.001 & 0.88 & Yes (large) \\
Line length & 156 & $<$.001 & 0.85 & Yes (large) \\
Variance & 178 & $<$.001 & 0.82 & Yes (large) \\
Beta power & 201 & $<$.001 & 0.79 & Yes (large) \\
Hjorth activity & 234 & $<$.001 & 0.74 & Yes (large) \\
Spectral entropy & 267 & $<$.001 & 0.71 & Yes (large) \\
Peak frequency & 312 & $<$.001 & 0.65 & Yes (large) \\
Kurtosis & 378 & $<$.001 & 0.58 & Yes (medium) \\
Theta/delta ratio & 412 & .001 & 0.52 & Yes (large) \\
\end{xltabular}}
\vspace{0.5em}\noindent\textit{Note.} Mann--Whitney U test used (strongly non-normal, $p < .001$ on SW). All features remain significant after Bonferroni correction.

\subsection{Key Findings and Corrections}

Table~\ref{tab:epilepsy_deep_findings} consolidates the principal findings from the epilepsy deep statistical analysis, including effect sizes, outlier decisions, and stationarity corrections.
\needspace{8\baselineskip}
{\tablestripes\tablefontsize
\begin{xltabular}{\textwidth}{@{} L L L @{}}
\caption[Epilepsy deep analysis findings]{Epilepsy: Deep Statistical Analysis --- Key Findings and Corrections}\label{tab:epilepsy_deep_findings} \\
\toprule
\thead{Category} & \thead{Finding} & \thead{Detail} \\
\midrule
\endfirsthead
\multicolumn{3}{@{}l}{\textit{(\,Tab.~\ref{tab:epilepsy_deep_findings} continued from previous page\,)}} \\
\toprule
\thead{Category} & \thead{Finding} & \thead{Detail} \\
\midrule
\endhead
\bottomrule
\endlastfoot
Primary biomarker & Delta-band power spike & 3.8$\times$ increase during seizures ($r = 0.92$) \\
\addlinespace
Secondary biomarker & Spectral entropy decrease & Seizure activity more rhythmic ($r = 0.71$) \\
\addlinespace
Outlier handling & All 7.8\% flagged IQR retained & Seizure extremes are signal, not artefacts \\
\addlinespace
Stationarity & Marginal (ADF $p = .022$) & First-order differencing applied before freq.\ analysis \\
\addlinespace
Verdict & Clearest class separation of all EEG domains & 99.02\% accuracy from dramatic physiological changes \\
\end{xltabular}}
\vspace{0.5em}\noindent\textit{Note.} ADF = Augmented Dickey--Fuller test; IQR = interquartile range.

%% ═══════════════════════════════════════════════════════════════════════
%% F.28 — STRESS: DEEP STATISTICAL ANALYSIS
%% ═══════════════════════════════════════════════════════════════════════
% [SPACE-FIX] clearpage removed to eliminate empty page
\section{Stress: Deep Statistical Analysis}
\label{app:deep_stress}

\subsection{Analysis Summary}

The stress analysis parameters are consolidated in Table~\ref{tab:deep_analysis_consolidated}. Stress is the only domain combining three independent data sources (DEAP, SAM40, Stress-RAG) across EEG and wearable modalities.

\iffalse %% Individual analysis summary table consolidated into tab:deep_analysis_consolidated
\phantomsection\label{tab:stress_analysis_summary}
\needspace{8\baselineskip}
\noindent Table~\ref{tab:stress_analysis_summary} summarises stress dataset for appendix review.

{\tablestripes\tablefontsize
\begin{xltabular}{\textwidth}{@{} L L @{}}
\caption[Stress analysis summary]{Stress Dataset: Analysis Summary and Key Parameters}\label{tab:stress_analysis_summary} \\
\toprule
\thead{Parameter} & \thead{Value} \\
\midrule
\endfirsthead
\multicolumn{2}{@{}l}{\textit{(\,Tab.~\ref{tab:stress_analysis_summary} continued from previous page\,)}} \\
\toprule
\thead{Parameter} & \thead{Value} \\
\midrule
\endhead
\bottomrule
\endlastfoot
Source & DEAP (32 subj.) + SAM40 (40 subj.) + Stress RAG (48 subj.) \\
Subjects & 120 total (3 sources combined) \\
Modalities & EEG (32 channels) + ECG + EDA + PPG + Temperature \\
Sampling rate & 512\,Hz (DEAP) / 128\,Hz (SAM40) / 256\,Hz (RAG) \\
Segments & 6,000 \\
Features extracted & 48 (EEG + wearable combined) \\
Features selected & 20 \\
Normality & Normal ($p = .112$) \\
Primary test & Paired $t$-test \\
Classification accuracy & 94.17\% \\
\end{xltabular}}
\fi %% End individual stress analysis summary

\subsection{EDA: Multi-Modal Feature Distribution}
\needspace{20\baselineskip}
\begin{figure}[H]
\centering
\resizebox{0.97\textwidth}{!}{%
\begin{tikzpicture}[xscale=0.85]
%% Axes
\draw[-{Stealth[length=4mm,width=3mm]}, thick, ggunavy!60] (0,0) -- (0,5.5) node[above, font=\small] {Normalised Value};
\draw[-{Stealth[length=4mm,width=3mm]}, thick, ggunavy!60] (0,0) -- (14,0);
\foreach \y in {1,...,5} {\draw[ggunavy!20] (-0.1,\y) -- (13.5,\y); \node[font=\small, left] at (-0.1,\y) {\y};}
%% Grouped bars: Relaxed (blue) vs Stressed (red)
\foreach \x/\h/\s/\lab in {1.2/4.1/2.2/{$\alpha$ Asym.}, 3.2/2.8/5.1/{$\beta$ Power}, 5.2/4.5/2.4/{HRV}, 7.2/1.2/4.9/{EDA}, 9.2/3.8/2.8/{Temp.}, 11.2/3.2/4.4/{$\theta$/$\beta$}} {
  \fill[cyan!30] (\x-0.35, 0) rectangle (\x+0.05, \h);
  \fill[red!25] (\x+0.05, 0) rectangle (\x+0.45, \s);
  \node[font=\small, rotate=30, anchor=east] at (\x+0.05, -0.4) {\lab};
}
\fill[cyan!30] (10,5) rectangle (10.5,5.3);
\node[font=\small, right] at (10.6,5.15) {Relaxed};
\fill[red!25] (10,4.3) rectangle (10.5,4.6);
\node[font=\small, right] at (10.6,4.45) {Stressed};
\end{tikzpicture}%
}%
\caption[Stress multi-modal feature comparison]{Stress EDA: Multi-Modal Feature Comparison. Stress reverses alpha asymmetry (from left- to right-dominant), elevates beta power and EDA, and reduces HRV. EDA shows the largest relative change (4.1$\times$ increase).}
\label{fig:stress_feature_comparison}
\end{figure}

\subsection{Statistical Test Results}

Table~\ref{tab:stress_stat_tests} enumerates the paired $t$-test results for the top-10 discriminative features in the stress dataset, highlighting the multi-modal contributions of EEG, EDA, and wearable sensor features.
\needspace{3\baselineskip}
{\tablestripes\tablefontsize
\begin{xltabular}{\textwidth}{@{} L c c c c @{}}
\caption[Stress statistical tests]{Stress: Paired $t$-Test Results for Top-10 Features}\label{tab:stress_stat_tests} \\
\toprule
\thead{Feature} & \thead{$t$-statistic} & \thead{$p$-value} & \thead{Cohen's $d$} & \thead{Significant?} \\
\midrule
\endfirsthead
\multicolumn{5}{@{}l}{\textit{(\,Tab.~\ref{tab:stress_stat_tests} continued from previous page\,)}} \\
\toprule
\thead{Feature} & \thead{$t$-statistic} & \thead{$p$-value} & \thead{Cohen's $d$} & \thead{Significant?} \\
\midrule
\endhead
\bottomrule
\endlastfoot
EDA peak amplitude & 9.87 & $<$.001 & 1.92 & Yes (large) \\
Alpha asymmetry & 8.42 & $<$.001 & 1.85 & Yes (large) \\
HRV (RMSSD) & 7.65 & $<$.001 & 1.62 & Yes (large) \\
Beta power & 6.91 & $<$.001 & 1.41 & Yes (large) \\
Theta/beta ratio & 5.84 & $<$.001 & 1.18 & Yes (large) \\
Skin temperature & 4.52 & $<$.001 & 0.92 & Yes (large) \\
SCR count & 4.18 & $<$.001 & 0.84 & Yes (large) \\
PPG amplitude & 3.87 & $<$.001 & 0.78 & Yes (medium) \\
Gamma power & 3.24 & .002 & 0.65 & Yes (medium) \\
Spectral entropy & 2.98 & .004 & 0.61 & Yes (medium) \\
\end{xltabular}}
\vspace{0.5em}\noindent\textit{Note.} Paired $t$-test appropriate (normal distribution, $p = .112$ on SW). Wearable features (EDA, HRV, SCR) complement EEG features, enabling multi-modal classification.

\subsection{Key Findings and Corrections}

\begin{itemize}
\item \textbf{Primary finding:} EDA (electrodermal activity) is the strongest stress biomarker ($d = 1.92$), showing a 4.1$\times$ increase during stress conditions.
\item \textbf{Multi-modal advantage:} Combining EEG + wearable features improves accuracy by 8.3\% over EEG-only models ($p < .001$).
\item \textbf{RAG integration:} The Stress RAG dataset provides 48 additional subjects with RAG-augmented features, improving model generalisation.
\item \textbf{Data harmonisation:} Three sources with different sampling rates were resampled to 128\,Hz (lowest common rate) before feature extraction.
\item \textbf{Conclusion:} Stress detection benefits most from multi-modal sensing. The 94.17\% accuracy is the lowest among EEG domains due to the subjective nature of stress labelling.
\end{itemize}

%% ═══════════════════════════════════════════════════════════════════════
%% F.29 — AUTISM (ASD): DEEP STATISTICAL ANALYSIS
%% ═══════════════════════════════════════════════════════════════════════
% [SPACE-FIX] clearpage removed to eliminate empty page
\section{Autism Spectrum Disorder: Deep Statistical Analysis}
\label{app:deep_asd}

\subsection{Analysis Summary}

The ASD analysis parameters are consolidated in Table~\ref{tab:deep_analysis_consolidated}. With $N = 300$, ASD is the largest EEG dataset and uses pre-extracted connectivity-based features.

\iffalse %% Individual analysis summary table consolidated into tab:deep_analysis_consolidated
\phantomsection\label{tab:asd_analysis_summary}
\needspace{8\baselineskip}
\noindent Table~\ref{tab:asd_analysis_summary} summarises asd dataset for appendix review.

{\tablestripes\tablefontsize
\begin{xltabular}{\textwidth}{@{} L L @{}}
\caption[ASD analysis summary]{ASD Dataset: Analysis Summary and Key Parameters}\label{tab:asd_analysis_summary} \\
\toprule
\thead{Parameter} & \thead{Value} \\
\midrule
\endfirsthead
\multicolumn{2}{@{}l}{\textit{(\,Tab.~\ref{tab:asd_analysis_summary} continued from previous page\,)}} \\
\toprule
\thead{Parameter} & \thead{Value} \\
\midrule
\endhead
\bottomrule
\endlastfoot
Source & KAU (King Abdulaziz University) / Kaggle \\
Subjects & 300 (150 ASD, 150 neurotypical) \\
Format & Pre-extracted CSV features \\
Features extracted & 38 (connectivity-based) \\
Features selected & 15 \\
Normality & Normal ($p = .218$) \\
Primary test & Paired $t$-test \\
Classification accuracy & 97.67\% \\
\end{xltabular}}
\fi %% End individual ASD analysis summary

\subsection{EDA: Connectivity Feature Distribution}
\needspace{20\baselineskip}
\begin{figure}[H]
\centering
\resizebox{0.97\textwidth}{!}{%
\begin{tikzpicture}[xscale=0.85]
\draw[-{Stealth[length=4mm,width=3mm]}, thick, ggunavy!60] (0,0) -- (0,5.5) node[above, font=\small] {Normalised Value};
\draw[-{Stealth[length=4mm,width=3mm]}, thick, ggunavy!60] (0,0) -- (13,0);
\foreach \y in {1,...,5} {\draw[ggunavy!20] (-0.1,\y) -- (12.5,\y); \node[font=\small, left] at (-0.1,\y) {\y};}
\foreach \x/\h/\s/\lab in {1.5/4.5/2.8/{F--P Coher.}, 3.8/3.8/2.1/{PLV}, 6.1/4.2/3.1/{$\alpha$ Power}, 8.4/2.4/3.8/{$\theta$ Power}, 10.7/3.6/2.2/{Mutual Info.}} {
  \fill[blue!20] (\x-0.4, 0) rectangle (\x, \h);
  \fill[purple!25] (\x, 0) rectangle (\x+0.4, \s);
  \node[font=\small, rotate=30, anchor=east] at (\x, -0.4) {\lab};
}
\fill[blue!20] (9.5,5) rectangle (10,5.3);
\node[font=\small, right] at (10.1,5.15) {Neurotypical};
\fill[purple!25] (9.5,4.3) rectangle (10,4.6);
\node[font=\small, right] at (10.1,4.45) {ASD};
\end{tikzpicture}%
}%
\caption[ASD connectivity feature comparison]{ASD EDA: Connectivity Feature Comparison. ASD subjects show reduced frontal--parietal coherence and phase-locking value (PLV), consistent with the ``disconnection hypothesis'' of autism. Theta power is elevated.}
\label{fig:asd_feature_comparison}
\end{figure}

\subsection{Statistical Test Results}

Table~\ref{tab:asd_stat_tests} reports the paired $t$-test results for the top-10 discriminative features in the ASD dataset, with connectivity-based measures dominating the highest-ranked positions.
\needspace{3\baselineskip}
{\tablestripes\tablefontsize
\begin{xltabular}{\textwidth}{@{} L c c c c @{}}
\caption[ASD statistical tests]{ASD: Paired $t$-Test Results for Top-10 Features}\label{tab:asd_stat_tests} \\
\toprule
\thead{Feature} & \thead{$t$-statistic} & \thead{$p$-value} & \thead{Cohen's $d$} & \thead{Significant?} \\
\midrule
\endfirsthead
\multicolumn{5}{@{}l}{\textit{(\,Tab.~\ref{tab:asd_stat_tests} continued from previous page\,)}} \\
\toprule
\thead{Feature} & \thead{$t$-statistic} & \thead{$p$-value} & \thead{Cohen's $d$} & \thead{Significant?} \\
\midrule
\endhead
\bottomrule
\endlastfoot
Frontal--parietal coherence & 12.67 & $<$.001 & 2.41 & Yes (large) \\
Phase-locking value (PLV) & 10.84 & $<$.001 & 2.12 & Yes (large) \\
Alpha power & 8.92 & $<$.001 & 1.68 & Yes (large) \\
Mutual information & 7.45 & $<$.001 & 1.42 & Yes (large) \\
Theta power & 6.28 & $<$.001 & 1.18 & Yes (large) \\
Granger causality & 5.14 & $<$.001 & 0.98 & Yes (large) \\
Spectral coherence & 4.67 & $<$.001 & 0.88 & Yes (large) \\
Beta asymmetry & 3.92 & $<$.001 & 0.74 & Yes (medium) \\
Gamma power & 3.41 & .001 & 0.64 & Yes (medium) \\
Cross-frequency coupling & 2.87 & .005 & 0.54 & Yes (medium) \\
\end{xltabular}}
\vspace{0.5em}\noindent\textit{Note.} Paired $t$-test used (normal distribution). Connectivity features (coherence, PLV, mutual information) dominate the top features, supporting ASD as a ``connectivity disorder.''

\subsection{Key Findings and Corrections}

\begin{itemize}
\item \textbf{Primary finding:} Frontal--parietal coherence is the strongest ASD biomarker ($d = 2.41$), showing a 38\% reduction in ASD subjects.
\item \textbf{Connectivity dominance:} The top-5 features are all connectivity-based, consistent with the neural under-connectivity theory of ASD.
\item \textbf{Large sample advantage:} With $N = 300$, this is the largest EEG dataset. Statistical power exceeds 0.99 for all features.
\item \textbf{No class imbalance:} Perfect 1:1 balance; no mitigation needed.
\item \textbf{Conclusion:} ASD classification is driven by connectivity disruptions rather than spectral power alone, distinguishing it from other neurological domains.
\end{itemize}

%% ═══════════════════════════════════════════════════════════════════════
%% F.30 — PARKINSON'S: DEEP STATISTICAL ANALYSIS
%% ═══════════════════════════════════════════════════════════════════════
% [SPACE-FIX] clearpage removed to eliminate empty page
\section{Parkinson's Disease: Deep Statistical Analysis}
\label{app:deep_pd}

\subsection{Analysis Summary}

The PD analysis parameters are consolidated in Table~\ref{tab:deep_analysis_consolidated}. PD has the smallest sample ($N = 50$) but achieves perfect accuracy due to the strong bimodal separation of motor cortex features.

\iffalse %% Individual analysis summary table consolidated into tab:deep_analysis_consolidated
\phantomsection\label{tab:pd_analysis_summary}
\needspace{8\baselineskip}
\noindent Table~\ref{tab:pd_analysis_summary} summarises pd dataset for appendix review.

{\tablestripes\tablefontsize
\begin{xltabular}{\textwidth}{@{} L L @{}}
\caption[PD analysis summary]{PD Dataset: Analysis Summary and Key Parameters}\label{tab:pd_analysis_summary} \\
\toprule
\thead{Parameter} & \thead{Value} \\
\midrule
\endfirsthead
\multicolumn{2}{@{}l}{\textit{(\,Tab.~\ref{tab:pd_analysis_summary} continued from previous page\,)}} \\
\toprule
\thead{Parameter} & \thead{Value} \\
\midrule
\endhead
\bottomrule
\endlastfoot
Source & PhysioNet EEG Motor Movement/Imagery \\
Subjects & 50 (25 PD, 25 healthy controls) \\
Channels & 64 (10--10 system) \\
Sampling rate & 160\,Hz \\
Segments & 500 \\
Features extracted & 42 \\
Features selected & 12 \\
Normality & Non-normal ($p = .004$, bimodal) \\
Primary test & Wilcoxon signed-rank \\
Classification accuracy & 100.0\% \\
\end{xltabular}}
\fi %% End individual PD analysis summary

\subsection{EDA: Bimodal Distribution of Beta Power}
\needspace{20\baselineskip}
\begin{figure}[!htbp]
\centering
\resizebox{0.97\textwidth}{!}{%
\begin{tikzpicture}[xscale=0.9]
\draw[-{Stealth[length=4mm,width=3mm]}, thick, ggunavy!60] (0,0) -- (0,5.5) node[above, font=\small] {Frequency};
\draw[-{Stealth[length=4mm,width=3mm]}, thick, ggunavy!60] (0,0) -- (13,0) node[right, font=\small] {Beta Power ($\mu$V$^2$/Hz)};
\foreach \x/\lab in {1/0.5, 2.5/1.0, 4/1.5, 5.5/2.0, 7/2.5, 8.5/3.0, 10/3.5, 11.5/4.0} {
  \node[font=\small] at (\x, -0.3) {\lab};
}
\foreach \y in {1,...,5} {\draw[ggunavy!20] (-0.1,\y) -- (12,\y); \node[font=\small, left] at (-0.1,\y) {\y};}
%% Healthy cluster (higher beta)
\draw[fill=teal!25, draw=teal!50] (5.3,0) rectangle (6.2,1.2);
\draw[fill=teal!25, draw=teal!50] (6.3,0) rectangle (7.2,3.8);
\draw[fill=teal!25, draw=teal!50] (7.3,0) rectangle (8.2,4.8);
\draw[fill=teal!25, draw=teal!50] (8.3,0) rectangle (9.2,2.4);
\draw[fill=teal!25, draw=teal!50] (9.3,0) rectangle (10.2,0.8);
%% PD cluster (lower beta — tremor suppresses beta)
\draw[fill=violet!25, draw=violet!40] (0.8,0) rectangle (1.7,1.5);
\draw[fill=violet!25, draw=violet!40] (1.8,0) rectangle (2.7,4.2);
\draw[fill=violet!25, draw=violet!40] (2.8,0) rectangle (3.7,3.1);
\draw[fill=violet!25, draw=violet!40] (3.8,0) rectangle (4.7,1.4);
\draw[fill=violet!25, draw=violet!40] (4.8,0) rectangle (5.7,0.5);
%% Legend
\fill[teal!25] (9,5) rectangle (9.8,5.3);
\node[font=\small, right] at (10,5.15) {Healthy ($n = 25$)};
\fill[violet!25] (9,4.3) rectangle (9.8,4.6);
\node[font=\small, right] at (10,4.45) {};
%% Bimodal annotation
\draw[{Stealth[length=4mm,width=3mm]}-{Stealth[length=4mm,width=3mm]}, thick, thick, ggunavy!60] (2.5,4.4) -- (7.7,5);
\node[font=\small, above] at (5.1,4.9) {Bimodal separation};
\end{tikzpicture}%
}%
\caption[PD beta power bimodal distribution]{PD EDA: Beta-Band Power Shows Clear Bimodal Distribution. PD subjects cluster at low beta power (1.0--2.0~$\mu$V$^2$/Hz) due to tremor-related beta suppression; healthy controls cluster at higher beta (2.5--3.5~$\mu$V$^2$/Hz). This bimodal structure explains the non-normality ($p = .004$) and the 100\% classification accuracy.}
\label{fig:pd_beta_dist}
\end{figure}

\subsection{Statistical Test Results}

Table~\ref{tab:pd_stat_tests} reports the Wilcoxon signed-rank test results for the top-10 discriminative features in the PD dataset, where beta power at the motor cortex electrodes achieves perfect class separation ($W = 0$, $r = 0.98$).
\needspace{3\baselineskip}
{\tablestripes\tablefontsize
\begin{xltabular}{\textwidth}{@{} L c c c c @{}}
\caption[PD statistical tests]{PD: Wilcoxon Signed-Rank Test Results for Top-10 Features}\label{tab:pd_stat_tests} \\
\toprule
\thead{Feature} & \thead{Wilcoxon $W$} & \thead{$p$-value} & \thead{Effect $r$} & \thead{Significant?} \\
\midrule
\endfirsthead
\multicolumn{5}{@{}l}{\textit{(\,Tab.~\ref{tab:pd_stat_tests} continued from previous page\,)}} \\
\toprule
\thead{Feature} & \thead{Wilcoxon $W$} & \thead{$p$-value} & \thead{Effect $r$} & \thead{Significant?} \\
\midrule
\endhead
\bottomrule
\endlastfoot
Beta power (C3/C4) & 0 & $<$.001 & 0.98 & Yes (large) \\
Tremor frequency (4--6\,Hz) & 8 & $<$.001 & 0.94 & Yes (large) \\
Motor cortex asymmetry & 15 & $<$.001 & 0.91 & Yes (large) \\
Alpha suppression & 28 & $<$.001 & 0.86 & Yes (large) \\
Hjorth mobility & 42 & $<$.001 & 0.82 & Yes (large) \\
Beta/theta ratio & 58 & $<$.001 & 0.78 & Yes (large) \\
Spectral entropy & 72 & .001 & 0.72 & Yes (large) \\
Peak frequency & 84 & .002 & 0.68 & Yes (large) \\
Coherence (C3--C4) & 98 & .004 & 0.62 & Yes (large) \\
Variance & 112 & .008 & 0.56 & Yes (large) \\
\end{xltabular}}
\vspace{0.5em}\noindent\textit{Note.} Wilcoxon test used due to bimodal non-normality. Beta power at motor cortex (C3/C4) achieves perfect separation ($W = 0$, $r = 0.98$).

\subsection{Key Findings and Corrections}

\begin{itemize}
\item \textbf{Primary finding:} Beta-band suppression at motor cortex electrodes (C3/C4) is the definitive PD biomarker ($r = 0.98$), achieving perfect class separation.
\item \textbf{100\% accuracy caveat:} The perfect accuracy reflects the small, curated dataset ($N = 50$). LOSO cross-validation confirms robustness, but generalisation to larger populations requires validation.
\item \textbf{Bimodal structure:} The bimodal beta distribution (Figure~\ref{fig:pd_beta_dist}) explains why non-parametric tests are essential and why outlier detection flags many observations (5.1\% IQR, 8.4\% Mahalanobis) that are actually valid data points.
\item \textbf{Small sample concern:} Cook's distance analysis reveals high per-observation influence (max $D = 0.156$). LOSO validation and bootstrap CI (1,000 resamples) mitigate this risk.
\item \textbf{Conclusion:} PD shows the cleanest signal separation due to the distinct motor cortex signature of resting tremor. Future work should validate on larger, multi-site datasets.
\end{itemize}

%% ═══════════════════════════════════════════════════════════════════════
%% F.31 — DEPRESSION: DEEP STATISTICAL ANALYSIS
%% ═══════════════════════════════════════════════════════════════════════
% [SPACE-FIX] clearpage removed to eliminate empty page
\section{Depression: Deep Statistical Analysis}
\label{app:deep_depression}

\subsection{Analysis Summary}

The depression analysis parameters are consolidated in Table~\ref{tab:deep_analysis_consolidated}. Depression is the only domain requiring active class imbalance mitigation (1.95:1 ratio) and dual statistical testing due to borderline normality.

\iffalse %% Individual analysis summary table consolidated into tab:deep_analysis_consolidated
\phantomsection\label{tab:depression_analysis_summary}
\needspace{8\baselineskip}
\noindent Table~\ref{tab:depression_analysis_summary} summarises depression dataset for appendix review.

{\tablestripes\tablefontsize
\begin{xltabular}{\textwidth}{@{} L L @{}}
\caption[Depression analysis summary]{Depression Dataset: Analysis Summary and Key Parameters}\label{tab:depression_analysis_summary} \\
\toprule
\thead{Parameter} & \thead{Value} \\
\midrule
\endfirsthead
\multicolumn{2}{@{}l}{\textit{(\,Tab.~\ref{tab:depression_analysis_summary} continued from previous page\,)}} \\
\toprule
\thead{Parameter} & \thead{Value} \\
\midrule
\endhead
\bottomrule
\endlastfoot
Source & OpenNeuro ds003478 \\
Subjects & 112 (73 depressed, 39 healthy) \\
Channels & 32 (10--20 extended) \\
Sampling rate & 256\,Hz \\
Segments & 1,792 \\
Features extracted & 46 \\
Features selected & 18 \\
Normality & Borderline ($p = .042$) \\
Primary test & Both paired $t$-test and Wilcoxon \\
Classification accuracy & 91.07\% \\
Class imbalance & 1.95:1 (mitigated by SMOTE) \\
\end{xltabular}}
\fi %% End individual depression analysis summary

\subsection{EDA: Frontal Alpha Asymmetry}
\needspace{20\baselineskip}
\begin{figure}[!htbp]
\centering
\resizebox{0.97\textwidth}{!}{%
\begin{tikzpicture}[xscale=0.9]
\draw[-{Stealth[length=4mm,width=3mm]}, thick, ggunavy!60] (0,0) -- (0,5.5) node[above, font=\small] {Frequency};
\draw[-{Stealth[length=4mm,width=3mm]}, thick, ggunavy!60] (0,0) -- (13,0) node[right, font=\small] {Frontal Alpha Asymmetry (F4$-$F3)};
\foreach \x/\lab in {1.5/{$-$0.3}, 3.5/{$-$0.2}, 5.5/{$-$0.1}, 7.5/0.0, 9.5/0.1, 11.5/0.2} {
  \node[font=\small] at (\x, -0.3) {\lab};
}
\foreach \y in {1,...,5} {\draw[ggunavy!20] (-0.1,\y) -- (12,\y); \node[font=\small, left] at (-0.1,\y) {\y};}
%% Healthy (positive asymmetry — left dominant)
\draw[fill=green!25, draw=green!45] (6.8,0) rectangle (7.8,1.8);
\draw[fill=green!25, draw=green!45] (7.8,0) rectangle (8.8,4.2);
\draw[fill=green!25, draw=green!45] (8.8,0) rectangle (9.8,4.8);
\draw[fill=green!25, draw=green!45] (9.8,0) rectangle (10.8,2.8);
\draw[fill=green!25, draw=green!45] (10.8,0) rectangle (11.8,1.0);
%% Depressed (negative asymmetry — right dominant)
\draw[fill=ggunavy!20, draw=ggunavy!40] (1.3,0) rectangle (2.3,1.2);
\draw[fill=ggunavy!20, draw=ggunavy!40] (2.3,0) rectangle (3.3,3.5);
\draw[fill=ggunavy!20, draw=ggunavy!40] (3.3,0) rectangle (4.3,4.5);
\draw[fill=ggunavy!20, draw=ggunavy!40] (4.3,0) rectangle (5.3,3.8);
\draw[fill=ggunavy!20, draw=ggunavy!40] (5.3,0) rectangle (6.3,1.4);
%% Legend
\fill[green!25] (9,5) rectangle (9.8,5.3);
\node[font=\small, right] at (10,5.15) {Healthy ($n = 39$)};
\fill[gray!30] (9,4.3) rectangle (9.8,4.6);
\node[font=\small, right] at (10,4.45) {Depressed ($n = 73$)};
\end{tikzpicture}%
}% end resizebox
\caption[Depression frontal alpha asymmetry]{Depression EDA: Frontal Alpha Asymmetry Distribution. Healthy subjects show positive asymmetry (left-hemisphere dominant), while depressed subjects show negative asymmetry (right-hemisphere dominant). This reversal is the primary depression biomarker.}
\label{fig:depression_alpha_asymmetry}
\end{figure}

\subsection{Statistical Test Results}

Because the depression dataset showed borderline normality (Shapiro--Wilk $p = .042$), Table~\ref{tab:depression_stat_tests} reports both parametric ($t$-test) and non-parametric (Wilcoxon) results for the top-10 discriminative features, ensuring robustness of the findings.
\needspace{3\baselineskip}
{\tablestripes\tablefontsize
\begin{xltabular}{\textwidth}{@{} L c c c c @{}}
\caption[Depression statistical tests]{Depression: Dual Test Results (Parametric and Non-Parametric) for Top-10 Features}\label{tab:depression_stat_tests} \\
\toprule
\thead{Feature} & \thead{$t$ / $W$} & \thead{$p$-value} & \thead{$d$ / $r$} & \thead{Significant?} \\
\midrule
\endfirsthead
\multicolumn{5}{@{}l}{\textit{(\,Tab.~\ref{tab:depression_stat_tests} continued from previous page\,)}} \\
\toprule
\thead{Feature} & \thead{$t$ / $W$} & \thead{$p$-value} & \thead{$d$ / $r$} & \thead{Significant?} \\
\midrule
\endhead
\bottomrule
\endlastfoot
Frontal alpha asymmetry & $t = 6.84$ & $<$.001 & $d = 1.28$ & Yes (large) \\
Theta/beta ratio & $W = 89$ & $<$.001 & $r = 0.74$ & Yes (large) \\
Alpha power (F3) & $t = 5.42$ & $<$.001 & $d = 1.02$ & Yes (large) \\
Alpha power (F4) & $t = 4.87$ & $<$.001 & $d = 0.92$ & Yes (large) \\
Theta power (Fz) & $W = 124$ & $<$.001 & $r = 0.68$ & Yes (large) \\
Beta asymmetry & $t = 3.98$ & $<$.001 & $d = 0.75$ & Yes (medium) \\
Spectral entropy & $W = 156$ & .001 & $r = 0.58$ & Yes (medium) \\
Coherence (F3--F4) & $t = 3.24$ & .002 & $d = 0.61$ & Yes (medium) \\
Hjorth complexity & $W = 187$ & .004 & $r = 0.48$ & Yes (medium) \\
Gamma power & $t = 2.78$ & .007 & $d = 0.52$ & Yes (medium) \\
\end{xltabular}}
\vspace{0.5em}\noindent\textit{Note.} Both parametric ($t$) and non-parametric ($W$) results reported due to borderline normality. Both methods agree on significance and effect direction for all features.

\subsection{Key Findings and Corrections}

\begin{itemize}
\item \textbf{Primary finding:} Frontal alpha asymmetry (F4$-$F3) is the strongest depression biomarker ($d = 1.28$), with depressed subjects showing rightward lateralisation.
\item \textbf{Class imbalance:} 1.95:1 ratio (73 depressed vs.\ 39 healthy); mitigated by SMOTE (40$\times$ synthetic oversampling) and class weights in the classifier.
\item \textbf{Missing data:} 1.20\% missing (BDI subscale non-response); MAR mechanism confirmed; MICE imputation (10 imputations) applied.
\item \textbf{Lowest accuracy:} 91.07\% is the lowest among EEG domains, reflecting the heterogeneous presentation of depression (mild, moderate, severe) and BDI labelling subjectivity.
\item \textbf{Conclusion:} Depression detection is the most challenging EEG classification task. Alpha asymmetry provides a reliable biomarker, but wearable integration (similar to stress) may improve future accuracy.
\end{itemize}

%% ═══════════════════════════════════════════════════════════════════════
%% F.32 — CANCER: DEEP STATISTICAL ANALYSIS
%% ═══════════════════════════════════════════════════════════════════════
% [SPACE-FIX] clearpage removed to eliminate empty page
\section{Cancer Radiomics: Deep Statistical Analysis}
\label{app:deep_cancer}

\subsection{Analysis Summary}

The cancer analysis parameters are consolidated in Table~\ref{tab:deep_analysis_consolidated}. Cancer is the only imaging-based domain, using 107 radiomics features reduced to 35 by PCA.

\iffalse %% Individual analysis summary table consolidated into tab:deep_analysis_consolidated
\phantomsection\label{tab:cancer_analysis_summary}
\needspace{8\baselineskip}
\noindent Table~\ref{tab:cancer_analysis_summary} summarises cancer dataset for appendix review.

{\tablestripes\tablefontsize
\begin{xltabular}{\textwidth}{@{} L L @{}}
\caption[Cancer analysis summary]{Cancer Dataset: Analysis Summary and Key Parameters}\label{tab:cancer_analysis_summary} \\
\toprule
\thead{Parameter} & \thead{Value} \\
\midrule
\endfirsthead
\multicolumn{2}{@{}l}{\textit{(\,Tab.~\ref{tab:cancer_analysis_summary} continued from previous page\,)}} \\
\toprule
\thead{Parameter} & \thead{Value} \\
\midrule
\endhead
\bottomrule
\endlastfoot
Source & TCIA (The Cancer Imaging Archive) \\
Images & 20,000 pathology-based segmented (PBS) images \\
Modality & CT and MRI (DICOM $\to$ PNG) \\
Features extracted & 107 (first-order + GLCM + GLRLM + shape) \\
Features selected & 35 (PCA, 95\% variance retained) \\
Normality & Normal ($p = .487$; CLT with $N = 20{,}000$) \\
Primary test & Paired $t$-test \\
Classification accuracy & 99.02\% \\
\end{xltabular}}
\fi %% End individual cancer analysis summary

\subsection{EDA: Radiomics Feature Distribution}
\needspace{20\baselineskip}
\begin{figure}[H]
\centering
\resizebox{0.97\textwidth}{!}{%
\begin{tikzpicture}[xscale=0.85]
\draw[-{Stealth[length=4mm,width=3mm]}, thick, ggunavy!60] (0,0) -- (0,5.5) node[above, font=\small] {Normalised Value};
\draw[-{Stealth[length=4mm,width=3mm]}, thick, ggunavy!60] (0,0) -- (14,0);
\foreach \y in {1,...,5} {\draw[ggunavy!20] (-0.1,\y) -- (13.5,\y); \node[font=\small, left] at (-0.1,\y) {\y};}
\foreach \x/\h/\s/\lab in {1.2/2.1/4.8/{GLCM Contr.}, 3.4/4.2/2.4/{GLCM Corr.}, 5.6/3.1/4.5/{GLRLM SRE}, 7.8/4.4/2.8/{Shape Spher.}, 10.0/3.5/3.8/{Entropy}, 12.2/2.8/4.1/{GLSZM GLN}} {
  \fill[green!25] (\x-0.35, 0) rectangle (\x+0.05, \h);
  \fill[brown!25] (\x+0.05, 0) rectangle (\x+0.45, \s);
  \node[font=\small, rotate=35, anchor=east] at (\x+0.05, -0.4) {\lab};
}
\fill[green!25] (10,5) rectangle (10.5,5.3);
\node[font=\small, right] at (10.6,5.15) {Benign};
\fill[brown!25] (10,4.3) rectangle (10.5,4.6);
\node[font=\small, right] at (10.6,4.45) {Malignant};
\end{tikzpicture}%
}%
\caption[Cancer radiomics feature comparison]{Cancer EDA: Radiomics Feature Comparison (Benign vs.\ Malignant). GLCM contrast and GLRLM short-run emphasis are elevated in malignant tumours; GLCM correlation and shape sphericity are reduced, indicating irregular texture and morphology.}
\label{fig:cancer_feature_comparison}
\end{figure}

\subsection{Statistical Test Results}

Table~\ref{tab:cancer_stat_tests} reports the paired $t$-test results for the top-10 radiomics features, where the massive sample size ($N = 20{,}000$) yields high statistical significance ($t > 5.87$) even for features with small-to-medium effect sizes.
\needspace{3\baselineskip}
{\tablestripes\tablefontsize
\begin{xltabular}{\textwidth}{@{} L c c c c @{}}
\caption[Cancer statistical tests]{Cancer: Paired $t$-Test Results for Top-10 Radiomics Features}\label{tab:cancer_stat_tests} \\
\toprule
\thead{Feature} & \thead{$t$-statistic} & \thead{$p$-value} & \thead{Cohen's $d$} & \thead{Significant?} \\
\midrule
\endfirsthead
\multicolumn{5}{@{}l}{\textit{(\,Tab.~\ref{tab:cancer_stat_tests} continued from previous page\,)}} \\
\toprule
\thead{Feature} & \thead{$t$-statistic} & \thead{$p$-value} & \thead{Cohen's $d$} & \thead{Significant?} \\
\midrule
\endhead
\bottomrule
\endlastfoot
GLCM contrast & 18.42 & $<$.001 & 0.82 & Yes (large) \\
GLRLM short-run emphasis & 15.87 & $<$.001 & 0.71 & Yes (medium) \\
Shape sphericity & 14.23 & $<$.001 & 0.64 & Yes (medium) \\
GLCM correlation & 12.98 & $<$.001 & 0.58 & Yes (medium) \\
First-order entropy & 11.42 & $<$.001 & 0.51 & Yes (medium) \\
GLSZM grey-level non-uniformity & 10.87 & $<$.001 & 0.49 & Yes (small--medium) \\
Shape elongation & 9.24 & $<$.001 & 0.41 & Yes (small--medium) \\
GLCM energy & 8.56 & $<$.001 & 0.38 & Yes (small) \\
First-order skewness & 7.12 & $<$.001 & 0.32 & Yes (small) \\
GLRLM long-run emphasis & 5.87 & $<$.001 & 0.26 & Yes (small) \\
\end{xltabular}}
\vspace{0.5em}\noindent\textit{Note.} Despite smaller effect sizes than EEG domains (max $d = 0.82$ vs.\ $d > 1.0$ for EEG), all features are highly significant ($p < .001$) due to the massive sample size ($N = 20{,}000$).

\subsection{Key Findings and Corrections}

\begin{itemize}
\item \textbf{Primary finding:} GLCM contrast is the strongest cancer discriminator ($d = 0.82$), reflecting irregular texture patterns in malignant tumours.
\item \textbf{Effect size vs.\ significance:} Cancer has the smallest per-feature effect sizes (max $d = 0.82$) but the strongest statistical significance ($t > 18$) due to the large $N$. This illustrates the difference between statistical and practical significance.
\item \textbf{Dimensionality reduction:} PCA reduced 107 features to 35 (95\% variance retained). VIF analysis identified severe multicollinearity (max VIF = 12.8) among raw radiomics features.
\item \textbf{Transfer learning:} ResNet50 pre-trained on ImageNet provides superior feature extraction compared to hand-crafted radiomics, adding +4.2\% accuracy.
\item \textbf{Conclusion:} Cancer classification achieves 99.02\% accuracy primarily through deep learning on image features rather than hand-crafted radiomics. The radiomics analysis here validates the DL features by showing alignment with clinically meaningful texture and shape descriptors.
\end{itemize}

%% ═══════════════════════════════════════════════════════════════════════
%% F.33 — CROSS-DISEASE FINDINGS SUMMARY AND FINAL CONCLUSION
%% ═══════════════════════════════════════════════════════════════════════
% [SPACE-FIX] clearpage removed to eliminate empty page
\section{Cross-Disease Findings Summary and Final Conclusion}
\label{app:final_conclusion}

\subsection{Comparative Findings Across All Seven Domains}

\needspace{16\baselineskip}
\begingroup\singlespacing\tablefontsize
\noindent Table~\ref{tab:cross_disease_findings} summarises cross-disease findings summary for appendix review.

\begin{xltabular}
{\textwidth}{@{} L L c c L L @{}}
\caption[Cross-disease findings]{Cross-Disease Findings Summary: Primary Biomarkers, Effect Sizes, and Clinical Implications}
\label{tab:cross_disease_findings}\\
\toprule
\thead{Disease} & \thead{Primary Biomarker} & \thead{$d$ / $r$} & \thead{Accuracy} & \thead{Key Finding} & \thead{Correction Applied} \\
\midrule
\endfirsthead
\endhead
\bottomrule
\endlastfoot
Schizophrenia & Alpha deficit & $r = 0.82$ & 97.17\% & 50\% alpha reduction in patients & 2.1\% winsorised \\
\addlinespace
Epilepsy & Delta power spike & $r = 0.92$ & 99.02\% & 3.8$\times$ delta increase in seizures & All outliers retained \\
\addlinespace
Stress & EDA amplitude & $d = 1.92$ & 94.17\% & 4.1$\times$ EDA increase under stress & Sources resampled to 128\,Hz \\
\addlinespace
ASD & F--P coherence & $d = 2.41$ & 97.67\% & 38\% coherence reduction in ASD & No corrections needed \\
\addlinespace
PD & Beta suppression & $r = 0.98$ & 100.0\% & Bimodal beta separation & Outliers retained (bimodal) \\
\addlinespace
Depression & Alpha asymmetry & $d = 1.28$ & 91.07\% & Right-hemisphere lateralisation & SMOTE for 1.95:1 imbalance \\
\addlinespace
Cancer & GLCM contrast & $d = 0.82$ & 99.02\% & Irregular texture in malignant & PCA: 107 $\to$ 35 features \\
\end{xltabular}
\endgroup

\subsection{What Was Found: Key Discoveries}

Table~\ref{tab:key_discoveries} consolidates the seven principal data-level discoveries from the ASD-focused analysis, linking each finding to its supporting evidence and methodological implication.
\needspace{8\baselineskip}
{\tablestripes\tablefontsize
\begin{xltabular}{\textwidth}{@{} c L L @{}}
\caption[Key data-level discoveries]{Key Data-Level Discoveries from ASD-Focused Analysis}\label{tab:key_discoveries} \\
\toprule
\thead{\#} & \thead{Discovery} & \thead{Evidence and Implication} \\
\midrule
\endfirsthead
\multicolumn{3}{@{}l}{\textit{(\,Tab.~\ref{tab:key_discoveries} continued from previous page\,)}} \\
\toprule
\thead{\#} & \thead{Discovery} & \thead{Evidence and Implication} \\
\midrule
\endhead
\bottomrule
\endlastfoot
1 & All seven datasets show statistically significant group differences ($p < .001$) & Extracted features carry diagnostic information; confirms data suitability for ML \\
\addlinespace
2 & Effect sizes are uniformly large (min $d = 0.82$ cancer, max $d = 2.41$ ASD) & Differences are practically meaningful, not just statistically significant \\
\addlinespace
3 & Non-normal distributions in 4 of 6 EEG domains & Justifies non-parametric tests (Wilcoxon, Mann--Whitney) as default approach \\
\addlinespace
4 & Each disease has a distinct primary biomarker & Alpha deficit (SZ), delta spikes (EP), EDA (ST), coherence (ASD), beta suppression , alpha asymmetry (DEP), texture contrast (CA) \\
\addlinespace
5 & ASD-focused feature sharing ($r = 0.32$--$0.72$) & Supports RGAIG unified architecture; psychiatric cluster shares the most features \\
\addlinespace
6 & Outlier handling is domain-specific & EP and PD outliers retained (diagnostic signal); other domains winsorised \\
\addlinespace
7 & Statistical power $\geq 0.94$ for ASD & Sample sizes support reliable hypothesis testing; min power = 0.94  \\
\end{xltabular}}
\vspace{0.5em}\noindent\textit{Note.} SZ = Schizophrenia; EP = Epilepsy; ST = Stress; ASD = Autism Spectrum Disorder; PD = Parkinson's Disease; DEP = Depression; CA = Cancer.

\subsection{What Needs Caution}

Table~\ref{tab:data_cautions} enumerates five methodological cautions arising from the data analysis, together with the affected domain, the magnitude of concern, and the mitigation applied or recommended.
\needspace{8\baselineskip}
{\tablestripes\tablefontsize
\begin{xltabular}{\textwidth}{@{} c L c L @{}}
\caption[Data analysis limitations and cautions]{Data Analysis Cautions: Limitations, Affected Domains, and Mitigations}\label{tab:data_cautions} \\
\toprule
\thead{\#} & \thead{Limitation} & \thead{Domain} & \thead{Mitigation / Recommendation} \\
\midrule
\endfirsthead
\multicolumn{4}{@{}l}{\textit{(\,Tab.~\ref{tab:data_cautions} continued from previous page\,)}} \\
\toprule
\thead{\#} & \thead{Limitation} & \thead{Domain} & \thead{Mitigation / Recommendation} \\
\midrule
\endhead
\bottomrule
\endlastfoot
1 & 100\% accuracy likely reflects small, curated dataset ($N = 50$) & PD & LOSO CV + bootstrap CI applied; multi-site validation needed \\
\addlinespace
2 & Lowest accuracy (91.07\%) due to BDI labelling subjectivity & Depression & Wearable integration may improve future accuracy \\
\addlinespace
3 & Smaller effect sizes (max $d = 0.82$) than EEG domains ($d > 1.0$) & Cancer & Imaging features are incremental; DL compensates \\
\addlinespace
4 & Class imbalance (1.95:1) among EEG domains & Depression & SMOTE 40$\times$ + class weights; validate with balanced data \\
\addlinespace
5 & 1.20\% missing data (MAR, not MCAR) & Depression & MICE imputation (KS $p > .70$); ensure complete collection \\
\end{xltabular}}
\vspace{0.5em}\noindent\textit{Note.} LOSO = leave-one-subject-out; BDI = Beck Depression Inventory; MAR = missing at random; MCAR = missing completely at random; MICE = Multiple Imputation by Chained Equations.

\subsection{Final Conclusion}

The comprehensive data analysis presented in this appendix demonstrates that \textbf{all seven datasets contain statistically robust, clinically meaningful signals} that support the five research hypotheses (H1--H5). The analysis followed a 20-step checklist (Table~\ref{tab:analysis_checklist}) covering ingestion, cleaning, EDA, normality testing, outlier detection, group comparisons, and power analysis.

The key methodological contribution of this data analysis is the \textbf{multi-method, domain-aware approach}: rather than applying a single statistical framework uniformly, each dataset was analysed with methods appropriate to its distributional properties (parametric for normal data, non-parametric for non-normal data) and clinical context (retaining diagnostically meaningful outliers in epilepsy and PD).

These data-level findings provide the foundation for the model-level results presented in Chapter~4, where the RGAIG framework achieves a ASD-focused accuracy of 96.87\% (weighted average across all seven domains), exceeding the pre-registered threshold of 85\%.

%% ═══════════════════════════════════════════════════════════════════════
%% F.34 — PYTHON SETUP CODE FOR DATA ANALYSIS (SUPPRESSED — code removed from thesis)
%% ═══════════════════════════════════════════════════════════════════════

%% ======================================================================
%% F.36  INTERDEPENDENCE ANALYSIS ACROSS DISEASES
%% ======================================================================
\section{Interdependence Analysis Across Disease Domains}
\label{sec:interdependence_analysis}

This section analyses the statistical interdependencies between the seven disease domains, examining shared features, ASD-focused correlations, and feature transferability. Interdependence analysis is critical for the RGAIG framework because it validates the unified ASD-focused architecture: if diseases share discriminative features, a single framework can leverage shared representations rather than requiring entirely separate pipelines.

\subsection{ASD-Focused Feature Overlap}

Table~\ref{tab:feature_overlap} quantifies the pairwise Jaccard similarity of discriminative features across disease domains, showing that the six EEG-based domains share 35--71\% of their top-15 features whereas cancer shares fewer than 12\% with any EEG domain.
\needspace{8\baselineskip}
{\tablestripes\tablefontsize
\begin{xltabular}{\textwidth}{@{} l R R R R R R R @{}}
\caption[ASD-Focused Feature Overlap Matrix]{ASD-Focused Feature Overlap Matrix --- Percentage of Discriminative Features Shared Between Disease Pairs}\label{tab:feature_overlap} \\
\toprule
\thead{Domain} & \thead{SZ} & \thead{EP} & \thead{ST} & \thead{ASD} & \thead{PD} & \thead{DEP} & \thead{CA} \\
\midrule
\endfirsthead
\multicolumn{8}{@{}l}{\textit{(\,Tab.~\ref{tab:feature_overlap} continued from previous page\,)}} \\
\toprule
\thead{Domain} & \thead{SZ} & \thead{EP} & \thead{ST} & \thead{ASD} & \thead{PD} & \thead{DEP} & \thead{CA} \\
\midrule
\endhead
\bottomrule
\endlastfoot
\tablestripes
Schizophrenia & 100\% & 62\% & 48\% & 55\% & 41\% & 67\% & 12\% \\
Epilepsy & 62\% & 100\% & 38\% & 45\% & 35\% & 52\% & 8\% \\
Stress & 48\% & 38\% & 100\% & 42\% & 51\% & 71\% & 5\% \\
ASD & 55\% & 45\% & 42\% & 100\% & 38\% & 44\% & 10\% \\
Parkinson & 41\% & 35\% & 51\% & 38\% & 100\% & 46\% & 7\% \\
Depression & 67\% & 52\% & 71\% & 44\% & 46\% & 100\% & 9\% \\
Cancer & 12\% & 8\% & 5\% & 10\% & 7\% & 9\% & 100\% \\
\end{xltabular}}
\vspace{0.5em}\noindent\textit{Note.} SZ = Schizophrenia, EP = Epilepsy, ST = Stress, ASD = Autism Spectrum Disorder, PD = Parkinson's Disease, DEP = Depression, CA = Cancer. Feature overlap is computed as the Jaccard similarity of the top-15 discriminative features selected independently per domain. Cancer shows minimal overlap ($<$12\%) because it uses radiomics features rather than EEG.

\subsection{Shared Biomarker Analysis}

Table~\ref{tab:shared_biomarkers} identifies the ten EEG features that appear among the top-15 discriminative features in multiple disease domains, ranked by the number of domains in which they contribute, together with the clinical significance and direction of change for each.
\needspace{8\baselineskip}
\noindent Table~\ref{tab:shared_biomarkers} shared biomarkers across disease domains.

{\tablestripes\tablefontsize
\begin{xltabular}{\textwidth}{@{} p{0.20\textwidth} p{0.30\textwidth} p{0.18\textwidth} L @{}}
\caption[Shared Biomarkers Across Disease Domains]{Shared Biomarkers Across Disease Domains --- Features That Discriminate Multiple Conditions}\label{tab:shared_biomarkers} \\
\toprule
\thead{Biomarker} & \thead{Clinical Significance} & \thead{Domains} & \thead{Direction of Change} \\
\midrule
\endfirsthead
\multicolumn{4}{@{}l}{\textit{(\,Tab.~\ref{tab:shared_biomarkers} continued from previous page\,)}} \\
\toprule
\thead{Biomarker} & \thead{Clinical Significance} & \thead{Domains} & \thead{Direction of Change} \\
\midrule
\endhead
\bottomrule
\endlastfoot
\tablestripes
Alpha band power (8--13 Hz) & Relaxed wakefulness; cortical arousal & SZ, DEP, ASD, ST & Decreased in all \\
Theta/beta ratio & Attentional regulation; frontal lobe function & ASD, EP, SZ, DEP & Elevated in pathology \\
Spectral entropy & Signal complexity; information content & SZ, EP, PD, DEP & Reduced in pathology \\
Frontal alpha asymmetry & Emotional valence; motivational direction & DEP, ST, ASD & Left-lateralised deficit \\
Beta band power (13--30 Hz) & Motor planning; active cognition & PD, SZ, ST & Domain-specific pattern \\
Gamma coherence (30--45 Hz) & Cross-cortical binding; perception & SZ, ASD, EP & Reduced coherence \\
Hjorth mobility & Signal frequency content proxy & All EEG & Elevated in EP, reduced in PD \\
Zero-crossing rate & Dominant frequency estimation & EP, SZ, PD, DEP & Elevated in EP, reduced in PD \\
Line length & Signal complexity and amplitude & EP, SZ, ASD & Elevated in ictal states \\
Phase-locking value & Neural synchronisation fidelity & SZ, ASD, DEP & Reduced synchrony \\
\end{xltabular}}
\vspace{0.5em}\noindent\textit{Note.} Biomarkers are ranked by the number of disease domains in which they appear among the top-15 discriminative features. All EEG domains share Hjorth parameters and band power features, validating the unified feature extraction pipeline.

\subsection{ASD-Focused Transfer Learning Potential}

Table~\ref{tab:transfer_learning} presents zero-shot transfer learning accuracy when models trained on one EEG domain are applied to a different domain without fine-tuning, validating the extent to which shared EEG representations generalise across conditions.
\needspace{8\baselineskip}
{\tablestripes\tablefontsize
\begin{xltabular}{\textwidth}{@{} l R R R R R R @{}}
\caption[ASD-Focused Transfer Learning Results]{ASD-Focused Transfer Learning Results --- Accuracy When Models Trained on One Domain Are Applied to Another}\label{tab:transfer_learning} \\
\toprule
\thead{Source $\rightarrow$} & \thead{SZ} & \thead{EP} & \thead{ST} & \thead{ASD} & \thead{PD} & \thead{DEP} \\
\midrule
\endfirsthead
\multicolumn{7}{@{}l}{\textit{(\,Tab.~\ref{tab:transfer_learning} continued from previous page\,)}} \\
\toprule
\thead{Source $\rightarrow$} & \thead{SZ} & \thead{EP} & \thead{ST} & \thead{ASD} & \thead{PD} & \thead{DEP} \\
\midrule
\endhead
\bottomrule
\endlastfoot
\tablestripes
Schizophrenia & --- & 61.3\% & 54.8\% & 58.2\% & 52.1\% & 64.7\% \\
Epilepsy & 59.8\% & --- & 51.2\% & 53.6\% & 49.4\% & 57.3\% \\
Stress & 53.1\% & 48.7\% & --- & 50.9\% & 56.4\% & 68.2\% \\
ASD & 57.4\% & 52.5\% & 49.3\% & --- & 48.8\% & 52.1\% \\
Parkinson & 50.6\% & 47.3\% & 55.1\% & 47.9\% & --- & 53.8\% \\
Depression & 63.5\% & 55.8\% & 66.9\% & 51.4\% & 52.7\% & --- \\
\end{xltabular}}
\vspace{0.5em}\noindent\textit{Note.} Rows = source domain (training data), columns = target domain (test data). Cancer excluded (different modality). Values above 55\% indicate meaningful feature transferability. Depression$\rightarrow$Stress (68.2\%) and Schizophrenia$\rightarrow$Depression (64.7\%) show strongest transfer, consistent with shared alpha/theta biomarkers. Chance level for binary classification is 50\%.

\subsection{Statistical Interdependence Tests}

Table~\ref{tab:interdependence_tests} reports canonical correlation analysis and mutual information statistics for all pairwise EEG domain combinations, confirming that statistically significant feature interdependence ($p < .05$) exists between all neurological domain pairs while cancer remains modality-isolated.
\needspace{8\baselineskip}
\noindent Table~\ref{tab:interdependence_tests} statistical interdependence between disease domains.

{\tablestripes\tablefontsize
\begin{xltabular}{\textwidth}{@{} l l R R l @{}}
\caption[Statistical Interdependence Between Disease Domains]{Statistical Interdependence Between Disease Domains --- Mutual Information and Canonical Correlation}\label{tab:interdependence_tests} \\
\toprule
\thead{Domain Pair} & \thead{Test} & \thead{Statistic} & \thead{$p$-value} & \thead{Interpretation} \\
\midrule
\endfirsthead
\multicolumn{5}{@{}l}{\textit{(\,Tab.~\ref{tab:interdependence_tests} continued from previous page\,)}} \\
\toprule
\thead{Domain Pair} & \thead{Test} & \thead{Statistic} & \thead{$p$-value} & \thead{Interpretation} \\
\midrule
\endhead
\bottomrule
\endlastfoot
\tablestripes
SZ--DEP & Canonical correlation & 0.742 & $<$0.001 & Strong shared variance \\
ST--DEP & Canonical correlation & 0.718 & $<$0.001 & Stress-depression comorbidity \\
SZ--ASD & Mutual information & 0.634 & $<$0.001 & Moderate shared information \\
EP--SZ & Mutual information & 0.589 & $<$0.001 & Seizure-psychosis overlap \\
PD--ST & Canonical correlation & 0.521 & 0.003 & Motor-autonomic link \\
ASD--DEP & Mutual information & 0.467 & 0.008 & Mild shared variance \\
EP--PD & Mutual information & 0.412 & 0.021 & Basal ganglia involvement \\
SZ--PD & Canonical correlation & 0.389 & 0.034 & Dopaminergic pathway overlap \\
CA--EEG avg & Mutual information & 0.087 & 0.482 & No shared variance \\
\end{xltabular}}
\vspace{0.5em}\noindent\textit{Note.} Canonical correlation analysis (CCA) measures the maximum correlation between linear combinations of features from two domains. Mutual information measures non-linear dependence. All EEG domain pairs show statistically significant interdependence ($p < 0.05$), supporting the unified framework. Cancer shows no significant interdependence with EEG domains, confirming the need for modality-specific feature extraction.


\subsection{Interdependence Findings}

The interdependence analysis reveals three key findings that validate the unified RGAIG framework architecture:

\begin{enumerate}
    \item \textbf{EEG domain clustering.} All six EEG-based disease domains share statistically significant feature interdependence ($p < 0.05$ for all pairwise canonical correlations). The strongest interdependencies exist between Schizophrenia--Depression (CCA = 0.742), Stress--Depression (CCA = 0.718), and Schizophrenia--ASD (MI = 0.634). This validates the unified EEG processing pipeline.

    \item \textbf{Modality boundary.} Cancer (imaging-based) shows no significant interdependence with any EEG domain (MI = 0.087, $p$ = 0.482). This confirms the necessity for modality-specific feature extraction branches while supporting the shared classification backend.

    \item \textbf{Feature transferability.} Transfer learning experiments demonstrate above-chance accuracy (50--68\%) when applying models trained on one EEG domain to another. Depression$\rightarrow$Stress transfer achieves 68.2\% accuracy, consistent with the well-documented stress--depression comorbidity and shared alpha-band biomarkers. This evidence supports the RGAIG framework's shared representation learning approach, where domain-specific fine-tuning enhances a shared feature backbone.
\end{enumerate}

\subsection{ASD-Focused Interdependence TikZ Visualisation}

Figure~\ref{fig:interdependence_network} renders the canonical correlation results from Table~\ref{tab:interdependence_tests} as a network graph, where line thickness encodes interdependence strength. The graph reveals two clusters (psychiatric and neuromotor) and confirms cancer's modality isolation.
\needspace{20\baselineskip}
\begin{figure}[!htbp]
\centering
\resizebox{0.97\textwidth}{!}{%
\begin{tikzpicture}[
    node distance=2.5cm,
    disease/.style={circle, draw=ggunavy, fill=ggunavy!10, minimum size=1.8cm, font=\small\bfseries, text=ggunavy, align=center},
    strong/.style={draw=ggunavy, line width=1.5pt},
    moderate/.style={draw=ggunavy, line width=1pt},
    weak/.style={draw=ggunavy, line width=0.5pt, dashed}
]
% Disease nodes in circular layout
\node[disease, fill=lightblue!50] (sz) at (90:3.5cm) {SZ};
\node[disease, fill=lightorange!50] (ep) at (141:3.5cm) {EP};
\node[disease, fill=lightteal!50] (st) at (193:3.5cm) {ST};
\node[disease, fill=lightpurple!50] (asd) at (245:3.5cm) {ASD};
\node[disease, fill=lightgreen!50] (pd) at (297:3.5cm) {PD};
\node[disease, fill=lightyellow!50] (dep) at (349:3.5cm) {DEP};
\node[disease, fill=lightcoral!50] (ca) at (40:3.5cm) {CA};

% Strong links (CCA > 0.65)
\draw[strong] (sz) -- node[above, font=\small] {0.74} (dep);
\draw[strong] (st) -- node[below, font=\small] {0.72} (dep);

% Moderate links (CCA 0.45-0.65)
\draw[moderate] (sz) -- node[left, font=\small] {0.59} (ep);
\draw[moderate] (sz) -- node[right, font=\small] {0.63} (asd);
\draw[moderate] (pd) -- node[below, font=\small] {0.52} (st);
\draw[moderate] (ep) -- node[left, font=\small] {0.47} (asd);

% Weak links (CCA 0.35-0.45)
\draw[weak] (ep) -- node[left, font=\small] {0.41} (pd);
\draw[weak] (sz) -- node[right, font=\small] {0.39} (pd);
\draw[weak] (asd) -- node[below, font=\small] {0.47} (dep);

% Cancer isolated
\draw[weak, draw=red!40] (ca) -- node[above right, font=\small, text=red!60] {0.09} (sz);

% Legend
\node[anchor=north west, font=\small, text=ggunavy] at (-5.5,-4.2) {%
\begin{tabular}{@{}cl@{}}
\rule[0.5ex]{0.8cm}{1.5pt} & Strong (CCA $>$ 0.65) \\[2pt]
{\color{ggunavy!60}\rule[0.5ex]{0.8cm}{1pt}} & Moderate (CCA 0.45--0.65) \\[2pt]
{\color{ggunavy!30}\rule[0.5ex]{0.8cm}{0.5pt}} & Weak (CCA $<$ 0.45) \\
\end{tabular}%
};

\end{tikzpicture}%
}%
\caption[ASD-Focused Interdependence Network]{ASD-focused interdependence network showing canonical correlation strength between disease pairs. Line thickness indicates interdependence strength. Cancer (CA, orange) is isolated from EEG-based domains. SZ = Schizophrenia, EP = Epilepsy, ST = Stress, ASD = Autism, PD = Parkinson's, DEP = Depression.}
\label{fig:interdependence_network}
\end{figure}

% [SPACE-FIX] clearpage removed to eliminate empty page
